\documentclass[final,12pt,a4paper,twoside,fleqn]{extbook}
\usepackage[T1,LGR]{fontenc}
\usepackage[utf8]{inputenc}
\usepackage[spanish]{babel}
\decimalpoint
\def\@makechapterhead#1{%
  \vspace*{50\p@}%
  {\parindent \z@ \raggedright \normalfont
    \ifnum \c@secnumdepth >\m@ne
      \if@mainmatter
        %\huge\bfseries \@chapapp\space \thechapter
        \huge\bfseries \thechapter.\space%
        %\par\nobreak
        %\vskip 20\p@
      \fi
    \fi
    \interlinepenalty\@M
    \huge \bfseries #1\par\nobreak
    \vskip 40\p@
  }
}
\makeatother
\makeatletter
\addto\captionsspanish{
    \renewcommand{\partname}{}
    \renewcommand{\chaptername}{}
}
\makeatother
\makeatletter
\usepackage[backend=biber]{biblatex}
\addbibresource{bibliografia.bib}
\usepackage{hyperref}
\usepackage{listings}
%\usepackage[noenc]{tipa}
%\usepackage{ly1}
\usepackage[fleqn]{mathtools}
\usepackage{lplfitch}
\usepackage{12many}
\usepackage{abraces}
%\usepackage{abstract}
%\usepackage{yfonts}
\usepackage{color}
\usepackage[usenames,dvipsnames,svgnames,table]{xcolor}
\usepackage{xxcolor}
\usepackage{verbatim}
\usepackage{anyfontsize}
\usepackage{anysize}
\usepackage{helvet}
\usepackage{textcomp}
\usepackage{dsfont}
\usepackage{fancybox}
\usepackage{pgf}
\usepackage{tikz}
\usetikzlibrary{arrows,decorations.shapes,automata,backgrounds,petri,cd}
\usetikzlibrary{positioning}
\usetikzlibrary{shadows}
\usetikzlibrary{chains}
\usepackage{bbm}
\usepackage[inline]{enumitem}
\usepackage{mathdesign}
\usepackage{calrsfs}
\usepackage{mathrsfs}
\usepackage{mathalfa}
\usepackage{amssymb}
%\usepackage{amsrefs}
\usepackage{amsmath}
\usepackage{amsthm}
%\usepackage{amsaddr}
\usepackage{amsfonts}
\usepackage{amsxtra}
\usepackage{amscd}
\usepackage{bm}
%\usepackage{wasysym}
\usepackage{bold-extra}
\usepackage{bbm}
\usepackage[mathscr]{euscript}
\usepackage{color}
\usepackage{units}
\usepackage{nomencl}
\providecommand{\printnomenclature}{\printglossary}
\providecommand{\makenomenclature}{\makeglossary}
\makenomenclature
\usepackage{microtype}
\usepackage{graphicx}
\makeatletter
%% EPISTEMOLOGIA DE LAS MATEM\'ATICAS
\newcommand{\sforall}{\ensuremath{\forall}}
\newcommand{\sexists}{\ensuremath{\exists}}
\newcommand{\snexists}{\ensuremath{\nexists}}
\newcommand{\sand}{\ensuremath{\wedge}}
\newcommand{\sor}{\ensuremath{\vee}}
\newcommand{\sthen}{\ensuremath{\Longrightarrow}}
\newcommand{\sif}{\ensuremath{\Leftarrow}}
\newcommand{\ssi}{\ensuremath{\Leftrightarrow}}
\newcommand{\setomega}{\ensuremath{\omega}}
\newcommand{\setmu}{\ensuremath{\mu}}
\newcommand{\gensig}{\ensuremath{\sigma}}
\newcommand{\bneg}{\ensuremath{\sim}}
\newcommand{\ssimeq}{\ensuremath{\simeq}}
\newcommand{\sunion}{\ensuremath{\cup}}
\newcommand{\sintersection}{\ensuremath{\cap}}
\newcommand{\sproduct}{\ensuremath{\times}}
\newcommand{\sminus}{\ensuremath{\setminus}}
\newcommand{\ssymdiff}{\ensuremath{\triangle}}
\newcommand{\sbin}{\ensuremath{\in}}
\newcommand{\snin}{\ensuremath{\notin}}
\newcommand{\seq}{\ensuremath{=}}
\newcommand{\sneq}{\ensuremath{\neq}}
\newcommand{\eps}{\ensuremath{\varepsilon}}
\newcommand{\beps}{\ensuremath{\backepsilon}}
\newcommand{\toward}{\ensuremath{\longrightarrow}}
\newcommand{\fromto}{\ensuremath{\longrightarrow}}
\newcommand{\elmapstoel}{\ensuremath{\mapsto}}
\newcommand{\absoluto}{\ensuremath{\boldsymbol{\mathbf{V}}}}
\newcommand{\nada}{\ensuremath{\boldsymbol{\mathbf{\Lambda}}}}
\newcommand{\nta}{\ensuremath{\varphi}}
\newcommand{\ntb}{\ensuremath{\psi}}
\newcommand{\ntc}{\ensuremath{\vartheta}}
\newcommand{\vnta}{\ensuremath{\Phi}}
\newcommand{\vntb}{\ensuremath{\Psi}}
\newcommand{\vntc}{\ensuremath{\Theta}}
\newcommand{\chtitletext}[1]{\textbf{\emph{#1}}}
\newcommand{\memph}[1]{\textbf{\emph{#1}}}
\newcommand{\mmemph}[1]{\large{\memph{#1}}}
\newcommand{\bfs}[1]{\ensuremath{\boldsymbol{#1}}}
\newcommand{\bfsf}[1]{\ensuremath{\bfs{\mathsf{#1}}}}
\newcommand{\funcsum}{\ensuremath{\bfs{\mathsf{s}}}}
\newcommand{\funcsumat}[2]{\ensuremath{\funcsum\left({#1,#2}\right)}}
\newcommand{\funcprod}{\ensuremath{\bfs{\mathsf{p}}}}
\newcommand{\funcprodat}[2]{\ensuremath{\funcprod\left({#1,#2}\right)}}
\newcommand{\funccompl}{\ensuremath{\bfs{\mathsf{c}}}}
\newcommand{\funccomplat}[1]{\ensuremath{\funccompl\left({#1}\right)}}
\newcommand{\can}[1]{\ensuremath{\mathfrak{#1}}}
\newcommand{\ntn}[1]{\ensuremath{\bfs{#1}}}
\newcommand{\exn}[1]{\ensuremath{\bfs{\mathscr{#1}}}}
\newcommand{\notas}[1]{\ensuremath{{
{}^{\displaystyle{\bfs{nt}}}{\!}{\bfs{#1}}
}}}
\newcommand{\casos}[1]{\ensuremath{{
    {}^{\displaystyle{\bfs{cs}}}{\!}{\bfs{#1}}
}}}
\newcommand{\notasn}[1]{\ensuremath{
    {\notas{\mathscr{#1}}}
}}
\newcommand{\casosn}[1]{\ensuremath{
    {\casos{\mathscr{#1}}}
}}
\newcommand{\Gen}[1]{\ensuremath{\bfs{
    {\mathcal{G}\mathbf{e}\mathbf{n}}\set{\exn{#1}}
}}}
\newcommand{\GenTo}[2]{\ensuremath{
    \set{#1}
    \bfs{\mathcal{G}\mathbf{e}\mathbf{n}{\_}\mathbf{a}}
    \set{#2}
}}
\newcommand{\ExtGenTo}[2]{\ensuremath{
    \GenTo{\exn{#1}}{\exn{#2}}
}}
\newcommand{\Par}[1]{\ensuremath{
        \bfs{
            {\mathcal{P}\mathbf{a}\mathbf{r}}\set{\mathscr{#1}}
        }
}}
\newcommand{\ParTo}[2]{\ensuremath{
    \set{#1}
    \bfs{\mathcal{P}\mathbf{a}\mathbf{r}\bfs{\_}\mathbf{a}}
    \set{#2}
}}
\newcommand{\ExtParTo}[2]{\ensuremath{
    \ParTo{\exn{#1}}{\exn{#2}}
}}
\newcommand{\tty}{\ensuremath{\mathtt{y}}}
\newcommand{\ttY}{\ensuremath{\mathtt{Y}}}
\newcommand{\ttx}{\ensuremath{\mathtt{x}}}
\newcommand{\entreparentesis}[1]{\ensuremath{
        \left({#1}\right)
}}
\newcommand{\entrecorchetes}[1]{\ensuremath{
        \left\lbrack{#1}\right\rbrack
}}
\newcommand{\entredoblescorchetes}[1]{\ensuremath{
        \left[\!{\left[{#1}\right]}\!\right]
}}
\newcommand{\entreangulos}[1]{\ensuremath{\left\langle{#1}\right\rangle}}
\newcommand{\doblelanglegrande}{\ensuremath{{\big{
                \left\langle\!\!\left\langle
}}}}
\newcommand{\dobleranglegrande}{\ensuremath{{\big{
                \right\rangle\!\!\right\rangle
}}}}
\newcommand{\entregrandesdoblesangulos}[1]{\ensuremath{
    \left\doblelanglegrande
    {#1}
    \right\dobleranglegrande
}}
\newcommand{\entrellaves}[1]{\ensuremath{\left\lbrace{#1}\right\rbrace}}
\newcommand{\entredoblesllaves}[1]{\ensuremath{
        \left\lbrace\!\left\lbrace{#1}\right\rbrace\!\right\rbrace
}}
\newcommand{\desederivaque}[2]{\ensuremath{{#1}\bfs{\;\vdash\;}{#2}}}
\newcommand{\denosederivaque}[2]{\ensuremath{{#1}\bfs{\;\nvdash\;}{#2}}}
\newcommand{\parts}{\ensuremath{\bfs{\mathcal{P}}}}
\newcommand{\partsof}[1]{\ensuremath{\parts{\entreparentesis{\exn{#1}}}}}
\newcommand{\universe}{\ensuremath{\bfs{\mathfrak{U}}}}
\newcommand{\semptyset}{\ensuremath{\emptyset}}
\newcommand{\cna}{\ensuremath{\can{a}}}
\newcommand{\cnb}{\ensuremath{\can{b}}}
\newcommand{\cnc}{\ensuremath{\can{c}}}
\newcommand{\cnd}{\ensuremath{\can{d}}}
\newcommand{\cne}{\ensuremath{\can{e}}}
\newcommand{\cnf}{\ensuremath{\can{f}}}
\newcommand{\cng}{\ensuremath{\can{g}}}
\newcommand{\cnh}{\ensuremath{\can{h}}}
\newcommand{\cni}{\ensuremath{\can{i}}}
\newcommand{\cnj}{\ensuremath{\can{j}}}
\newcommand{\cnk}{\ensuremath{\can{k}}}
\newcommand{\cnl}{\ensuremath{\can{l}}}
\newcommand{\cnm}{\ensuremath{\can{m}}}
\newcommand{\cnn}{\ensuremath{\can{n}}}
\newcommand{\cno}{\ensuremath{\can{o}}}
\newcommand{\cnp}{\ensuremath{\can{p}}}
\newcommand{\cnq}{\ensuremath{\can{q}}}
\newcommand{\cnr}{\ensuremath{\can{r}}}
\newcommand{\cns}{\ensuremath{\can{s}}}
\newcommand{\cnt}{\ensuremath{\can{t}}}
\newcommand{\cnu}{\ensuremath{\can{u}}}
\newcommand{\cnw}{\ensuremath{\can{w}}}
\newcommand{\cnx}{\ensuremath{\can{x}}}
\newcommand{\cny}{\ensuremath{\can{y}}}
\newcommand{\cnz}{\ensuremath{\can{z}}}
\newcommand{\enA}{\ensuremath{\exn{A}}}
\newcommand{\enB}{\ensuremath{\exn{B}}}
\newcommand{\enC}{\ensuremath{\exn{C}}}
\newcommand{\enD}{\ensuremath{\exn{D}}}
\newcommand{\enE}{\ensuremath{\exn{E}}}
\newcommand{\enF}{\ensuremath{\exn{F}}}
\newcommand{\enG}{\ensuremath{\exn{G}}}
\newcommand{\enH}{\ensuremath{\exn{H}}}
\newcommand{\enI}{\ensuremath{\exn{I}}}
\newcommand{\enJ}{\ensuremath{\exn{A}}}
\newcommand{\enK}{\ensuremath{\exn{B}}}
\newcommand{\enL}{\ensuremath{\exn{C}}}
\newcommand{\enM}{\ensuremath{\exn{A}}}
\newcommand{\enN}{\ensuremath{\exn{B}}}
\newcommand{\enO}{\ensuremath{\exn{C}}}
\newcommand{\enP}{\ensuremath{\exn{A}}}
\newcommand{\enQ}{\ensuremath{\exn{B}}}
\newcommand{\enR}{\ensuremath{\exn{C}}}
\newcommand{\enS}{\ensuremath{\exn{A}}}
\newcommand{\enT}{\ensuremath{\exn{B}}}
\newcommand{\enU}{\ensuremath{\exn{C}}}
\newcommand{\enV}{\ensuremath{\exn{A}}}
\newcommand{\enW}{\ensuremath{\exn{B}}}
\newcommand{\enX}{\ensuremath{\exn{C}}}
\newcommand{\enY}{\ensuremath{\exn{A}}}
\newcommand{\enZ}{\ensuremath{\exn{B}}}
\newcommand{\predSet}{\ensuremath{\can{S}\!\can{e}\!\can{t}}}
\newcommand{\predClass}{\ensuremath{
        \can{C}\!\can{l}\!\can{a}\!\can{s}\!\can{s}
}}
\newcommand{\mathdef}{\ensuremath{\stackrel{\mathsf{def}}{{}={}}}}
\newcommand{\pfmathdef}[2]{\ensuremath{{#1}{\mathdef}{#2}}}
\newcommand{\logicdef}{\ensuremath{
        \stackrel
            {\mathsf{def}}
            {\Longleftrightarrow}
}}
\newcommand{\pflogicdef}[2]{\ensuremath{{#1}{\logicdef}{#2}}}
\newcommand{\Real}{\ensuremath{\mathds{R}}}
\newcommand{\logicand}{\ensuremath{\sand }}
\newcommand{\logicor}{\ensuremath{\sor}}
\newcommand{\then}{\ensuremath{\sthen}}
\newcommand{\entonces}{\ensuremath{\sthen}}
\newcommand{\pflogicand}[2]{\ensuremath{{#1}\logicand{#2}}}
\newcommand{\pflogicor}[2]{\ensuremath{{#1}\logicor{#2}}}
\newcommand{\pfthen}[2]{\ensuremath{{#1}\then{#2}}}
\newcommand{\pfentonces}[2]{\ensuremath{{#1}\entonces{#2}}}
\newcommand{\pfsetunion}[2]{\ensuremath{{#1}\sunion{#2}}}
\newcommand{\pfexnunionexn}[2]{\ensuremath{
        \pfsetunion{\exn{#1}}{\exn{#2}}
}}
\newcommand{\pfcsnexnunionexn}[2]{\ensuremath{
        \pfsetunion{\casosn{#1}}{\casosn{#2}}
}}
\newcommand{\pfsetintersection}[2]{\ensuremath{
        {#1}\sintersection{#2}
}}
\newcommand{\pfexnintersectionexn}[2]{\ensuremath{
        \pfsetintersection{\exn{#1}}{\exn{#2}}
}}
\newcommand{\pfcsnexnintersectionexn}[2]{\ensuremath{
        \pfsetintersection{\casosn{#1}}{\casosn{#2}}
}}
\newcommand{\pfsetproduct}[2]{\ensuremath{{#1}{\sproduct}{#2}}}
\newcommand{\pfexnproductexn}[2]{\ensuremath{
        \pfsetproduct{\exn{#1}}{\exn{#2}}
}}
\newcommand{\pfcsnexnproductexn}[2]{\ensuremath{
        \pfsetproduct{\casosn{#1}}{\casosn{#2}}
}}
\newcommand{\pfsetminus}[2]{\ensuremath{{#1}{\sminus}{#2}}}
\newcommand{\pfexnminusexn}[2]{\ensuremath{
        \pfsetminus{\exn{#1}}{\exn{#2}}
}}
\newcommand{\pfcsnexnminusexn}[2]{\ensuremath{
        \pfsetminus{\casosn{#1}}{\casosn{#2}}
}}
\newcommand{\pfsetsymdiff}[2]{\ensuremath{{#1}{\ssymdiff}{#2}}}
\newcommand{\pfexnsymdiffexn}[2]{\ensuremath{
        \pfsetsymdiff{\exn{#1}}{\exn{#2}}
}}
\newcommand{\pfcsnexnsymdiffexn}[2]{\ensuremath{
        \pfsetsymdiff{\casosn{#1}}{\casosn{#2}}
}}
\newcommand{\pfin}[2]{\ensuremath{{#1}\sbin{#2}}}
\newcommand{\pfcsnin}[2]{\ensuremath{{#1}\sbin\casosn{#2}}}
\newcommand{\pfnotin}[2]{\ensuremath{{#1}\snin{#2}}}
\newcommand{\pfcsnnotin}[2]{\ensuremath{{#1}\snin\casosn{#2}}}
\newcommand{\pfeq}[2]{\ensuremath{{#1}\seq{#2}}}
\newcommand{\pfneq}[2]{\ensuremath{{#1}\sneq{#2}}}
\newcommand{\pfexninexn}[2]{\ensuremath{\pfin{\exn{#1}}{\exn{#2}}}}
\newcommand{\pfcaninexn}[2]{\ensuremath{\pfin{\can{#1}}{\exn{#2}}}}
\newcommand{\pfcanincan}[2]{\ensuremath{\pfin{\can{#1}}{\can{#2}}}}
\newcommand{\pfexnincan}[2]{\ensuremath{\pfin{\exn{#1}}{\can{#2}}}}
\newcommand{\pfexneqexn}[2]{\ensuremath{\pfeq{\exn{#1}}{\exn{#2}}}}
\newcommand{\pfcaneqexn}[2]{\ensuremath{\pfeq{\can{#1}}{\exn{#2}}}}
\newcommand{\pfcaneqcan}[2]{\ensuremath{\pfeq{\can{#1}}{\can{#2}}}}
\newcommand{\pfexneqcan}[2]{\ensuremath{\pfeq{\exn{#1}}{\can{#2}}}}
\newcommand{\pfexnnotinexn}[2]{\ensuremath{\pfnotin{\exn{#1}}{\exn{#2}}}}
\newcommand{\pfcannotinexn}[2]{\ensuremath{\pfnotin{\can{#1}}{\exn{#2}}}}
\newcommand{\pfcannotincan}[2]{\ensuremath{\pfnotin{\can{#1}}{\can{#2}}}}
\newcommand{\pfexnnotincan}[2]{\ensuremath{\pfnotin{\exn{#1}}{\can{#2}}}}
\newcommand{\pfexnneqexn}[2]{\ensuremath{\pfneq{\exn{#1}}{\exn{#2}}}}
\newcommand{\pfcanneqexn}[2]{\ensuremath{\pfneq{\can{#1}}{\exn{#2}}}}
\newcommand{\pfcanneqcan}[2]{\ensuremath{\pfneq{\can{#1}}{\can{#2}}}}
\newcommand{\pfexnneqcan}[2]{\ensuremath{\pfneq{\exn{#1}}{\can{#2}}}}
\newcommand{\pfsubseteq}[2]{\ensuremath{{#1} \subseteq {#2}}}
\newcommand{\pfnsubseteq}[2]{\ensuremath{{#1} \nsubseteq {#2}}}
\newcommand{\pfsupseteq}[2]{\ensuremath{{#1} \supseteq {#2}}}
\newcommand{\pfnsupseteq}[2]{\ensuremath{{#1}\nsupseteq{#2}}}
\newcommand{\pfexnsubexneq}[2]{\ensuremath{\pfsubseteq{\exn{#1}}{\exn{#2}}}}
\newcommand{\pfcansubexneq}[2]{\ensuremath{\pfsubseteq{\can{#1}}{\exn{#2}}}}
\newcommand{\pfcansubcaneq}[2]{\ensuremath{\pfsubseteq{\can{#1}}{\can{#2}}}}
\newcommand{\pfexnsubcaneq}[2]{\ensuremath{\pfsubseteq{\exn{#1}}{\can{#2}}}}
\newcommand{\pfexnsupexneq}[2]{\ensuremath{\pfsupseteq{\exn{#1}}{\exn{#2}}}}
\newcommand{\pfcansupexneq}[2]{\ensuremath{\pfsupseteq{\can{#1}}{\exn{#2}}}}
\newcommand{\pfcansupcaneq}[2]{\ensuremath{\pfsupseteq{\can{#1}}{\can{#2}}}}
\newcommand{\pfexnsupcaneq}[2]{\ensuremath{\pfsupseteq{\exn{#1}}{\can{#2}}}}
\newcommand{\pfcansubcan}[2]{\ensuremath{\pfsubset{\can{#1}}{\can{#2}}}}
\newcommand{\pfcansubexn}[2]{\ensuremath{\pfsubset{\can{#1}}{\exn{#2}}}}
\newcommand{\pfexnsubcan}[2]{\ensuremath{\pfsubset{\exn{#1}}{\can{#2}}}}
\newcommand{\pfexnsubexn}[2]{\ensuremath{\pfsubset{\exn{#1}}{\exn{#2}}}}
\newcommand{\pfcansupcan}[2]{\ensuremath{\pfsupset{\can{#1}}{\can{#2}}}}
\newcommand{\pfcansupexn}[2]{\ensuremath{\pfsupset{\can{#1}}{\exn{#2}}}}
\newcommand{\pfexnsupcan}[2]{\ensuremath{\pfsupset{\exn{#1}}{\can{#2}}}}
\newcommand{\pfexnsupexn}[2]{\ensuremath{\pfsupset{\exn{#1}}{\exn{#2}}}}
\newcommand{\existecan}{\ensuremath{\sexists}}
\newcommand{\existecande}[1]{\ensuremath{\existecan{\can{#1}}}}
\newcommand{\existecanen}[2]{\ensuremath{
    \entreparentesis{\existecande{#1}}
    \entreparentesis{\pfin{\can{#1}}{#2}}
}}
\newcommand{\existeexnen}[2]{\ensuremath{
    \entreparentesis{\existecande{#1}}
    \entreparentesis{\pfin{\exn{#1}}{#2}}
}}
\newcommand{\existeexn}{\ensuremath{\sexists}}
\newcommand{\existeexnde}[1]{\ensuremath{\existeexn{\exn{#1}}}}
\newcommand{\existeexnencan}[2]{\ensuremath{
    \left(\existeexnde{#1}\right)
    \left(\pfexnincan{#1}{#2}\right)
}}
\newcommand{\existeexnencanque}[3]{\ensuremath{
\left(\existeexnde{#1}\right)
\entreparentesis{
        \pflogicand{\left(\pfexnincan{#1}{#2}\right)}
        {\entreparentesis{#3}}
    }
}}
\newcommand{\existecanenexn}[2]{\ensuremath{
    \left(\existecande{#1}\right)
    \left(\pfcaninexn{#1}{#2}\right)
}}
\newcommand{\existecanenexnque}[3]{\ensuremath{
\left(\existecande{#1}\right)
\entreparentesis{
        \pflogicand{\left(\pfcaninexn{#1}{#2}\right)}
        {\entreparentesis{#3}}
    }
}}
\newcommand{\existecanencan}[2]{\ensuremath{
    \left(\existecande{#1}\right)
    \left(\pfcanincan{#1}{#2}\right)
}}
\newcommand{\existecanencanque}[3]{\ensuremath{
    \left(\existecande{#1}\right)
    \entreparentesis{
        \pflogicand{\left(\pfcanincan{#1}{#2}\right)}
        {\entreparentesis{#3}}
    }
}}
\newcommand{\existeexnenexn}[2]{\ensuremath{
    \left(\existeexnde{#1}\right)
    \left(\pfexninexn{#1}{#2}\right)
}}
\newcommand{\existeexnenexnque}[3]{\ensuremath{
    \left(\existeexnde{#1}\right)
    \entreparentesis{
          \pflogicand{\left(\pfexninexn{#1}{#2}\right)}
          {\entreparentesis{#3}}
    }
}}
\newcommand{\noexistecanen}[2]{\ensuremath{\noexistecan\pfin{#1}{#2}}}
\newcommand{\noexisteexten}[2]{\ensuremath{\noexistecasoen{\exn{#1}}{\exn{#2}}}}
\newcommand{\noexistecaso}{\ensuremath{\snexists}}
\newcommand{\noexistecasode}[1]{\ensuremath{\noexistecaso{#1}}}
\newcommand{\noexisteextde}[1]{\ensuremath{\noexistecasode{\exn{#1}}}}
\newcommand{\paratodocaso}{\ensuremath{\sforall}}
\newcommand{\paratodocan}{\ensuremath{\sforall}}
\newcommand{\paratodocasode}[1]{\ensuremath{\paratodocaso{#1}}}
\newcommand{\paratodocande}[1]{\ensuremath{\paratodocan{\can{#1}}}}
\newcommand{\paratodocasonombde}[1]{\ensuremath{\paratodocasode{\can{#1}}}}
\newcommand{\paratodocasoextde}[1]{\ensuremath{\paratodocasode{\exn{#1}}}}
\newcommand{\paraninguncasode}[1]{\ensuremath{\paraninguncaso{#1}}}
\newcommand{\paraningunaextde}[1]{\ensuremath{\paraninguncasode{\exn{#1}}}}
\newcommand{\paratodocasoen}[2]{\ensuremath{\paratodocasode{#1}\pfin{#1}{#2}}}
\newcommand{\paratodocanen}[2]{\ensuremath{
        \entreparentesis{
            \paratodocande{#1}
        }
        \entreparentesis{
            \pfin{\can{#1}}{#2}
        }
}}
\newcommand{\esoytalque}[2]{\ensuremath{
        \entreparentesis{
            \pflogicand{#1}{#2}
        }
}}
\newcommand{\esoimplicaque}[2]{\ensuremath{
    \entreparentesis{
        \pfentonces{#1}{#2}
    }
}}
\newcommand{\paratodocanenexnque}[3]{\ensuremath{
    \entreparentesis{\paratodocasode{\can{#1}}}
    \esoytalque{\entreparentesis{\pfin{\can{#1}}{\exn{#2}}}}{#3}
}}
\newcommand{\paratodocanenexnimplicaque}[3]{\ensuremath{
    \entreparentesis{\paratodocasode{\can{#1}}}
    \esoimplicaque{\entreparentesis{\pfin{\can{#1}}{\exn{#2}}}}{#3}
}}
\newcommand{\paratodocanencanque}[3]{\ensuremath{
    \entreparentesis{\paratodocasode{\can{#1}}}
    \esoytalque{\entreparentesis{\pfin{\can{#1}}{\can{#2}}}}{#3}
}}
\newcommand{\paratodoexnencanque}[3]{\ensuremath{
    \entreparentesis{\paratodocasode{\exn{#1}}}
    \esoytalque{\entreparentesis{\pfin{\exn{#1}}{\can{#2}}}}{#3}
}}
\newcommand{\paratodoexnenexnque}[3]{\ensuremath{
\entreparentesis{\paratodocasode{\exn{#1}}}
\esoytalque{\entreparentesis{\pfin{\exn{#1}}{\exn{#2}}}}{#3}
}}
\newcommand{\paratodocanenexn}[2]{\ensuremath{
    \entreparentesis{\paratodocasode{\can{#1}}}
    \entreparentesis{\pfin{\can{#1}}{\exn{#2}}}
}}
\newcommand{\paratodocanencan}[2]{\ensuremath{
    \entreparentesis{\paratodocasode{\can{#1}}}
    \entreparentesis{\pfin{\can{#1}}{\can{#2}}}
}}
\newcommand{\paratodoexnencan}[2]{\ensuremath{
    \entreparentesis{\paratodocasode{\exn{#1}}}
    \entreparentesis{\pfin{\exn{#1}}{\can{#2}}}
}}
\newcommand{\paratodoexnenexn}[2]{\ensuremath{
    \entreparentesis{
        \paratodocasode{\exn{#1}}
    }
    \entreparentesis{
        \pfin{\exn{#1}}{\exn{#2}}
    }
}}
\newcommand{\paraninguncasoen}[2]{\ensuremath{
        \paratodocaso\pfnotin{#1}{#2}
}}
\newcommand{\pfgensig}[1]{\ensuremath{\gensig{\left({#1}\right)}}}
\newcommand{\set}[1]{\ensuremath{\left\lbrace{#1}\right\rbrace}}
\newcommand {\setCompr}[2]{
	\ensuremath{
		\left\lbrace{ 
			#1 {\;\big{\vert}\;} #2
        }\right\rbrace
    }
}

\newcommand{\extset}[2]{\ensuremath{
        \left\lbrace{#1}\,\vert\,{#2}\right\rbrace
}}
\newcommand{\ext}[1]{\ensuremath{\left\lbrace{#1}\right\rbrace}}
\newcommand{\extdext}[1]{\ensuremath{\left\lbrace{\exn{#1}}\right\rbrace}}
\newcommand{\predext}[2]{\ensuremath{\left\lbrace{#1}\,∣\,{#2}\right\rbrace}}
%\newcommand{\rfrac}[2]{\ensuremath{{{}^{#1}}{\!}{{/{\!}}_{#2}}}}
\DeclareMathOperator{\lcm}{lcm}
\newcommand{\norm}[1]{\ensuremath{\left\Vert#1\right\Vert}}
\newcommand{\abs}[1]{\ensuremath{\left\vert#1\right\vert}}
\newcommand{\BX}{\ensuremath{\bfs{B}(X)}}
\newcommand{\A}{\ensuremath{\mathcal{A}}}
\newcommand{\setN}{\ensuremath{\mathds{N}}}
\newcommand{\setR}{\ensuremath{\mathds{R}}}
\newcommand{\setQ}{\ensuremath{\mathds{Q}}}
\newcommand{\setZ}{\ensuremath{\mathbb{Z}}}
\newcommand{\setC}{\ensuremath{\mathbb{C}}}
\newcommand{\setH}{\ensuremath{\mathbb{H}}}
\newcommand{\setO}{\ensuremath{\mathbb{O}}}
\newcommand{\setB}{\ensuremath{\mathbb{B}}}
\newcommand{\algB}{\ensuremath{\mathscr{B}}}
\newcommand{\relR}{\ensuremath{\mathscr{R}}}
\newcommand{\ntna}{\ensuremath{\bfs{\nta}}}
\newcommand{\ntnb}{\ensuremath{\bfs{\ntb}}}
\newcommand{\ntnc}{\ensuremath{\bfs{\ntc}}}
\newcommand{\vntna}{\ensuremath{\bfs{\vnta}}}
\newcommand{\vntnb}{\ensuremath{\bfs{\vntb}}}
\newcommand{\vntnc}{\ensuremath{\bfs{\vntc}}}
\makeatother
\makeatletter
% ANILLOS Y \'ALGEBRAS
\renewcommand{\norm}[1]{\ensuremath{\left\Vert#1\right\Vert}}
\renewcommand{\abs}[1]{\ensuremath{\left\vert#1\right\vert}}
\newcommand{\card}[1]{\ensuremath{\abs{#1}}}
\newcommand{\conj}[1]{\ensuremath{\left\lbrace{#1}\right\rbrace}}
\newcommand{\extconj}[2] {\ensuremath{
        \left\lbrace{#1}\,∣\,{#2}\right\rbrace
}}
\newcommand{\seps}{\ensuremath{\varepsilon }}
\newcommand{\To}{\ensuremath{         \longrightarrow         }}
\newcommand{\setfunct}[2]{\ensuremath{\mathbf{#1}{\left({#2}\right)}}}
\newcommand{\bcomplement}[1]{\ensuremath{\complement\left[{#1}\right]}}
\newcommand{\NullaryFunct}[1]{\ensuremath{#1\left(\right)\tighht}}
\newcommand{\UnaryFunct}[2]{\ensuremath{#1\left(#2\right)}}
\newcommand{\Idealz}{\ensuremath{\left(\mathfrak{0}\right)}}
\newcommand{\BinaryFunct}[3]{\ensuremath{#1\left(#2,#3\right)}}
\newcommand{\ThreearyFunct}[4]{\ensuremath{#1\left(#2,#3,#4\right)}}
\newcommand{\FouraryFunct}[5]{\ensuremath{#1\left(#2,#3,#4,#5\right)}}
\newcommand{\RingLocalizedOn}[2]{\ensuremath{
        \NullaryFunct{{#1}_{\Ideal{#2}}}
}}
\newcommand{\Ideal}[1]{\ensuremath{\mathfrak{#1}}}
\newcommand{\Idealx}[2]{\ensuremath{{{\Ideal{#1}}_{#2}}}}
\newcommand{\IdealGeneratorWithSet}[1] {\ensuremath{
        \left\langle\left\lbrace{#1}\right\rbrace\right\rangle
}}
\newcommand{\IdealGeneratorWithExtSet}[2]
{\ensuremath{
        \left\langle{\left\lbrace{#1}\,|\,{#2}\right\rbrace}\right\rangle
}}
\newcommand{\IdealGenerator}[1] {\ensuremath{
        \left\langle{#1}\right\rangle
}}
\newcommand{\IdealGeneratorExt}[2] {\ensuremath{
        \left\langle{{#1}\,|\,{#2}}\right\rangle
}}
\newcommand{\idealof}[2]{\ensuremath{\Ideal{#1}{\unlhd}\,{#2}}}
\newcommand{\primeidealof}[2]{\ensuremath{\Ideal{#1}{\lhd}_{prime}}{#2}}
\newcommand{\maximalidealof}[2]{\ensuremath{\Ideal{#1}{\lhd}_{max}}{#2}}
\newcommand{\idealofx}[3]{\ensuremath{\Idealx{#1}{#2}{\unlhd}\,{#3}}}
\newcommand{\primeidealofx}[3]{\ensuremath{\Idealx{#1}{#2}{\lhd}_{prime}}{#3}}
\newcommand{\maximalidealofx}[3]{\ensuremath{\Idealx{#1}{#2}{\lhd}_{max}}{#3}}
\newcommand{\Krullht}{\ensuremath{\textbf{ht}}}
\newcommand{\KrullHeight}[1]{\ensuremath{\Krullht\left({#1}\right)}}
\newcommand{\Kht}[1]{\ensuremath{\KrullHeight{\Ideal{#1}}}}
\newcommand{\Khtx}[2]{\ensuremath{\Krullht\left(\Idealx{#1}{#2}\right)}}
\newcommand{\ChainOfIdealsOnV}[2]{\ensuremath{\mathcal{#1}\left({#2}\right)}}
\newcommand{\ChainOfIdealsOn}[3] {\ensuremath{\mathcal{#1}\left({#2},\Ideal{#3}\right)}}
\newcommand{\ChainOfPrimeIdealsOn}[3] {\ensuremath{\mathcal{#1}_{prime}\left({#2},\Ideal{#3}\right)}}
\newcommand{\IsAChain}[3] {\ensuremath{\textsc{Chain}\left\lbrace \ChainOfIdealsOn{#1}{#2}{#3}\right\rbrace}}
\newcommand{\AscendentChain} {\ensuremath{\textsc{AscChain}}}
\newcommand{\AscendentChainV} {\ensuremath{\AscendentChain}}
\newcommand{\IsAnAscendentChainV}[2] {\ensuremath{\AscendentChainV\left\lbrace \ChainOfIdealsOnV{#1}{#2}\right\rbrace}}
\newcommand{\IsAnAscendentChain}[3] {\ensuremath{\AscendentChain\left\lbrace \ChainOfIdealsOn{#1}{#2}{\Ideal{#3}}\right\rbrace}}
\newcommand{\IsAnAscendentChainx}[4] {\ensuremath{\AscendentChain\left\lbrace \ChainOfIdealsOn{#1}{#2}{\Idealx{#3}{#4}}\right\rbrace}}
\newcommand{\IsAPrimeChain}[3] {\ensuremath{\textsc{Chain}_{prime} \left\lbrace\mathcal{#1}\left({#2},\Ideal{#3}\right)\right\rbrace}}
\newcommand{\IsAPrimeChainV}[2] {\ensuremath{\textsc{Chain}_{prime} \left\lbrace\mathcal{#1}\left({#2}\right)\right\rbrace}}
\newcommand{\FamilyOfChainsOf}[2] {\ensuremath{\mathcal{F}\mathcal{C}_{\left({#1},{\Ideal{#2}}\right)}}}
\newcommand{\FamilyOfChainsOfV}[1] {\ensuremath{\mathcal{F}\mathcal{C}_{\left({#1}\right)}}}
\newcommand{\FamilyOfPrimeChainsOf}[2] {\ensuremath{
        \mathcal{F}
        \mathcal{C}_{
            \left({#1},{\Ideal{#2}}\right)}^{prime}
}}
\newcommand{\FamilyOfPrimeChainsOfV}[1]{\ensuremath{
        \mathcal{F}
        \mathcal{C}_{\left({#1}\right)}^{prime}
}}
\newcommand{\PowerSetOf}[1]{\ensuremath{
        \boldmath{\mathfrak{P}}
        \left({#1}\right)
}}
\newcommand{\IdealsOf}[1]{\ensuremath{\mathcal{I}{deals}\left({#1}\right)}}
\newcommand{\SpecOf}[1]{\ensuremath{\mathcal{S}{pec}\left({#1}\right)}}
\newcommand{\SpecOnOf}[2]{\ensuremath{\mathcal{V}_{#2}\left({#1}\right)}}
\newcommand{\MaxSpecOf}[1]{\ensuremath{\mathcal{M}{ax}\left({#1}\right)}}
\newcommand{\MaxSpecOnOf}[2]{\ensuremath{\mathcal{M}_{#2}\left({#1}\right)}}
\newcommand{\Then}{\ensuremath{          \Rightarrow          }}
\newcommand{\If}{\ensuremath{          \Leftarrow          }}
\newcommand{\Iff}{\ensuremath{\Leftrightarrow }}
\newcommand{\sbneg}{\ensuremath{∼}}
\newcommand{\IsAnUCRing}[1] {\ensuremath{
        \textsc{UCRng}
        \left\lbrace
            {#1}
        \right\rbrace
}}
\newcommand{\IsACRing}[1]{\ensuremath{
        \textsc{CRng}
        \left\lbrace{#1}\right\rbrace
}}
\newcommand{\IsARing}[1] {\ensuremath{
        \textsc{Rng}
        \left\lbrace{#1}\right\rbrace
}}
\newcommand{\IsAField}[1]{\ensuremath{
        \textsc{Field}\left\lbrace{#1}\right\rbrace
}}
\newcommand{\IsAGroup}[1] {\ensuremath{
        \textsc{Grp}
        \left\lbrace{#1}\right\rbrace
}}
\newcommand{\IsACGroup}[1] {\ensuremath{
        \textsc{AbGrp}
        \left\lbrace{#1}\right\rbrace
}}
\newcommand{\IsAModuleOver}[2] {\ensuremath{
        {#2}-\textsc{Mod}
        \left\lbrace{#1}\right\rbrace
}}
\newcommand{\IsALinearVectorialSpaceOver}[2] {\ensuremath{
        \textsc{LinVect}_{#2}
        \left\lbrace{#1}\right\rbrace
}}
\newcommand{\IsAnAlgebraOver}[3] {\ensuremath{
        \textsc{Alg}_{\IsAModuleOver{#2}{#3}}
        \left\lbrace{#1}\right\rbrace
}}
\newcommand{\FamilyOfAllChainsOf}[1] {\ensuremath{
        \mathcal{F}\mathcal{F}\mathcal{C}\left({#1}\right)
}}
\newcommand{\FamilyOfAllPrimeChainsOf}[1] {\ensuremath{
        {\mathcal{F}\mathcal{F}\mathcal{C}}^{prime}
        \left({#1}\right)
}}
\newcommand{\IsNoether}[1] {\ensuremath{
        \textsc{Noether}
        \left\lbrace{#1}\right\rbrace
}}
\newcommand{\IsNoetherian}[1] {\ensuremath{
        \textsc{Noether}
        \left\lbrace{#1}\right\rbrace
}}
\newcommand{\dimK}[1] {\ensuremath{
        {\textbf{dim}}_{\textsf{Krull}}
        \left(R\right)
}}
\newcommand{\lo}{\ensuremath{ \vee}}
\newcommand{\ly}{\ensuremath{         \wedge         }}
\newcommand{\Union}{\ensuremath{\bigcup}}
\newcommand{\DisjointUnion}{\ensuremath{\biguplus}}
\newcommand{\Intersection}{\ensuremath{\bigcap}}
\newcommand{\DirectSum}{\ensuremath{\bigoplus}}
\newcommand{\TensorProduct}{\ensuremath{\bigotimes}}
\newcommand{\FractionsOfRngOn}[2]{\ensuremath{\RingLocalizedOn{#1}{#2}}}
\newcommand{\NaturalInclusion}[1]{\ensuremath{{\imath}_{\Ideal{#1}}}}
\newcommand{\NaturalInclusionMapsTo}[2] {\ensuremath{
        \UnaryFunct{{\imath}_{\Ideal{#1}}}
        {#2}
}}
\newcommand{\KernelOfNaturalInclusion}[1] {\ensuremath{
        \UnaryFunct{\Ideal{K}}
        {\Ideal{#1}}
}}
\newcommand{\KernelOf}[1]{\ensuremath{
        \UnaryFunct{\textbf{\ker}}
        {#1}
}}
\newcommand{\length}[1]{\ensuremath{\mathbf{length}{\left({#1}\right)}}}
\newcommand{\IsADomain}[1]{\ensuremath{\textsc{Dom}{\lbrace{#1}\rbrace}}}
\newcommand{\IsSubNoetherian}[1]{\ensuremath{
        \textsc{SubNoether}
        {\lbrace{#1}\rbrace}
}}
\newcommand{\IsFG}[2]{\ensuremath{{
            {\boldsymbol{\textsc{fg}}}_{{#2}
            \textsf{-mod}}}{\lbrace{#1}\rbrace}
}}
\newcommand{\IsASet}[1]{\ensuremath{
        {{\textsc{Set}}
        {\lbrace{#1}\rbrace}}
}}
\newcommand{\GPolynomialRingOfOnOver}[3]{\ensuremath{
        {#1}
        \left[
            {
            \lbrace{
                {\mathtt{#2}}_{\imath}
                }
            \rbrace
            }_{
                \imath          \in {\mathcal{#3}}
                }
        \right]
}}
\newcommand{\quotientofideals}[2]{\ensuremath{
        \rfrac{\Ideal{#2}}{\Ideal{#1}}
}}
\newcommand{\quotientringbyideal}[2]{\ensuremath{
        \rfrac{{#2}}{\Ideal{#1}}
}}
\newcommand{\TuplesOfNearbyPrimes}[2]{\ensuremath{
        \mathfrak{#2}{-}\mathcal{N}{ear}\mathcal{P}{rimes}
        {\left({#1}\right)}
}}
\newcommand{\SubChain}[3]{\ensuremath{
        \overline{{\mathcal{#1}}_{\Ideal{#2}}^{\Ideal{#3}}}
}}
\newcommand{\myfrac}[2]{\ensuremath{
        \genfrac{}{}{0pt}{0}{#1}{#2}
}}
\newcommand{\MakeFiniteChains}[3]
{
    \ensuremath{
		\textsc{FiniteChains}
			{\lbrace{
				{#1}_{\Ideal{#2}}^{\Ideal{#3}}
			}\rbrace}
	}
}
\newcommand{\AnilloGenerado}{\ensuremath{\boldsymbol{\mathcal{A}}}}
\newcommand{\AnilloGeneradoPor}[2]{\ensuremath{
        {\boldsymbol{\mathcal{A}}}_{\left(#1,\mathcal{#2}\right)}
}}
\newcommand{\SetOfSetsOfIdealsOn}[2]{
	\ensuremath{
		\boldsymbol{\mathcal{P}}
		\left(
			\IdealsOf{\left({\AnilloGeneradoPor{#1}{#2}}\right)}
		\right)
	}
}
\newcommand{\bmid}{\ensuremath{\boldsymbol{\;∣\;}}}
\newcommand{\ascchain}[1]{\ensuremath{{\boldsymbol{\mathcal{#1}}}}}
\newcommand{\TOrderWMin}{\ensuremath{
    \underset{{with\,\mathcal{M}{in}}}
    {\mathcal{L}{in}\mathcal{O}{rder}}
}}
\newcommand{\SetOfAscChainsOn}[1]{\ensuremath{
        \mathcal{A}\mathcal{C}\left({#1}\right)
}}
\newcommand{\OrderOn}[3]
{
	\ensuremath
	{
		\left(
			{
				{#1},{#2},{#3}
			}
		\right)
	}
}
\newcommand{\OrderOnIndexes}[2]{\ensuremath{
        \left({\mathcal{#1},{#2}_{0},\leq}\right)
}}
\newcommand{\QueCumplaCon}{\ensuremath{
        \,\boldsymbol{:}\,
}}
\newcommand{\OrderOnIdealsChain}[2]{\ensuremath{
        \left(\ascchain{#1},{\mathfrak{#1}}_{{#2}_{0}},\subseteq \right)
}}
\newcommand{\OrderIsomorphism}{\ensuremath{
        {\cong}_{\TOrderWMin}
}}
\newcommand{\sumab}[3]{\ensuremath{
        \overset{#3}{\underset{{#1}={#2}}{\sum}}
}}
\newcommand{\diff}[2]{\ensuremath{\mathbf{d}{{#1}^{#2}}}}
\newcommand{\tangentfibered}[1]{\ensuremath{
        \mathbf{T}{\left({#1}\right)}
}}
\newcommand{\tangentspace}[2]{\ensuremath{
        {{\mathbf{T}}_{\mathbf{#1}}}{\left({#2}\right)}
}}
\newcommand{\dualtangentspace}[2]{\ensuremath{
        {{\mathbf{T}}_{\mathbf{#1}}}^{\boldsymbol{*}}
        {\left({#2}\right)}
}}
\newcommand{\diffof}[3]{\ensuremath{
        \frac{\boldsymbol{\partial}{#1}}
        {\boldsymbol{\partial}{{#2}^{#3}}}
}}
\newcommand{\esfera}[1]{\ensuremath{\mathbb{S}^{#1}}}
\newcommand{\bola}[1]{\ensuremath{\mathbb{B}^{#1}}}
\newcommand{\espacioproyectivoreal}[1]{\ensuremath{{
            \mathbb{R}\mathbb{P}}^{#1}
}}
\newcommand{\espacioproyectivocomplejo}[1]{\ensuremath{
        {\mathbb{C}\mathbb{P}}^{#1}
}}
% TEOR\'IA DE GRUPOS FINITOS
\newcommand{\MW}[1]{\ensuremath{{\mathsf{MW}}_{#1}}}
\newcommand{\nil}[1]{\ensuremath{{\mathtt{nil}}-{#1}}}
\newcommand{\setP}{\ensuremath{}\mathbb{P}}
\newcommand{\setPgt}[1]{\ensuremath{{\setP}_{>{#1}}}}
\newcommand{\setNgt}[1]{\ensuremath{{\setN}_{>{#1}}}}
\newcommand{\generadoporZitems}{\ensuremath{
        \entreangulos{\pbm{1}}
}}
\newcommand{\generadoporUitem}[1]{\ensuremath{
        \entreangulos{#1}
}}
\newcommand{\generadopor}[1]{\ensuremath{
        \generadoporUitem{#1}
}}
\newcommand{\generadoporDitems}[2]{\ensuremath{
        \entreangulos{#1,#2}
}}
\newcommand{\generadoporTitems}[3]{\ensuremath{
        \entreangulos{ #1,#2,#3}
}}
\newcommand{\generadoporNitems}[3]{\ensuremath{
        \entreangulos{#1,#2,\ldots,#3}
}}
\newcommand{\grA}{\ensuremath{\mathbf{A}}}
\newcommand{\grB}{\ensuremath{\mathbf{B}}}
\newcommand{\grC}{\ensuremath{\mathbf{C}}}
\newcommand{\grD}{\ensuremath{\mathbf{D}}}
\newcommand{\grE}{\ensuremath{\mathbf{E}}}
\newcommand{\grF}{\ensuremath{\mathbf{F}}}
\newcommand{\grG}{\ensuremath{\mathbf{G}}}
\newcommand{\grH}{\ensuremath{\mathbf{H}}}
\newcommand{\grI}{\ensuremath{\mathbf{I}}}
\newcommand{\grJ}{\ensuremath{\mathbf{J}}}
\newcommand{\grK}{\ensuremath{\mathbf{K}}}
\newcommand{\grN}{\ensuremath{\mathbf{N}}}
\newcommand{\grM}{\ensuremath{\mathbf{M}}}
\newcommand{\grL}{\ensuremath{\mathbf{L}}}
\newcommand{\grP}{\ensuremath{\mathbf{P}}}
\newcommand{\grQ}{\ensuremath{\mathbf{Q}}}
\newcommand{\grR}{\ensuremath{\mathbf{R}}}
\newcommand{\grS}{\ensuremath{\mathbf{S}}}
\newcommand{\grT}{\ensuremath{\mathbf{T}}}
\newcommand{\grU}{\ensuremath{\mathbf{U}}}
\newcommand{\subgr}{\ensuremath{<}}
\newcommand{\subgreq}{\ensuremath{\leqslant}}
\newcommand{\charsubgr}{\ensuremath{\stackrel{\mathtt{char}}{\_}}}
\newcommand{\charsubgreq}{\ensuremath{\stackrel{\mathtt{char}}{{}_{=}}}}
\newcommand{\normsubgr}{\ensuremath{         \vartriangleleft         }}
\newcommand{\normsubgreq}{\ensuremath{\trianglelefteq}}
\newcommand{\setofsets}[1]{\ensuremath{\mathcal{#1}}}
\newcommand{\biy}[1]{\ensuremath{{\mathtt{Biy}}{\left({#1}\right)}}}
\newcommand{\biyuno}[1]{\ensuremath{
        {{\mathtt{Biy}}_{1\mapsto1}}
        {\left({#1}\right)}
}}
\newcommand{\graut}[1]{\ensuremath{{\mathtt{Aut}}_{\mathtt{Grp}}{\left({#1}\right)}}}
\newcommand{\grinn}[1]{\ensuremath{{\mathtt{Inn}}_{\mathtt{Grp}}{\left({#1}\right)}}}
\newcommand{\grout}[1]{\ensuremath{{\mathtt{Out}}_{\mathtt{Grp}}{\left({#1}\right)}}}
\newcommand{\grendo}[1]{\ensuremath{{\mathtt{End}}_{\mathtt{Grp}}{\left({#1}\right)}}}
\newcommand{\grmonoendo}[1]{\ensuremath{
        {{\mathtt{End}}_{\mathtt{Mono}}}^{\mathtt{Grp}}{\left({#1}\right)}
}}
\newcommand{\grhom}[2]{\ensuremath{
        {{\mathtt{Hom}}_{Grp}}
        {\left({#1,#2}\right)}
}}
\newcommand{\griso}[2]{\ensuremath{
        {{\mathtt{Iso}}_{Grp}}
        {\left({#1,#2}\right)}
}}
\newcommand{\grbiy}[2]{\ensuremath{
        {\mathtt{Biy}}
        {\left({#1,#2}\right)}
}}
\newcommand{\grlattice}[1]{\ensuremath{
        {\mathcal{L}}_{#1}
}}
\newcommand{\pgrlattice}[2]{\ensuremath{
        {\mathcal{P}}_{#2}
        \entreparentesis{#1}
}}
\newcommand{\pnormgrlattice}[2]{\ensuremath{
        {\mathcal{N}}_{#2}
        \entreparentesis{#1}
}}
\newcommand{\psylowlattice}[2]{\ensuremath{
        {{\mathcal{S}\mathsf{yl}}_{#2}}
        \entreparentesis{#1}
}}
\newcommand{\grcoc}[2]{\ensuremath{
        \nicefrac{#1}{#2}
}}
\newcommand{\normgrlattice}[1]{\ensuremath{{\mathcal{N}}_{#1}}}
\newcommand{\chargrlattice}[1]{\ensuremath{{\mathcal{C}har}_{#1}}}
\newcommand{\maxgrlattice}[1]{\ensuremath{{\mathcal{M}ax}_{#1}}}
\newcommand{\mingrlattice}[1]{\ensuremath{{\mathcal{M}in}_{#1}}}
\newcommand{\supcoregrlattice}[2]{\ensuremath{
        \overline{\mathcal{K}_{#1}}
        {\left({#2}\right)}
}}
\newcommand{\subcoregrlattice}[2]{\ensuremath{
        \underline{\mathcal{K}_{#1}}
        {\left({#2}\right)}
}}
\newcommand{\pSyllattice}[2]{\ensuremath{
        {{\mathcal{S}yl}_{#2}}
        {\left({#1}\right)}
}}
\newcommand{\grclass}[1]{\ensuremath{\mathsf{#1}}}
\newcommand{\cyclicgr}[1]{\ensuremath{{\grclass{C}}_{#1}}}
\newcommand{\trivialgr}{\ensuremath{{\grclass{C}}_{\mathbf{1}}}}
\newcommand{\pd}[2]{\ensuremath{{#1}\!          \times          \!{#2}}}
\newcommand{\pdpot}[2]{\ensuremath{{\left({#1}\right)}^{#2}}}
\newcommand{\psd}[2]{\ensuremath{{#1}\!\rtimes\!{#2}}}
\newcommand{\psdef}[3]{\ensuremath{{#1}\!{\rtimes}_{#3}\!{#2}}}
\newcommand{\pactsd}[3]{\ensuremath{{#1}\!{\rtimes}_{#3}\!{#2}}}
\newcommand{\pxd}[2]{\ensuremath{{#1}\!\bowtie\!{#2}}}
\newcommand{\pactxd}[4]{\ensuremath{{#1}\! {_{#3\!}}{\bowtie}_{#4}\!{#2}}}
\newcommand{\pwr}[2]{\ensuremath{{#1}\wr{#2}}}
\newcommand{\pw}[2]{\ensuremath{{#1}\wr{#2}}}
\newcommand{\pdcyclics}[2]{\ensuremath{\pd{\cyclicgr{#1}}{\cyclicgr{#2}}}}
\newcommand{\pdtrescyclics}[3]{\ensuremath{{\cyclicgr{#1}}\!         \times         \!{\cyclicgr{#2}}\!         \times         \!{\cyclicgr{#3}}}}
\newcommand{\pdquatrocyclics}[4]{\ensuremath{{\cyclicgr{#1}}\!         \times         \!{\cyclicgr{#2}}\!         \times         \!{\cyclicgr{#3}}\!         \times         \!{\cyclicgr{#4}}}}
\newcommand{\psdcyclics}[2]{\ensuremath{\psd{\cyclicgr{#1}}{\cyclicgr{#2}}}}
\newcommand{\psdefcyclics}[3]{\ensuremath{\psdef{\cyclicgr{#1}}{\cyclicgr{#2}}{#3}}}
\newcommand{\pxdcyclics}[2]{\ensuremath{\pxd{\cyclicgr{#1}}{\cyclicgr{#2}}}}
\newcommand{\pwcyclics}[2]{\ensuremath{\pw{\cyclicgr{#1}}{\cyclicgr{#2}}}}
\newcommand{\dihgr}[1]{\ensuremath{{\grclass{Dih}}_{#1}}}
\newcommand{\dicgr}[1]{\ensuremath{{\grclass{Dic}}_{#1}}}
\newcommand{\gdihgr}[2]{\ensuremath{{\grclass{GDih}}_{#1,#2}}}
\newcommand{\sdihgr}[2]{\ensuremath{{\grclass{SDih}}_{#1,#2}}}
\newcommand{\gdicgr}[2]{\ensuremath{{\grclass{GDic}}_{#1,#2}}}
\newcommand{\simgr}[2]{\ensuremath{{\grclass{Sym}}_{#1,#2}}}
\newcommand{\altgr}[2]{\ensuremath{{\grclass{Alt}}_{#1,#2}}}
\newcommand{\gKleingr}[3]{\ensuremath{{\grclass{gK}}_{#1,#2}^{#3}}}
\newcommand{\Egr}[2]{\ensuremath{{\grclass{E}}_{#1}^{#2}}}
\newcommand{\quaterniongr}[1]{\ensuremath{{\grclass{Q}}_{#1}}}
\newcommand{\squaterniongr}[1]{\ensuremath{{\grclass{SQ}}_{#1}}}
\newcommand{\gquaterniongr}[1]{\ensuremath{{\grclass{GQ}}_{#1}}}
\newcommand{\normalizador}[2]{\ensuremath{
        {{\mathsf{N}}_{#1}}
        {\left({#2}\right)}
}}
\newcommand{\centralizador}[2]{\ensuremath{
        {{\mathsf{C}}_{#1}}
        {\left({#2}\right)}
}}
\newcommand{\centro}[1]{\ensuremath{
        {\mathsf{Z}}
        {\left({#1}\right)}
}}
\newcommand{\centronum}[2]{\ensuremath{
        {{\mathsf{Z}}_{#2}}
        {\left({#1}\right)}
}}
\newcommand{\frattini}[1]{\ensuremath{{\mathbf{\Phi }}_{#1}}}
\newcommand{\fitting}[1]{\ensuremath{{\bm{\mathsf{F}}}_{#1}}}
\newcommand{\gfitting}[1]{\ensuremath{{\bm{\mathsf{F}^{*}}}_{#1}}}
\newcommand{\indice}[2]{\ensuremath{\entrecorchetes{{#1}\;:\;{#2}}}}
\newcommand{\grderivado}[2]{\ensuremath{
        {#1}^{\entreparentesis{#2}}
}}
\newcommand{\grLderivado}[2]{\ensuremath{
        {{\mathsf{L}}_{\entreparentesis{#2}}}
        {\entreparentesis{#1}}
}}
\newcommand{\grabelianizado}[1]{\ensuremath{
        {#1}^{\mathsf{ab}}
}}
\newcommand{\conmutador}[2]{\ensuremath{
        \entrecorchetes{#1,#2}
}}
\newcommand{\grconmutador}[2]{\ensuremath{
        \generadopor{
            \entrecorchetes{#1,#2}
        }
}}
\newcommand{\supcore}[2]{\ensuremath{
        {\overline{{\bm{\mathsf{Core}}}_{#1}}}
        {\left({#2}\right)}
}}
\newcommand{\subcore}[2]{\ensuremath{
        {\underline{{\bm{\mathsf{Core}}}_{#1}}}
        {\left({#2}\right)}
}}
\newcommand{\neutromult}{\ensuremath{\mathbf{u}}}
\newcommand{\unity}{\ensuremath{\neutromult}}
\newcommand{\neutroadd}{\ensuremath{\mathbf{e}}}
\newcommand{\zero}{\ensuremath{\neutroadd}}
\newcommand{\gfrac}[2]{\ensuremath{\genfrac{}{}{0pt}{0}{#1}{#2}}}
\newcommand{\leqB}{\ensuremath{{\leq}_{\algB}}}
\newcommand{\geqB}{\ensuremath{{\geq}_{\algB}}}
\newcommand{\lneqB}{\ensuremath{{\lneq}_{\algB}}}
\newcommand{\gneqB}{\ensuremath{{\gneq}_{\algB}}}
\newcommand{\nleqB}{\ensuremath{{\nleq}_{\algB}}}
\newcommand{\ngeqB}{\ensuremath{{\ngeq}_{\algB}}}

\setcounter{secnumdepth}{4}
\setcounter{tocdepth}{4}
\synctex=1
\newtheorem{thm}{Teorema}[chapter]
\newtheorem{example}[thm]{Ejemplo}
\newtheorem{exercise}[thm]{Ejercicio}
\newtheorem{axiom}[thm]{Axioma}
\newtheorem{cor}[thm]{Corolario}
\newtheorem{lem}[thm]{Lema}
\newtheorem{prop}[thm]{Proposición}
\theoremstyle{definition}
\newtheorem{defn}[thm]{Definición}
\theoremstyle{remark}
\newtheorem{rem}[thm]{Nota}
\newtheorem{fixed}[thm]{Fijación de Notación}
\newtheorem{formalproof}[thm]{Demostración}
\newtheorem{solution}[thm]{Solución}
\makeatother
\makeatletter
%\bibliography{bibliografia}
%\usepackage{polyglossia}
\usepackage{csquotes}
\usepackage{fancyvrb}
\usepackage{listings}
\lstset{ %
    language=C++,                	% choose the language of the code
    basicstyle=\footnotesize,       % the size of the fonts that are used for the code
    numbers=left,                   % where to put the line-numbers
    numberstyle=\footnotesize,      % the size of the fonts that are used for the line-numbers
    stepnumber=1,                   % the step between two line-numbers.
                                    % If it is 1 each line will be numbered
    numbersep=5pt,                  % how far the line-numbers are from
                                    %the code
    backgroundcolor=\color{white},  % choose the background color.
                                    %You must add \usepackage{color}
    showspaces=false,               % show spaces adding
                                    %particular underscores
    showstringspaces=false,         % underline spaces within strings
    showtabs=false,                 % show tabs within strings
                                    %adding particular underscores
    frame=single,           		% adds a frame around the code
    tabsize=2,          			% sets default tabsize to 2 spaces
    captionpos=b,           		% sets the caption-position to bottom
    breaklines=true,        		% sets automatic line breaking
    breakatwhitespace=false,    	% sets if automatic breaks should
                                    %only happen at whitespace
    escapeinside={\%*}{*)}          % if you want to add a
                                    %comment within your code
}

\makeatother

\makeatletter

\begin{document}

\title{
    Matemáticas fundamentales en la Electrónica Digital
}

\author{
    Julián Calderón Almendros
}

\date{\today}

\maketitle


\makeatletter

\hypersetup{pdftitle={Fundamentos matemáticos de la Electrónica Digital},
pdfauthor={Julián Calderón Almendros},
pdfsubject={Álgebras de Boole, Funciones Booleanas, Minimización de Funciones Lógicas, Lógica Secuencial, Autómatas Finitos},
pdfkeywords={Álgebra, Retículos, Lógica, Matemáticas Discretas, Álgebra Computacional}}
\tableofcontents{}
\makeatletter
\listoffigures
\makeatletter
\listoftables
\makeatletter

%\printnomenclature{}
%\makeatletter


\newcommand {\setNz}{\ensuremath{{\setN}_{\mathsf{0}}}}
\newcommand {\setL}{\ensuremath{\mathbb{L}}}
\newcommand {\retL}{\ensuremath{\mathscr{L}}}

\chapter{Introducción}



La idea que persigo al pasar a limpio estos apuntes de clase y estudio no es
la de escribir un libro guía para los fundamentos de la electrónica digital,
finalidad bastante bien cubierta por una infinidad de buenos textos estándar.
Más bien es la de exponer de manera clara -y bastante prolija- todo lo que he
ido aprendiendo sobre estos fundamentos más bien por mi cuenta, el conjunto
de aledaños a este contenido, ya sea por el lado matemático, lógico,
histórico e incluso filosófico. Así que en ningún momento me propongo como
fin el hacer un manual pedagógicamente adecuado, sino más bien poner claridad
en lo que sé, establecer de forma accesible razonamientos difíciles de
encontrar, etc.



\section{Algunos convenios}



\(0 \notin \setN\) Este es el convenio que tomamos para la definición de los
naturales. Es el convenio más común entre matemáticos.

\(\setNz \mathdef \setN \sunion \set{0}\) Es el símbolo que utilizaremos
para los naturales con el elemento \(0\) incluido.

Definiremos los conjuntos \({[0,1]}_{\setQ}\) y \({[0,1]}^{n}_{\setQ}\) dónde
\(n\in\setN\) y \(n > 1\), de forma recursiva:



\begin{align}
&{\left[{0,1}\right]}_{\setQ} \mathdef {\left[{0,1}\right]}\cap{\setQ}\\
&{\left[{0,1}\right]}^{1}_{\setQ} \mathdef {\left[{0,1}\right]}_{\setQ}\\
&{\left[{0,1}\right]}^{n+1}_{\setQ} \mathdef {\left[{0,1}\right]}^{n}_{\setQ}
\times {\left[{0,1}\right]}^{1}_{\setQ} \qquad \forall
n\in{\setN\setminus\set{1}}\\
&\qquad\mbox{\textsf{   dónde    }} \\
&\qquad\qquad\left[0,1\right] 
\mathdef 
\set{{x \in \setR} {\;|\;} {0 \leq x} \textsf{ y } {x \leq 1}}
\end{align}



	Por lo general los elementos de un conjunto los representaremos por
letras minúsculas, tanto del alfabeto griego como del latino, e
indistintamente con o sin subscriptores, que podrán ser a su vez ristras de
letras y/o dígitos decimales, por ejemplo: \(a_3 , b , \gamma_{1547} ,
\delta_{\alpha 10 A}\). Este criterio será válido a excepción que se exprese
explícitamente otro nombre para elementos o conjuntos.



	A las denominaciones de conjunto le corresponderán letras mayúsculas de
los alfabetos griego o latino, y como en el caso anterior podrán llevar o no
como subscriptores cadenas de letras y/o dígitos decimales, por ejemplo:
\(A_3\) , \(B\) , \(\Gamma_{1547}\) , \(\Delta_{{\Pi}10A}\). Un caso especial
que se añadirá para los nombres de conjuntos serán los conjuntos que vengan
definidos por un elemento, pongamos \(e\). Este conjunto se denotará con
alguna acentuación, por ejemplo: \(\widehat{e}\), \(\widetilde{e}\) o
\(\overline{e}\) y se habrá definido anteriormente el significado preciso.
Todo esto será válido si no está explicitado otro mecanismo alternativo de
denotación.



	A veces aparecerá una operación binaria, digamos \(\cdot * \cdot\), dónde
los puntos se sustituyen con los argumentos, esto es, la operación así
expresada es puesta en forma infija. Cuando expresemos \(a * B\), o sea, el
elemento \(a\) operado mediante \(\cdot * \cdot\) (operación infija entre
elementos) con el conjunto \(B\), se tratará de un expresión no originalmente
definida entre argumentos no homogéneos. Se interpretará como \(a*B \mathdef
\conj{a*b{\,|\,} b\in B}\). Igualmente puede ocurrir que esa operación
aparezca entre dos conjuntos, con los que su interpretación habitual es \(A *
B \mathdef \conj{a*b {\,|\,} a\in A \text{ y } b\in B}\).


	El universo en el que nos moveremos será habitualmente \(\setB\), con o
sin algún subscriptor. Si en una expresión no aparece el signo de pertenencia
de un elemento se supone que pertenece a este conjunto universo, de la misma
forma que un conjunto será por defecto subconjunto de este conjunto universo.
En ocasiones trabajaremos explícitamente sobre el conjunto potencia de
de un conjunto \(\textsf{B}\), esto es, sobre \large\( \mathscr{P} \)
\normalsize \( \textsf{B} \).


\chapter{Álgebras de Boole y postulados de Huntington}


Axiomas para un álgebra de Boole
\(\algB \mathdef \left\langle
\setB,
op \set{ bin : \set{ + , * } }
\right\rangle\).

La operación uno-aria que se utilizará frecuentemente (la complementación
buleana), será introducida como definición en los axiomas y no supuesta
inicialmente más que como la existencia de un conjunto de elementos no vacío
asociados a otro que cumplen unos requisitos. Sin embargo esta operación es
tan importante en las álgebras de Boole que la forma más normal de escribirse
la estructura es



\begin{equation}
\algB \mathdef \left\langle{\setB,
op\set{bin:\set{+,*},una:\set{\overline{\phantom{A}}}}}
\right\rangle
\end{equation}



	Además, la distinción entre los dos neutros \(0_{\setB} \neq 1_{\setB}\)
no siempre se requiere. Esto tiene poco impacto, ya que el único álgebra de
Boole con un solo elemento es el álgebra de Boole trivial.



	Nosotros partiremos de unos postulados para el álgebra de Boole muy
naturales, los de Huntington (artículos en 1904, 1933 - este último
desarrolla un conjunto de tres axiomas, uno de ellos específicamente llamado
axioma de Huntington, y no se corresponde con el caso aquí expuesto, en esta
primera formalización-) definen que sea un álgebra de Boole \(\algB\),
\(\left\langle{\setB,\set{+,*}}\right\rangle\).



\section{Axiomas para un álgebra de Boole: Huntington 1904}



	Estos son los axiomas de Huntington de 1904 para definir un álgebra de
Boole. Notar que los primeros axiomas no son postulados por Huntington, sino
que son una expresión de condiciones fundamentales para que los axiomas
propiamente dichos tengan sentido formal.



Estos pre-axiomas son que \(\setB\) es un conjunto, que las operaciones
binarias $+\, ,\, \cdot$ son funciones y que existen los elementos \(0\) y \(1\)
en el conjunto \(\setB\). Más detalladamente:



\begin{axiom}{\textbf{\textsf{H0}}}\label{axm:H0}
\begin{flalign}\label{eqn:H0}
        {\bfs{\qquad\qquad \setB\ es\ un\ conjunto}}
\end{flalign}
\end{axiom}



	\(\setB\) será un conjunto en algún sistema formalizado -por ejemplo en
la axiomatización \textsf{Zermelo-Fraenkel-Skolem}, \textbf{ZFS} o, digamos,
en la más extensa de \textsf{Neumann-Gödel-Bernays}, abreviadamente
\textbf{NGB} -. Aunque bastante formal, este requisito nos quita el peso de
trabajar en clases propias. En principio tratamos con conjuntos
habitualmente, en el instante que tratamos problemas de la matemática que no
sean en álgebra homológica o en teoría de las categorías.



\begin{axiom}{\textbf{\textsf{H1s}}}\label{axm:H1s}
	La operación binaria interna suma \(+\) es función:
\begin{flalign}\label{eqn:H1s}
        \bfs{ \qquad \cdot \qquad + \qquad \cdot \qquad : \qquad
        \setB\times\setB \quad \fromto \qquad \setB}\\
        \qquad \forall (x,y) \in \setB\times\setB \quad \exists! z \in \setB
        \qquad x+y=z
\end{flalign}
\end{axiom}



	Siempre que se pide en una estructura algebraica una operación interna,
lo que se pide es que sea una función que acepta dos elementos del conjunto y
devuelve otro del mismo conjunto. Aquí es importante el hecho que debe ser
una función: la operación de dos elementos siempre debe dar otro elemento del
conjunto. No puede haber una pareja de números que no se puedan operar. En
este caso hemos adoptado el signo \(+\) y llamamos a la operación suma
booleana o lógica, y muchas veces abreviaremos con suma, distinguiendo
siempre esta suma de la suma aritmética.



\begin{axiom}{\textbf{\textsf{H1p}}}\label{axm:H1p}
	La operación binaria interna producto \(*\) es función:
\begin{flalign}\label{eqn:H1p}
        \bfs{ \qquad \cdot \qquad * \qquad \cdot \qquad : \qquad
        \setB\times\setB \quad \fromto \qquad \setB}\\
        \qquad \forall (x,y) \in \setB\times\setB \quad \exists! z \in \setB
        \qquad x * y=z
\end{flalign}
\end{axiom}



	Como en el caso anterior, se pide la existencia de otra operación binaria
interna. Debe ser una función. En este caso hemos adoptado el signo
\(*\) y llamamos a la operación producto booleano o lógico, y muchas
veces abreviaremos con producto, distinguiendo siempre esta suma de la suma
aritmética. Se ha optado por el signo \(*\) para poder utilizar el
signo \(\cdot\) como posición de los operandos, pero el signo natural sería
el de \(\cdot\) en vez de \(*\).



	A continuación vienen los axiomas de existencia de elemento neutro, en
sus dos versiones. Como veréis todos los axiomas tienen dos sabores, uno para
la suma y otro para el producto, un elemento neutro \(0\) para la suma y otro
\(1\) para el producto. Este detalle tendrá un desarrollo posterior muy
interesante: la dualidad de teoremas y expresiones.



\begin{axiom}{\textbf{\textsf{H2s}}}\label{axm:H2s}
\begin{flalign}\label{eqn:H2s}
          \bfs{\qquad \forall x \in \setB \qquad x+0=x}
\end{flalign}
\end{axiom}



\begin{axiom}{\textbf{\textsf{H2p}}}\label{axm:H2p}
\begin{flalign}\label{eqn:H2p}
          \bfs{\qquad \forall x \in \setB \qquad x * 1=x}
\end{flalign}
\end{axiom}



	De nuevo, los axiomas piden la conmutatividad de ambas operaciones
binarias internas.



\begin{axiom}{\textbf{\textsf{H3s}}}\label{axm:H3s}
\begin{flalign}\label{eqn:H3s}
          \bfs{\qquad \forall xy \in \setB \qquad x+y=y+x}
\end{flalign}
\end{axiom}



\begin{axiom}{\textbf{\textsf{H3p}}}\label{axm:H3p}
\begin{flalign}\label{eqn:H3p}
          \bfs{\qquad \forall xy \in \setB \qquad x* y=y* x}
\end{flalign}
\end{axiom}



	También las propiedades distributivas de una operación respecto a la otra
vendrán en tandem. La propiedades distributiva del producto respecto a la
suma tendrá una forma familiar. A la propiedad \emph{distributiva del
producto respecto de la suma} en el sentido inverso de la igualdad la
llamaremos \emph{sacar factor común}:



\begin{axiom}{\textbf{\textsf{H4s}}}\label{axm:H4s}
\begin{flalign}\label{eqn:H4s}
          \bfs{
          \qquad \forall xyz \in \setB 
          \qquad x*(y+z)=(x* y)+(x* z)
          }
\end{flalign}
\end{axiom}



	Propiedad distributiva de la suma respecto al producto. Esta tiene un
aspecto extraño: en aritmética no se cumple. A la propiedad
\emph{distributiva de la suma respecto del producto} en el sentido inverso de
la igualdad la llamaremos \emph{sacar sumando común}:



\begin{axiom}{\textbf{\textsf{H4p}}}\label{axm:H4p}
\begin{flalign}\label{eqn:H4p}
          \bfs{\qquad \forall xyz \in \setB \qquad x+(y*
          z)=(x+y)*(x+z)}
\end{flalign}
\end{axiom}



	Éste es el último axioma del álgebra de Boole. Su forma no pide
condiciones independientes para la suma y el producto, para el \(0\) y el
\(1\), sino que es una forma combinada que pide la existencia para cada
elemento de al menos un elemento que llamaremos complementario.



\begin{axiom}{\textbf{\textsf{H5}}}\label{axm:H5}
\begin{flalign}\label{eqn:H5}
&\bfs{\qquad\forall x \in \setB \quad \exists y \in \setB} \\
&\bfs{\qquad\qquad\qquad\qquad\qquad x+y=1} \\
&\bfs{\qquad\qquad\qquad\qquad\qquad x* y=0}
\end{flalign}
\end{axiom}



	Suelo utilizar este axioma en una forma por completo equivalente, pero un 
poco más conveniente para la forma en que trabajo.



\begin{axiom}{\textbf{\textsf{H5'}}}\label{axm:H5'}
\begin{flalign}\label{eqn:H5'}
&\bfs{\qquad\forall x \in \setB \quad \exists \widehat{x} \subseteq \setB}\\
&\bfs{\qquad\qquad\qquad\qquad\qquad\quad \widehat{x}\neq\emptyset}\\
&\bfs{\qquad\qquad\forall y \in \widehat{x}\qquad\qquad x+y=1}\\
&\bfs{\qquad\qquad\forall y \in \widehat{x}\qquad\qquad x* y=0}
\end{flalign}
\end{axiom}



	En el artículo original de 1904, Huntington pone un axioma sexto, en el
que se asegura que no estamos tratando con el álgebra de Boole trivial,
\(\algB\) tal que \(\exists a,b \in \setB\) que \(a\neq b\).



\begin{axiom}{\textbf{\textsf{H6}}}\label{axm:H6}
\begin{flalign}\label{eqn:H6}
&\bfs{\qquad\exists\; x,y\; \in \setB \quad x\; \neq\; y}
\end{flalign}
\end{axiom}




\section*{}



	Es importante darse cuenta que los axiomas están dispuestos de la
siguiente manera: si un axioma se predica de la suma y/o el elemento cero,
habrá otro que se predique de manera virtualmente idéntica del producto y el
elemento uno. Si el axioma relaciona las dos operaciones, habrá dos, uno de
ellos será una copia del otro cambiando las sumas por productos y los
productos por sumas (axioma de la propiedad distributiva).



	El axioma restante, el de existencia de algún elemento complementario,
también tiene una forma parecida: un complementario ha de cumplir dos
subfórmulas, una implica la suma y el elemento uno, y la otra implica el
producto y el elemento cero.



\textbf{Como conclusión podemos decir que, en los axiomas de Huntington, que
debe cumplir cualquier álgebra de Boole, las fórmulas que se cumplen son
siempre duales, esto es, si tenemos una en la que aparecen sumas, productos y
elementos cero y unidad, si cambiamos sumas por productos y productos por
sumas, elementos ceros por unidades y unidades por ceros, obtenemos un
axioma/teorema/fórmula dual.}



\section{Independencia de los axiomas de Huntington de 1904.}



	Para esto demostrar que los axiomas antes dados son independientes entre
sí, primero daremos un modelo consistente y sencillo, en el que haremos
variaciones y obtendremos modelos (ejemplos) de sistemas que cumplan todos
los axiomas excepto uno de ellos. Si se logra quedará claro que no podemos
deducir el axioma fallido del resto de axiomas que sí que se cumplen: el
axioma fallido es lógicamente independiente del resto de axiomas. Esto es
fácil de conseguir, comenzando con modelos de álgebras con base en un
conjunto de dos elementos \({\setB}_2 = \set{0,1}\). En esto copiamos los
modelos dados por Huntington en su artículo de 1904.



\subsection{Modelo mínimo de álgebra de Boole.}



	Definimos:
\begin{flalign}
{\setB}_2 &\mathdef \set{0,1} = \set{0} \cup \set{1} \\
0 &\mathdef \emptyset \\
1 &\mathdef \set{\emptyset} \cup \set{ \set{ \emptyset } } = \set{ \emptyset,
\set{ \emptyset } } \\
0 &\;\in\; 1 \\
0 &\;\subsetneq\; 1 \\
0 &\;\neq\; 1 \\
\end{flalign}
\begin{flalign}
\bfsf{s} &:  {\setB}_2 \times {\setB}_2{\qquad}\fromto{\qquad} {\setB}_2\\
\end{flalign}
\begin{flalign}
\bfsf{s} &:: \left({0,0}\right) {\qquad} \mapsto {\qquad} 0\\
\bfsf{s} &:: \left({0,1}\right) {\qquad} \mapsto {\qquad} 1\\
\bfsf{s} &:: \left({1,0}\right) {\qquad} \mapsto {\qquad} 1\\
\bfsf{s} &:: \left({1,1}\right) {\qquad} \mapsto {\qquad} 1
\end{flalign}
\begin{flalign}
\bfsf{p} &:  {\setB}_2 \times {\setB}_2{\qquad}\fromto{\qquad} {\setB}_2\\
\end{flalign}
\begin{flalign}
\bfsf{p} &:: \left({0,0}\right) {\qquad} \mapsto {\qquad} 0\\
\bfsf{p} &:: \left({0,1}\right) {\qquad} \mapsto {\qquad} 0\\
\bfsf{p} &:: \left({1,0}\right) {\qquad} \mapsto {\qquad} 0\\
\bfsf{p} &:: \left({1,1}\right) {\qquad} \mapsto {\qquad} 1
\end{flalign}
\begin{flalign}
\bfsf{c} &:  {\setB}_2 {\qquad} \fromto {\qquad} {\setB}_2
\end{flalign}
\begin{flalign}
\bfsf{c} &:: 0 {\qquad} \mapsto {\qquad} 1\\
\bfsf{c} &:: 1 {\qquad} \mapsto {\qquad} 0
\end{flalign}



	El conjunto \({\setB}_2\) es conjunto en \textbf{ZFS}, desde el momento
que la el axioma de unión nos asegura que \({\setB}_2\) es conjunto unión de
dos conjuntos de un solo elemento, \({\setB}_2 \mathdef { \set{0} } \cup {
\set{1} } \). Los elementos \(0 \mathdef \emptyset\) y  \(1 \mathdef \set{
\emptyset, \set{ \emptyset } } \). De nuevo para definir \(1\) como conjunto
necesitamos el axioma de unión (o de pares no ordenados) de \textbf{ZFS},
siendo \(1 \mathdef \set{0}\cup\set{\set{0}}\). Podemos ver que \(0 \cap 1 =
\emptyset \) por lo que \(0 \neq 1\). También \(0 \in 1\). Todo este párrafo
constituye la satisfacción de los requerimientos \ref{axm:H0}.



	Tal como hemos definido nuestro modelo de álgebra de Boole, lo primero
que queda claro es que \(\funcsum\) y \(\funcprod\) son funciones binarias
bien definidas, y esta última es además biyectiva. Luego los axiomas
\ref{axm:H1s} y \ref{axm:H1p} quedan satisfechos. 



	Los axiomas \ref{axm:H2s} y \ref{axm:H2p} (existencia del elemento
neutro) quedan directamente satisfechos por simple inspección de las tablas. 



	Para el axioma \ref{axm:H2s} observamos que \(\funcsumat{0}{0} = 0\) y
\(\funcsumat{0}{1} = \funcsumat{1}{0} = 1\) nos muestra el elemento neutro de
la suma, el elemento \(0\).



	Para el axioma \ref{axm:H2p} observamos que \( \funcprodat{1}{0} =
\funcprodat{0}{1} = 0\) y \(\funcprodat{1}{1} = 1\) nos muestra el elemento
neutro del producto, el elemento \(1\).



	Para los axiomas de conmutatividad solo hay que observar en la definición
de las funciones binarias \(\funcsum,\funcprod : \setB_2 \times \setB_2
\fromto \setB_2\), que \(\funcsumat{0}{1} = \funcsumat{1}{0} = 1\) y se
cumple \ref{axm:H3s}, que \(\funcprodat{0}{1} = \funcprodat{1}{0} = 0\) y se
cumple \ref{axm:H3p}. 



	Para las distribuciones de una operación interna sobre otra elegiré
construir unas tablas que muestren la igualdades necesarias.



	Primero veremos la distribución de la suma sobre el producto, o, lo que
es lo mismo, su inversa, sacar factor común. 



	Lo haremos rellenando y comparando la siguiente tabla de verdad.



	El método de relleno es fácil. Primero rellenamos las tres primeras
columnas, corresponden a las variables independientes, que al ser tres con
dos valores posibles en cada ocasión, son un total de ocho combinaciones
distintas. Esas columnas están marcadas con \(1\) en la última fila (fila que
se ve separada). A continuación rellenamos las columnas correspondientes a
una operación binaria interna de dos variables independientes ya rellenas.
Para rellenarlas, solo tenemos que buscar los valores en las tablas dadas
para las operaciones internas en ka definición de este álgebra de Boole.
Estas están marcadas con un \(2\) en la última línea, la línea separada.
Corresponden con la cuarta columna, la suma \( \funcsumat{ \textsf{b} } {
\textsf{c} } \), y con las columnas siete y ocho: la ocho es un producto
\(\funcprodat{\textsf{a}}{\textsf{c}}\) y la siete otro \( \funcprodat{
\textsf{a} } { \textsf{b} } \). Las columnas centrales, la cinco y la seis
corresponden a todas las valoraciones posibles (y en el mismo orden de valor
de las variables independientes) de las dos expresiones que queremos
comparar, las afirmadas por el postulado \ref{axm:H4s}. Podemos ver que ambas
columnas son idénticas, luego se cumple el citado postulado.



\begin{flalign}
\begin{matrix}
\hline
\textsf{a} & \textsf{b} & \textsf{c} &
\funcsumat{\textsf{b}}{\textsf{c}} &
\funcprodat{\textsf{a}}{\funcsumat{\textsf{b}}{\textsf{c}}} &
\funcsumat{\funcprodat{\textsf{a}} {\textsf{b}}} {\funcprodat{\textsf{a}}
{\textsf{c}}}
&
{\funcprodat{\textsf{a}}{\textsf{b}}}
&
{\funcprodat{\textsf{a}}{\textsf{c}}}
\\
\hline
0 & 0 & 0 & 0 & 0 & 0 & 0 & 0 \\
0 & 0 & 1 & 1 & 0 & 0 & 0 & 0 \\
0 & 1 & 0 & 1 & 0 & 0 & 0 & 0 \\
0 & 1 & 1 & 1 & 0 & 0 & 0 & 0 \\
1 & 0 & 0 & 0 & 0 & 0 & 0 & 0 \\
1 & 0 & 1 & 1 & 1 & 1 & 0 & 1 \\
1 & 1 & 0 & 1 & 1 & 1 & 1 & 0 \\
1 & 1 & 1 & 1 & 1 & 1 & 1 & 1 \\
\hline
1 & 1 & 1 & 2 & 3 & 3 & 2 & 2 \\ 
\hline
\end{matrix}
\end{flalign}



	Por último veremos la distribución del producto sobre la suma, o, lo que
es lo mismo, su inversa, sacar sumando común. El método es en todo idéntico
al seguido para la otra distributiva.



\begin{flalign}
\begin{matrix}
\hline
\textsf{a} & \textsf{b} & \textsf{c} &
\funcprodat{\textsf{b}}{\textsf{c}} &
\funcsumat {\textsf{a}}
           {\funcprodat{\textsf{b}}{\textsf{c}}} &
\funcprodat{\funcsumat{\textsf{a}}{\textsf{b}}}
           {\funcsumat{\textsf{a}}{\textsf{c}}}
&
{\funcsumat{\textsf{a}}{\textsf{b}}}
&
{\funcsumat{\textsf{a}}{\textsf{c}}}
\\
\hline
0 & 0 & 0 & 0 & 0 & 0 & 0 & 0 \\
0 & 0 & 1 & 0 & 0 & 0 & 0 & 1 \\
0 & 1 & 0 & 0 & 1 & 0 & 1 & 0 \\
0 & 1 & 1 & 1 & 1 & 1 & 1 & 1 \\
1 & 0 & 0 & 0 & 1 & 1 & 1 & 1 \\
1 & 0 & 1 & 0 & 1 & 1 & 1 & 1 \\
1 & 1 & 0 & 0 & 1 & 1 & 1 & 1 \\
1 & 1 & 1 & 1 & 1 & 1 & 1 & 1 \\
\hline
1 & 1 & 1 & 2 & 3 & 3 & 2 & 2 \\ 
\hline
\end{matrix}
\end{flalign}




	Para el axioma \ref{axm:H5} observamos que \(\funccomplat{0} = 1\) y
\(\funccomplat{1} = 0\) y por inspección en las tablas vemos que
\(\funcsumat{0}{\funccomplat{0}} = \funcsumat{0}{1} = 1\) y que
\(\funcprodat{0}{\funccomplat{0}} = \funcprodat{0}{1} = 0\) cumpliendo para
el elemento \(0\) se cumple que \(\exists 1 = \funccomplat{0} \in \setB_2 \)
\(\funcsumat{1}{\funccomplat{1}} = \funcsumat{1}{0} = 1\), que coincide con
la ecuación segunda (ecuación: 2.15) de \ref{eqn:H5} y
\(\funcprodat{1}{\funccomplat{1}}=\funcprodat{1}{0} = 0\) que coincide con la
tercera ecuación (ecuación: 2.16) de \ref{eqn:H5}, e igualmente para el
elemento \( 1 \), \(\exists 0 = \funccomplat{1} \in \setB_2 \) \(
\funcsumat{0}{\funccomplat{0}} = \funcsumat{0}{1} = 1\) (ecuación: 2.15) y
\(\funcprodat{0}{\funccomplat{0}} =\funcprodat{0}{1} = 0\) (ecuación: 2.16)
de \ref{eqn:H5}.




	Como las ecuaciones 2.15 y 2.16 se cumplen para \(0\) y \(1\), esto es,
\(\forall x \in \setB_2\) como se requiere en la ecuación 2.14, queda
satisfecho el postulado \ref{axm:H5}.



\subsection{Independencia de la suma está siempre definida.}



Modelo en el que solo falla \ref{axm:H1s}, esto es, que la operación suma
no es operación interna: no es función, pero en el sentido que \(1 + 1 \notin
\setB\).


\[
    \begin{matrix}
        {\cdot + \cdot} & \big{\vert} & {0} & {1} & \big{\vert} \\
        \hline
        {0}             & \big{\vert} & {0} & {1} & \big{\vert} \\
        {1}             & \big{\vert} & {1} & {x} & \big{\vert} \\
        \hline
    \end{matrix}
    \qquad\qquad
    \begin{matrix}
        {\cdot * \cdot} & \big{\vert} & {0} & {1} & \big{\vert} \\
        \hline
        {0}                   & \big{\vert} & {0} & {0} & \big{\vert} \\
        {1}                   & \big{\vert} & {0} & {1} & \big{\vert} \\
        \hline
    \end{matrix}
\]
%% (H1s) 
%% 0 + 0 = 0; 0 + 1 = 1; 1 + 0 = 1; 1 + 1 = x; 
%% 0 * 0 = 0; 0 * 1 = 0; 1 * 0 = 0; 1 * 1 = 1; 



	Al ser \(1 + 1 = x\), dónde \(x\) es ningún elemento de $\setB$, o dicho
de otro modo, \(1 + 1 \notin \setB\), vemos que sigue existiendo el elemento
neutro de la suma, \(0\), que la suma sigue siendo conmutativa y las
distributivas siguen valiendo mientras la suma tenga sentido (mientras no
aparezca \(x\)). El complementario de \(0\) es \(1\) y el de \(1\) es \(0\).



\subsection{Independencia de la unicidad de la definición de la suma.}

Modelo en el que no se cumple \ref{axm:H1s}, en el sentido que \(1 + 1
\mapsto 1\) y \(1 + 1 \mapsto 0\).


\[
    \begin{matrix}
        {\cdot + \cdot} & \big{\vert} & {0} & {1} & \big{\vert} \\
        \hline
        {0}             & \big{\vert} & {0} & {1} & \big{\vert} \\
        {1}             & \big{\vert} & {1} & {x=0 \vee x=1} & \big{\vert} \\
        \hline
    \end{matrix}
    \qquad\qquad
    \begin{matrix}
        {\cdot * \cdot} & \big{\vert} & {0} & {1} & \big{\vert} \\
        \hline
        {0}                   & \big{\vert} & {0} & {0} & \big{\vert} \\
        {1}                   & \big{\vert} & {0} & {1} & \big{\vert} \\
        \hline
    \end{matrix}
\]
%% (H1s) 
%% 0 + 0 = 0; 0 + 1 = 1; 1 + 0 = 1; 1 + 1 \mapsto 0,1; 
%% 0 * 0 = 0; 0 * 1 = 0; 1 * 0 = 0; 1 * 1 = 1; 



Exactamente igual que en la versión anterior, se cumplen todos los demás
postulados, sea cual sea el valor tomado para \(1 + 1\).



\subsection{Independencia de el producto está siempre definido.}



	Modelo en el que solo falla \ref{axm:H1p}, esto es, que la operación
producto no es operación interna: no es función. Este es el caso en que
\( 0 \cdot 0 \notin \setB \). Ponemos en ese caso \(0 \cdot 0 = x\) pero
igualmente podríamos haber dejado en blanco ese lugar.



\[
    \begin{matrix}
        {\cdot + \cdot} & \big{\vert} & {0} & {1} & \big{\vert} \\
        \hline
        {0}             & \big{\vert} & {0} & {1} & \big{\vert} \\
        {1}             & \big{\vert} & {1} & {1} & \big{\vert} \\
        \hline
    \end{matrix}
    \quad
    \begin{matrix}
        {\cdot * \cdot} & \big{\vert} & {0} & {1} & \big{\vert} \\
        \hline
        {0}                   & \big{\vert} & {x} & {0} & \big{\vert} \\
        {1}                   & \big{\vert} & {0} & {1} & \big{\vert} \\
        \hline
    \end{matrix}
\]
%% (H1p) 
%% 0 + 0 = 0; 0 + 1 = 1; 1 + 0 = 1; 1 + 1 = 1; 
%% 0 * 0 = x; 0 * 1 = 0; 1 * 0 = 0; 1 * 1 = 1; 



	Como es fácil ver, existe elemento neutro tanto de la suma como del
producto. Ambas operaciones son conmutativas y se dan las dos propiedades de
distribución, siempre que tenga sentido el producto, esto es, no aparezca \(0 \cdot 0\). El complementario del \(0\) es \(\overline 0 = 1\) y el del
\(1\) es \(\overline 1 = 0\).



\subsection{Independencia de la unicidad de la definición del producto.}



Modelo en el que solo falla \ref{axm:H1s}, esto es, en el sentido que \(0
\cdot 0 \mapsto 0\) y también \(0 \cdot 0 \mapsto 1\).



\[
    \begin{matrix}
        {\cdot + \cdot} & \big{\vert} & {0} & {1} & \big{\vert} \\
        \hline
        {0}             & \big{\vert} & {0} & {1} & \big{\vert} \\
        {1}             & \big{\vert} & {1} & {1} & \big{\vert} \\
        \hline
    \end{matrix}
    \quad
    \begin{matrix}
        {\cdot * \cdot} & \big{\vert} & {0} & {1} & \big{\vert} \\
        \hline
        {0}                   & \big{\vert} & {x=0 \vee x=1} & {0} & \big{\vert} \\
        {1}                   & \big{\vert} & {0} & {1} & \big{\vert} \\
        \hline
    \end{matrix}
\]
%% (H2s) 
%% 0 + 0 = 0  ; 0 + 1 = 1; 1 + 0 = 1; 1 + 1 = 1; 
%% 0 * 0 = 0,1; 0 * 1 = 0; 1 * 0 = 0; 1 * 1 = 1;



\subsection{Independencia de la existencia de elemento neutro de la suma.}



Modelo en el que solo falla \ref{axm:H2s}, esto es, la existencia de
elemento neutro en la operación binaria interna suma.



\[
    \begin{matrix}
        {\cdot + \cdot} & \big{\vert} & {0} & {1} & \big{\vert} \\
        \hline
        {0}             & \big{\vert} & {0} & {0} & \big{\vert} \\
        {1}             & \big{\vert} & {0} & {0} & \big{\vert} \\
        \hline
    \end{matrix}
    \quad
    \begin{matrix}
        {\cdot * \cdot} & \big{\vert} & {0} & {1} & \big{\vert} \\
        \hline
        {0}                   & \big{\vert} & {0} & {0} & \big{\vert} \\
        {1}                   & \big{\vert} & {0} & {1} & \big{\vert} \\
        \hline
    \end{matrix}
\]
%% II.a H2s 
%% 0 + 0 = 0; 0 + 1 = 0; 1 + 0 = 0; 1 + 1 = 0; 
%% 0 * 0 = 0; 0 * 1 = 0; 1 * 0 = 0; 1 * 1 = 1;


	Primero, está claro que las operaciones internas están bien construidas
en cuanto a que son funciones.


	Queda claro que no hay elemento neutro de la suma pues \(\forall xy \in
\setB \quad x + y = 0\).


	Segundo, existe el elemento neutro del producto: \(0 \cdot 1 = 1 \cdot 0
= 0 \) y \( 1 \cdot 1 = 1 \). Luego se cumple \ref{aciom:H2p}.



	La conmutatividad se hace patente al ver las diagonales inversas de las
tablas de operación, que muestran un único valor. \( 0 + 1 = 1 + 0 = 0\) y
\(0 \cdot 1 = 1 \cdot 0 = 0\). Se cumplen \ref{axm:H3s} y \ref{axm:H3p}. 



	En cuanto a la distribución del producto sobre la suma, veamos si podemos
comprobarla de forma sencilla: \( a \cdot ( b + c ) = (a \cdot b) + (a \cdot
c) \). Sabemos que \(b + c = 0\) siempre, y que, pongamos que \(a \cdot b = x
\in \setB\) y que \(a \cdot c = y \in \setB\). Ahora bien \(x + y = 0\). Así
que todo lo que tenemos que probar es que \(a \cdot 0 = 0\), pero esto es
claro en la table del producto. Luego se cumple \ref{axm:H4s}.



	Ahora la distribución de la suma sobre el producto. \( a + ( b \cdot c )
= (a + b) \cdot (a + c) \). Sabemos que \(a + (b \cdot c) = 0\) siempre, y
que, pongamos que \( a + b = 0 \) y que \(a + c = 0 \). Ahora bien \(0 \cdot
0 = 0\). Luego se cumple \ref{axm:H4p}.


	Nos queda encontrar un complemento para el \(0\). Pero encontrar el
complemento solo tiene sentido si existen los dos elementos neutros, pero en
este caso no existe el neutro de la suma: no tiene sentido buscar el
complementario.



\subsection{Independencia de la existencia de elemento neutro del producto.}



	Exponemos un modelo en el que solo falla \ref{axm:H2p}, esto es, la existencia de elemento neutro en la operación binaria interna producto.



\[
    \begin{matrix}
        {\cdot + \cdot} & \big{\vert} & {0} & {1} & \big{\vert} \\
        \hline
        {0}             & \big{\vert} & {0} & {1} & \big{\vert} \\
        {1}             & \big{\vert} & {1} & {1} & \big{\vert} \\
        \hline
    \end{matrix}
    \quad
    \begin{matrix}
        {\cdot * \cdot} & \big{\vert} & {0} & {1} & \big{\vert} \\
        \hline
        {0}                   & \big{\vert} & {1} & {1} & \big{\vert} \\
        {1}                   & \big{\vert} & {1} & {1} & \big{\vert} \\
        \hline
    \end{matrix}
\]
%% (H2p) 
%% 0 + 0 = 0; 0 + 1 = 1; 1 + 0 = 1; 1 + 1 = 1; 
%% 0 * 0 = 1; 0 * 1 = 1; 1 * 0 = 1; 1 * 1 = 1;


	Este modelo es completamente simétrico al anterior, por lo que no nos
hará falta demostrar que se dan el resto de postulados. Solo hay que seguir
el esquema de la sub-sección anterior.



\subsection{Independencia de la conmutatividad de la suma.}



	Modelo en que la conmutatividad de la suma \ref{axm:H3s} no se dá, pero
si que se dan el resto de postulados.



\begin{equation}
    \begin{matrix}
        {\cdot + \cdot} & \big{\vert} & {0} & {1} & \big{\vert} \\
        \hline
        {0}             & \big{\vert} & {0} & {0} & \big{\vert} \\
        {1}             & \big{\vert} & {1} & {1} & \big{\vert} \\
        \hline
    \end{matrix}
    \quad
    \begin{matrix}
        {\cdot * \cdot} & \big{\vert} & {0} & {1} & \big{\vert} \\
        \hline
        {0}                   & \big{\vert} & {0} & {0} & \big{\vert} \\
        {1}                   & \big{\vert} & {0} & {1} & \big{\vert} \\
        \hline
    \end{matrix}
\end{equation}
%% (H3s) 
%% 0 + 0 = 0; 0 + 1 = 0; 1 + 0 = 1; 1 + 1 = 1; 
%% 0 * 0 = 0; 0 * 1 = 0; 1 * 0 = 0; 1 * 1 = 1;



	Que la suma y el producto están bien definidos como función es inmediato.



	Podemos comprobar que tanto \(\forall x \; x + 0 = x\) como que \(\forall
x \; x + 1 = x\). Esto quiere decir que la suma tiene dos neutros, el \(0\) y
el \(1\). Esto tendrá repercusiones sobre el elemento complementario. Podemos
ver por otra parte que \(1 + 0 \neq 0 + 1\), por lo que efectivamente la suma
no es conmutativa. Para el producto todo es idéntico, tiene elemento neutro
general y único \(0\) y \(1\cdot 0 = 0\cdot 1\) por lo que el producto es
conmutativo.



	Por otro lado, vemos que el complementario de \(0\) es \(1\), \(0 + 1 =
1\) -es neutro del producto- y \(0 * 1 = 0\) -es neutro de la suma-. Para ver
que \(1\) tiene complementario, hay que ir con más cuidado, \(1 + 1 = 1\) -
\(1\) es neutro del producto - y \(1 * 1 = 1\) - \(1\) es también neutro del
producto -. Luego existe complementario de \(1\) y es \(1\).



	Vamos ahora con las distribuciones:



	Nos preguntamos si se cumple \( x \cdot ( y + z ) = ( x \cdot y ) + ( x
\cdot z ) \) en el actual modelo. Si \(x = 0\) entonces \( 0 \cdot (x + y) =
0 \) y la parte derecha de la igualdad se resuelve en \( ( 0 \cdot y ) + ( 0
\cdot z ) \) pero como \( 0 \cdot x = 0 \) obtenemos que la distribución
\ref{axm:H4s} será verdad para \(x = 0\) si \( 0 = 0 + 0 \), cosa que es
cierta. Si \( x = 1 \), la igualdad a verificar quedaría \( 1 \cdot ( y + z )
= ( 1 \cdot y ) + ( 1 \cdot z ) \) y por \ref{axm:H2p} queda \( y + z  =  y
+ z \) que no es más que la identidad lógica de la igualdad. Se satisface
\ref{axm:H4s}.



	Nos preguntamos por la satisfacción de \ref{axm:H4p} \( x + ( y \cdot z
) = ( x + y ) \cdot ( x + z ) \) en el actual modelo. Procedemos como en el
párrafo anterior, por casos. Si \(x = 0\) entonces \( x + ( y \cdot z ) = ( x
+ y ) \cdot ( x + z ) \) \(\Longrightarrow\) \(0 + ( y \cdot z ) = ( 0 + y )
\cdot ( 0 + z )\) \(\Longrightarrow\) \( 0 = 0 \cdot 0\) lo que es cierto.
Para el caso \(x = 1\),  obtenemos \( x + ( y \cdot z ) = ( x + y ) \cdot ( x
+ z ) \) \(\Longrightarrow\) \( 1 + ( y \cdot z ) = ( 1 + y ) \cdot ( 1 + z
)\) \(\Longrightarrow\) \( 1 = 1 \cdot 1 \). Por lo tanto también se verifica
\ref{axm:H4p}.



	Queda comprobado que la conmutatividad de la suma \ref{axm:H3s} es independiente del resto de postulados.



\subsection{Independencia de la conmutatividad del producto.}



	Modelo en que la conmutatividad del producto \ref{axm:H3p} no se da,
pero, si se satisfacen el resto de postulados.



\begin{equation}
    \begin{matrix}
        {\cdot + \cdot} & \big{\vert} & {0} & {1} & \big{\vert} \\
        \hline
        {0}             & \big{\vert} & {0} & {1} & \big{\vert} \\
        {1}             & \big{\vert} & {1} & {1} & \big{\vert} \\
        \hline
    \end{matrix}
    \quad
    \begin{matrix}
        {\cdot * \cdot} & \big{\vert} & {0} & {1} & \big{\vert} \\
        \hline
        {0}                   & \big{\vert} & {0} & {0} & \big{\vert} \\
        {1}                   & \big{\vert} & {1} & {1} & \big{\vert} \\
        \hline
    \end{matrix}
\end{equation}
%% H3p 
%% 0 + 0 = 0; 0 + 1 = 1; 1 + 0 = 1; 1 + 1 = 1; 
%% 0 * 0 = 0; 0 * 1 = 0; 1 * 0 = 1; 1 * 1 = 1;



	El modelo es totalmente simétrico al anterior por lo que se puede
comprobar todo siguiendo los mismos pasos que en la sub-sección anterior.



\subsection{Independencia de la distribución de la suma sobre el producto.}



	El postulado del título no se cumple. Existen valores del modelo \(
\exists x \in \setB \) que no cumplen \ref{axm:H4p}, \( x + ( y \cdot z )
\neq ( x + y ) \cdot (x + z) \).



\begin{equation}
    \begin{matrix}
        {\cdot + \cdot} & \big{\vert} & {0} & {1} & \big{\vert} \\
        \hline
        {0}             & \big{\vert} & {0} & {1} & \big{\vert} \\
        {1}             & \big{\vert} & {1} & {0} & \big{\vert} \\
        \hline
    \end{matrix}
    \quad
    \begin{matrix}
        {\cdot * \cdot} & \big{\vert} & {0} & {1} & \big{\vert} \\
        \hline
        {0}                   & \big{\vert} & {0} & {0} & \big{\vert} \\
        {1}                   & \big{\vert} & {0} & {1} & \big{\vert} \\
        \hline
    \end{matrix}
\end{equation}
%% (H4s) 
%% 0 + 0 = 0; 0 + 1 = 1; 1 + 0 = 1; 1 + 1 = 0;
%% 0 * 0 = 0; 0 * 1 = 0; 1 * 0 = 0; 1 * 1 = 1;



	Este modelo es virtualmente idéntico al cuerpo sobre \(\setZ\) de restos
módulo 2, \(\setZ/{\mod2}\). Ese es el cambio que se da en la suma. Ahora \(1
+ 1 = 0\).



	Se dan por lo tanto todos los teoremas conocidos incluida la distribución
\ref{axm:H4s} e incluso sabemos que las operaciones son asociativas.



	Pero en aritmética no se da la dualidad que buscamos y la distribución
del producto sobre la suma no es una verdad.



	Solo nos queda ver si cumple la existencia de elemento complementario que
se encuentra visualmente por inspección de las tablas: \(\overline{1} = 0\) y
\(\overline{0} = 1\).


	Nos queda verificar que efectivamente no se cumple la distribución de la suma sobre el producto (que intuitivamente vemos). Encontramos un caso, \( 1 + ( y \cdot z ) \neq ( 1 + y ) \cdot ( 1 + z ) \). Lo vemos a continuación.


\begin{flalign}
  x = 1 &{} \\
  1 + ( y \cdot z ) &= ( 1 + y ) \cdot ( 1 + z ) \\
  \overline{y \cdot z} &= \overline{y} \cdot \overline{z} \\
  y = 0 &{} \\
  \overline{0 \cdot z} &= \overline{0} \cdot \overline{z} \\
  \overline{0} &= 1 \cdot \overline{z} \\
  1 &= \overline{z} \\
  z &= 0\\
  x=1 \wedge y=0 \wedge z=1 &\Longrightarrow \neg\ref{axm:H4p}
\end{flalign}



	De dónde efectivamente este modelo no distribuye la suma sobre un
producto. Y la independencia de \ref{axm:H4p} queda probada.



\subsection{Independencia de la distribución del producto sobre la suma.}


\begin{equation}
    \begin{matrix}
        {\cdot + \cdot} & \big{\vert} & {0} & {1} & \big{\vert} \\
        \hline
        {0}             & \big{\vert} & {0} & {1} & \big{\vert} \\
        {1}             & \big{\vert} & {1} & {1} & \big{\vert} \\
        \hline
    \end{matrix}
    \quad
    \begin{matrix}
        {\cdot * \cdot} & \big{\vert} & {0} & {1} & \big{\vert} \\
        \hline
        {0}                   & \big{\vert} & {1} & {0} & \big{\vert} \\
        {1}                   & \big{\vert} & {0} & {1} & \big{\vert} \\
        \hline
    \end{matrix}
\end{equation}
%% (H4s) 
%% 0 + 0 = 0; 0 + 1 = 0; 1 + 0 = 0; 1 + 1 = 1; 
%% 0 * 0 = 1; 0 * 1 = 0; 1 * 0 = 0; 1 * 1 = 1;



	Completamente simétrico al caso anterior.



\subsection{Independencia de la existencia del elemento complementario.}



	Por último presentamos un modelo que no cumple la existencia del
complementario y si satisface el resto de postulados.



\begin{equation}
    \begin{matrix}
        {\cdot + \cdot} & \big{\vert} & {0} & {1} & \big{\vert} \\
        \hline
        {0}             & \big{\vert} & {0} & {1} & \big{\vert} \\
        {1}             & \big{\vert} & {1} & {1} & \big{\vert} \\
        \hline
    \end{matrix}
    \quad
    \begin{matrix}
        {\cdot * \cdot} & \big{\vert} & {0} & {1} & \big{\vert} \\
        \hline
        {0}                   & \big{\vert} & {0} & {1} & \big{\vert} \\
        {1}                   & \big{\vert} & {1} & {1} & \big{\vert} \\
        \hline
    \end{matrix}
\end{equation}
%% (H5) 
%% 0 + 0 = 0; 0 + 1 = 1; 1 + 0 = 1; 1 + 1 = 1;
%% 0 * 0 = 0; 0 * 1 = 1; 1 * 0 = 1; 1 * 1 = 1;



	Como podemos ver en este modelo la suma y el producto son idénticos, por
lo que el \(0\) es el único neutro, tanto de la suma como del producto. La
conmutatividad es evidente a primera vista. 


	Veamos la distribución del producto respecto de la suma. Com la suma y el
producto son idénticos, bastará con probar este caso únicamente.



\begin{equation}
x \cdot (y + z) = x + ( y + z ) = ( x + x ) + ( y + z ) = ( x + y ) + ( x + z
) = ( x + y ) \cdot ( x + z )
\end{equation}



	Solo queda probar que \(x = x + x\) (que se hace por simple inspección
visual de las tablas) y la propiedad asociativa para este modelo, que
tendremos que probar. Lo haremos mediante una tabla.



\begin{equation}
  \begin{matrix}
  \hline
  \mathsf{a} & \mathsf{b} & \mathsf{c} & \mathsf{b}\cdot\mathsf{c} &
  \mathsf{a}\cdot\left( \mathsf{b} \cdot \mathsf{c} \right) &
  \left(\mathsf{a}\cdot \mathsf{b}\right) \cdot \mathsf{c} &
  \mathsf{a}\cdot\mathsf{b}\\
  \hline
%%a   b   c   bc a(bc)(ab)c ab
  0 & 0 & 0 & 0 & 0    & 0 & 0 \\
  0 & 0 & 1 & 1 & 1    & 1 & 0 \\
  0 & 1 & 0 & 1 & 1    & 1 & 1 \\
  0 & 1 & 1 & 1 & 1    & 1 & 1 \\
  1 & 0 & 0 & 0 & 1    & 1 & 1 \\
  1 & 0 & 1 & 1 & 1    & 1 & 1 \\
  1 & 1 & 0 & 1 & 1    & 1 & 1 \\
  1 & 1 & 1 & 1 & 1    & 1 & 1 \\
  \hline
  1 & 1 & 1 & 2 & 3    & 3 & 2 \\
  \hline
  \end{matrix}
\end{equation}


\section{Algunos ejemplos}



	El interés de esta sección está en dar ejemplos variados de álgebras de
Boole, sin estudiar ya la independencia de los postulados entre sí, cosa que
quedó clara ya en la sección anterior. Además muestran la consistencia de los
axiomas, puesto que al encontrar modelos (ejemplos) dónde todos los axiomas
se cumplen, de camino completamos el punto anterior viendo que los axiomas
son consistentes: no son contradictorios entre sí.



\subsection{Ejemplo: el álgebra de las partes de un conjunto.}



	El ejemplo primero y más sencillo de ver es el álgebra de las partes de
un conjunto $\mathsf{U}$. Dado que para un conjunto cualquiera,
\large {$\mathscr{P}$} \normalsize $\mathsf{U}$ $\mathdef$ $\setCompr{
\mathsf{ X } } { \mathsf{X} \subseteq \mathsf{ U } }$ es también un conjunto
en cualquiera de los dos sistemas axiomáticos de teoría de conjuntos
mencionados.



	Además, dónde se verifica que $\emptyset \in$ \large $\mathscr{P}$
\normalsize $\mathsf{U}$ y que  $\mathsf{U} \in $ \large $\mathscr{P}$
\normalsize $\mathsf{U}$.



	Definiremos, $\setB \mathdef $ \large{$\mathscr{P}$} \normalsize
$\mathsf{U}$, \normalsize $0_{\setB} \mathdef \emptyset$ y $1_{\setB}
\mathdef \mathsf{U}$. Como producto lógico pondremos la intersección de
conjuntos $* \mathdef \bigcap$ y como suma lógica pondremos la unión de
conjuntos $+ \mathdef \bigcup$.



	Las conmutatividades de la suma y el producto se intercambian
directamente por la conmutatividad de la unión y la intersección.



	El elemento neutro de la suma es trivialmente $\emptyset$ ya que es
neutro para la unión. Igualmente el neutro del producto es $\mathsf{U}$ por
ser el neutro de la intersección en \large $\mathscr{P}$ \normalsize
$\mathsf{U}$.



	Las distributividades recaen directamente en las de la unión sobre la
intersección y de la intersección sobre la unión.



	Sea $\forall \mathsf{X} \in \;$ \large
$\mathscr{P}$\normalsize$\mathsf{U}$$\quad \exists \mathsf{!}
{\mathsf{Y}}_{\mathsf{X}}$,  ${\mathsf{Y}}_{\mathsf{X}} \mathdef
\left(\mathsf{U} \setminus \mathsf{X}\right)$ y a partir de ahí sabemos que
puesto que $\mathsf{X} \cup \left(\mathsf{U}\setminus\mathsf{X}\right) =
\mathsf{U} = 1_{\setB}$ y $\mathsf{X} \cap \left( \mathsf{U} \setminus
\mathsf{X} \right) = \emptyset = 0_{\setB}$, que quiere decir que existe un
subconjunto complementario para cada subconjunto, y éste es único. 
Normalmente denotaremos ${ {\complement}^{\mathsf{U} } }$ \Large $\mathsf{x}$
\normalsize $ \mathdef \left({\mathsf{U}\setminus\mathsf{X}}\right) \mathdef
{\overline{\mathsf{U}}}$.



\subsection{Ejemplo: álgebra de productos de números primos.}



	Un ejemplo interesante, fácil de construir, es el álgebra de Boole de los
números que son producto de los primeros números primos (cantidad finita de
ellos) y sus divisores. Para algún $n \in \setN$ consideramos el conjunto



\begin{equation}
\mathsf{P}_{n} \mathdef \set{p_1 = 2, p_2 = 3, p_3 = 5,p_4 = 7, p_5 = 11,
p_6 = 13,\ldots,p_n} \subsetneq \setP
\end{equation}



, definimos el conjunto booleano cómo



\begin{equation}
\setB \mathdef {
\setCompr
{m \in \setN\;} 
{\;m \mid {\prod}_{\kappa=1}^{n} p_{\kappa},\;p_{\kappa} \in \mathsf{P}_{n}}
}
\end{equation}



	Consideramos $1_{\setB} \mathdef {\prod}_{q \in \mathsf{P}_{n}} q$, 
$0_{\setB} \mathdef 1\in\setN$ y las operaciones serán el
mínimo común múltiplo como suma booleana \(a + b \mathdef
\lcm\left(a,b\right)\) y el máximo común divisor a como el producto
lógico $a* b\mathdef\gcd\left(a,b\right)$. En cuanto al elemento
complementario lo definimos como 



\begin{equation}
\forall k \in \setB 
\quad 
{\overline k} 
\mathdef 
\frac{1_{\setB}}{k} = 
{\prod \setCompr{p \in \mathsf{P}_n}{p \nmid k } }
\end{equation}



	Éste es un modelo fácil de desarrollar para poner ejemplos.



\subsection{Ejemplo: álgebras definidas a partir de otra y un elemento.}



	Partimos de un álgebra de Boole cualquiera \(\setB\) que tenga más de dos
elementos, y de un elemento \(x \in \setB\setminus\set{0,1} \). Definimos
ahora dos conjuntos, el primero como \({\setB}_{\leq x} \mathdef \setCompr{x
* y}{y \in \setB}\) y el segundo como \({\setB}_{\geq x} \mathdef
\setCompr{x + y}{y \in \setB}\). Las álgebras de Boole nuevas a considerar
son \({\setB}_{\leq x}\) y \({\setB}_{\geq x}\). Consideramos el mismo
producto booleano que en el conjunto inicial e idénticamente hacemos con la
suma booleana (solo que restringidos al subconjunto elegido). Ahora nos
queda, para el álgebra de los mayores que \(x\), que el cero lógico en
\(0_{{\setB}_{\geq x}} \mathdef x\) y el uno lógico queda como
\(1_{{\setB}_{\geq x}} \mathdef 1_{\setB}\), y para el álgebra de los menores
que \(x\), \(0_{{\setB}_{\leq x}} \mathdef 0_{\setB}\) y \(1_{{\setB}_{\geq
x}} \mathdef x\). El complemento varía, en el conjunto de los mayores que
\(x\), de la siguiente forma, \(y \in {\setB}_{\geq x} \; y' \mathdef
\overline{y} + x\) dónde el complemento de \(y\) es \(y^\prime\). En el
álgebra de los menores que \(x\), el complemento es \(y' \mathdef
\overline{y} * x\).



\subsection{Ejemplo: álgebras definidas a partir de otra y dos elementos.}



	Partimos de un álgebra de Boole cualquiera \(\algB\), cuyo conjunto base
tenga más de cuatro elementos. Sea \(a\in\setB\) y \(b\in\setB\), de forma
que \(a * b \neq 0_{\setB}\) y \(a + b \neq 1_{\setB}\). Ahora consideraremos
\(0_{\setB_{ a \leq \cdot \leq b} } \mathdef a * b \) y \( 1_{ \setB_{a
\leq \cdot \leq b}} \mathdef a + b\). Podemos considerar el conjunto 



\begin{equation}
{\setB}_{a\leq \cdot\leq b} 
\mathdef 
	\setCompr{ 
		1_{\setB_{a\leq\cdot\leq b}} * z
		}
		{z \in \setB} 
	\cap 
	\setCompr{
		0_{\setB_{a\leq\cdot\leq b}} + z
		}
		{z \in \setB}
=
	\setCompr{
		\left({1_{\setB_{a\leq\cdot\leq b}}}* x\right)
		+{0_{\setB_{a\leq\cdot\leq b}}}
	}
	{
		x \in \setB
	}
\end{equation}



Dados 
\(
x \in \setB_{a \leq \cdot \leq b}
\Rightarrow 
\exists x^{\prime} \in \setB
\quad 
x = 1_{\setB_{a \leq \cdot \leq b}} * x^{\prime} + 0_{\setB_{a \leq
\cdot \leq b}} 
\), e idénticamente con
\(
y \in \setB_{a \leq \cdot \leq b} 
\Rightarrow 
\exists y^{\prime} 
\quad 
y = 1_{\setB_{a \leq \cdot \leq b}} * y^{\prime} + 0_{\setB_{a \leq
\cdot
\leq b}} \) las operaciones quedan:

\begin{equation}
  x {*}_{a \leq\cdot\leq b} y \mathdef 
  1_{a \leq\cdot\leq b} * 
  \left(
  x^\prime * y^\prime + 
  \left(
  0_{a \leq\cdot\leq b} * \left(x^\prime + y^\prime \right)
  \right)
  \right) + 0_{a \leq\cdot\leq b}
\end{equation}



\begin{equation}
  x {+}_{a \leq\cdot\leq b} y \mathdef 
  1_{a \leq\cdot\leq b} * 
  \left(
  x^\prime + y^\prime 
  \right) 
  + 
  0_{a \leq\cdot\leq b}
\end{equation}



Cogiendo la definición anterior de \(x \in \setB_{a\leq\cdot\leq b}\) que
conlleva la existencia de \(x^\prime \in \setB\), tal que \(x =
1_{a\leq\cdot\leq b} * x^\prime + 0_{a\leq\cdot\leq b}\), entonces la
inversa en \(\setB_{a\leq\cdot\leq b}\) de \(x\) será \(\tilde x \mathdef
1_{a\leq\cdot\leq b} * \overline{x^\prime} + 0_{a\leq\cdot\leq b}\).



\subsection{Ejemplo: el álgebra de las proposiciones.} 



	Éste es sin duda, el primer desarrollo que se hizo del álgebra de Boole,
hecha por el propio George Boole en mitad del siglo XIX (edición 1851). Lo
desarrolló como una “Una investigación en las leyes del pensamiento”. Amigo
suyo que lo ayudó a penetrar los ambientes académicos es Augustus De Morgan.
Desde Aristóteles (que funda por primera vez la lógica como ciencia
analítica) en el siglo IV a. C. no había habido ningún adelanto sustancial en
la lógica. Kant (medio siglo antes de Boole) había considerado que la lógica
era un cuerpo de doctrina cerrado y completo (esto es, no había nada más que
decir que lo que ya había desarrollado y escrito Aristóteles en sus
“Tratados de Lógica” u “Organon”, hacía ya 2.100 años). Aunque la verdadera
revolución se da algunos años más tarde con Frege, el lógico más importante
desde Aristóteles. Si doy estos datos sobre la historia de la lógica que
todos asociaréis más a la filosofía, que parece queda muy lejos del propósito
de unos apuntes de matemáticas discretas que cubran de la forma más amplia
posible los Fundamentos de Electrónica en su parte Digital, es porque en
realidad no queda tan lejos. La idea de hacer un lenguaje dónde el
razonamiento siguiera unas pautas claras de forma que siempre quedara todo
tan cierto como en las matemáticas era ya antigua. Aristóteles ya advertía de
una cierta in-definición insuperable de los términos más importantes de la
filosofía (en realidad de casi todos los conceptos de la vida ordinaria):
“existen conceptos o ideas que corresponden con la realidad que no se usan de
forma equívoca – esto es, su uso no es equívoco, este mismo concepto, palabra
o idea que hablamos no se refiere a realidades distintas y diferenciadas, de
forma que nos llevan a confusión – pero tampoco de forma unívoca – como las
definiciones desde axiomas en un lenguaje formal matemático, así tenemos que
existen conceptos análogos” (es una glosa de palabras de Aristóteles). En la
Modernidad, dado que el concepto de analogía lleva aparejado un tratamiento
difícil que no lleva fácilmente a certeza, se intentan buscar criterios de
certeza absoluta y unas definiciones que aparentemente son unívocas y se
tratan como tales. Es significativo el nombre (y la estructura interna) de
una importante obra de Spinoza: “Ética demostrada según el orden
geométrico”. Será Leibnitz quién escriba ya cumplidamente sobre la necesidad
de establecer un léxico completamente unívoco (una tarea mastodonte, o mejor,
imposible) y un “cálculo” del pensamiento, de forma que “una cuestión como la
existencia de Dios pueda ser resuelta mediante la resolución de unas
ecuaciones de pensamiento” (de nuevo es una glosa). Se empezaba a buscar con
ansiedad una mecanización del pensamiento. Este motivo fue muy fructífero,
dando un primer paso hacia él, George Boole que hace un álgebra de las
proposiciones. Esta álgebra no era más amplia que la de Aristóteles (de
hecho, era solo un subconjunto de la lógica aristotélica), pero permitía el
cálculo al modo algebraico. De aquí a la llegada de Frege, ya, Charles
Babbage diseña y realiza (sin éxito debido al trabajo de mecanizado
excesivamente minucioso que requería el diseño) una computadora universal
mecánica (mediante engranajes) prácticamente similar al modelo de Von
Neumann. La condesa de Lovelace (Ada) es la matemática que hace los primeros
programas en lenguaje ensamblador de la máquina de Babbage. La primera
programadora de la historia. Después viene Frege, al que debemos el invento
de la Lógica-Matemática, o lógica extensional en aquel nacer. Ampliaba
inmensamente la lógica aristotélica pues la estructura de la proposición no
seguía ya el esquema sujeto-predicado que se mostraba insuficiente por lo que
adopta el esquema función-argumentos, además de introducir los
cuantificadores (“existe al menos un” \(\exists\) y “para todo” \(\forall\)).
Después de Frege siguen los desarrollos con gente como Bertrand Russell,
David Hilbert y otros, hasta los increíbles resultados de Gödel (y
posteriormente Turing) que ponen punto final a muchas de las pretensiones de
mecanización del pensamiento, pero que son ya la base de la computación
moderna, siendo los trabajos definitivos los de Alan Turing. Como veis el
camino recorrido es largo y complicado, siendo el momento crucial para el
arranque de la ingeniería digital los trabajos de George Boole. No he
mencionado intentos de mecanización más antiguos como los del mallorquín
Raimundo Lulio, ni la calculadora binaria de Leibnitz, ni mecanismos remotos
como el mecanismo de Anticitera, así como el papel fundamental que jugó la
crisis de las matemáticas de finales del siglo XIX y principios del XX ni el
papel de las máquinas de cifrado y descifrado de mensajes en la Gran Guerra y
la II Guerra Mundial (en las que intervinieron muchas de las mentes antes
mencionadas).



	De manera un tanto informal podemos ver una proposición como una frase
que afirma o niega una propiedad de un objeto, una relación entre objetos o
la existencia del mismo, una frase que ha de ser, o verdadera \(\top \equiv
\mathsf{verdadero} \equiv \mathsf{true} \equiv 1\), o falsa \(\bot \equiv
\mathsf{falso} \equiv \mathsf{false} \equiv 0\), sin lugar a la 
interpretación. Este conjunto de proposiciones pueden ser operadas mediante 
la conjunción ‘y’ y los símbolos para la operación (normalmente infija): \(
\cdot\wedge\cdot \equiv \cdot*\cdot \equiv \cdot\&\!\&\cdot 
\equiv \cdot\mathsf{and}\cdot \equiv \cdot\mathsf{y}\cdot \), es verdadero si
es verdadero y es verdadero a la vez y falso en cualquier otro caso. La
disyunción sería la ‘o’ y sus símbolos (infijos): \(\cdot \vee \cdot  \equiv 
\cdot + \cdot \equiv \cdot |\!| \cdot \equiv \cdot \mathsf{or} \cdot \), 
siendo verdadera con que cualquiera de las proposiciones sea verdadera, y 
siendo falsa sólo cuando todas las proposiciones son falsas y lo son a la 
vez.  Para el ‘no’ usamos los siguientes símbolos: \(\neg\cdot \equiv !\cdot 
\equiv \overline{\;\cdot\;}\equiv \mathsf{not} \). En cuanto a la notación 
más habitual, en electrónica digital es \(\cdot * \cdot\), \(\cdot+\cdot\) y 
\(\overline{\;\cdot\;}\) para los funtores, y para las constantes \(1\equiv 
\mathsf{H}\) y \(0 \equiv \mathsf{L}\), en computación, por ejemplo en 
lenguajes derivados de la familia de \textsf{C}, \(\cdot\&\!\&\cdot\), 
\(\cdot |\!|\cdot\) y \(!\cdot\) para el nivel de \emph{tipo booleano}, y  
\(\cdot\&\cdot\), \(\cdot |\cdot\) y \(\sim\cdot\) para el nivel de
\emph{tipo bit} y operaciones \emph{bitwise}, y en cuanto a las constantes
hay más variedad, por ejemplo en \(\textsf{C}\!+\!+\), \(\mathsf{true}\) y
\(\mathsf{false}\) a nivel del \emph{tipo booleano}, en \textsf{C}, el valor
falso recae en cualquier entero igual a \(0\) y el verdadero en ser distinto
de \(0\), en \textsf{Python} o \textsf{Haskell}, los valores son
\(\mathsf{True}\) y \(\mathsf{False}\). En matemáticas (teoría de retículos)
y lógica, el estándar es \(\cdot\wedge\cdot\), \(\cdot\vee\cdot\) y 
\(\neg\cdot\) para las operaciones, \(\top\) para el valor verdadero y 
\(\bot\) para el valor falso. El conjunto de Boole es el conjunto de las
proposiciones de la que partamos.



	De manera más formal se considera un elemento del álgebra de Boole de la
lógica a las clases de equivalencia de las proposiciones equivalentes
lógicamente (en su valor de verdad o falsedad) entre sí.



\subsection{Ejemplo: el álgebra de conmutación de Shannon.}



	Esta es el álgebra de Boole más sencilla que hay: \(\setB \mathdef
{\setB}_2 = \set{0,1}\). Las operaciones las concretaremos en tablas:



\[
  \begin{matrix}
   a + b & {} & \vline & b & {} \\
   {} & {} & \vline & 0 &  1 \\
    \hline
   {a} & {0} & \vline  &  0 & 1 \\
   {} & {1} & \vline  &  1 & 1 \\
  \end{matrix}
  \qquad
  \begin{matrix}
   a \cdot b & {} & \vline & b & {} \\
   {} & {} & \vline & 0 &  1 \\
    \hline
   {a} & {0} & \vline  &  0 & 0 \\
   {} & {1} & \vline  &  0 & 1 \\
  \end{matrix}
  \qquad
  \begin{matrix}
    {} & {} & \vline & {\overline a} \\
    {} & {} & \vline & {} \\
    \hline
    a & 0 & \vline &  1 \\
    {} & 1 & \vline &  0 \\
  \end{matrix}
\]



Podréis comprobar fácilmente que se cumplen todos los postulados de
Huntington tomando como álgebra la definida en las anteriores tablas. Éstas
serán usadas frecuentemente durante el curso. Esta álgebra está contenida en
todo álgebra de Boole. Cumple de la manera más escueta posible con todos los
postulados de Huntington.



Como diagrama de Hasse (dónde se hace patente un orden parcial en el álgebra
de Boole) tenemos el siguiente:

\begin{center}
\begin{tikzpicture}
    \node (top) at (0,0) {$ \mathfrak{1} $};
    \node [below of=top](bot)  {$ \mathfrak{0} $};
    \draw [blue,  thick, -{stealth}] (top) -- (bot);
\end{tikzpicture}
\end{center}


	Pero veamos algo más sobre este álgebra. Supongamos un conjunto \(
\mathsf{U}_1 = \set{ \bigstar } \). Pongamos el álgebra de conjuntos sobre el
conjunto potencia del anterior: \large\(\mathscr{P}\) \normalsize \(
\mathsf{U}_1 \;=\; \set{ \emptyset , \set{ \bigstar } } \).



	Ahora hagamos una biyección entre \( \setB_2 \) y \large \( \mathscr{P}\!
\) \normalsize \( \mathsf{U}_1 \). \(\mathfrak{f}\;:\; \setB_2 \fromto \)
\large \( \mathscr{P} \) \normalsize \( \mathsf{U}_1 \) tal que \(
\mathfrak{f} (0) = \emptyset \) y \( \mathfrak{f} (0) = \mathsf{U}_1 \).
Podremos observar las siguientes relaciones:



\newcommand {\fbat}[1]{\ensuremath{\mathfrak{f}\left({#1}\right)}}
\begin{flalign}
\fbat{0} &= \emptyset \\
\fbat{1} &= {\mathsf{U}}_1 \\
\fbat{x + y} &= \fbat{x} \cup \fbat{y} \\
\fbat{x \cdot y} &= \fbat{x} \cap \fbat{y} \\
\fbat{\overline{x}} &= {\mathsf{U}}_1 \setminus \fbat{x}
\end{flalign}



	Esto quiere decir que los dos sistemas son, en cuanto álgebras de Boole,
la misma cosa. Esto es un isomorfismo de álgebras de Boole. El álgebra de
Boole de conmutación o de Shannon es virtualmente idéntica al álgebra de los
conjuntos sobre el conjunto potencia de un conjunto de un solo elemento. A
decir verdad, todas las álgebras de Boole de \(2\) elementos son isomorfas
entre sí. Esto es algo que ampliaremos para muchos en otros casos.



\subsection{Ejemplo: el álgebra sobre conjuntos de cardinalidad 2.}



	Esta es el siguiente álgebra de Boole más sencilla: \(\setB \mathdef
{\setB}_4 = \set{\set{a,b},\set{a},\set{b},\set{}=\emptyset}\). Las
operaciones las concretaremos en tablas, proceden de la definición del
álgebra de las partes de un conjunto ya explicada anteriormente, aunque la
detallaremos aquí también.



	Llamamos \( \setB_4 \mathdef \) \large \( \mathscr{P} \) \normalsize \( 
\mathsf{U}_2 \), dónde \( \mathsf{U}_2 = \set{a,b} \). De esta forma nos 
queda que \large \( \mathscr{P} \) \normalsize \( \mathsf{U}_2 = \set{ 
\mathsf{U}_2, \set{a}, \set{b}, \emptyset} \). Nombraremos, por comodidad, \( 
\mathfrak{a} \mathdef \set{a} \) y \( \mathfrak{b} \mathdef \set{b} \).
Claramente \( \mathfrak{1} \mathdef 1_{\setB_4} \mathdef \mathsf{U}_2 \) y \(
\mathfrak{0} \mathdef 0_{\setB_4} \mathdef \emptyset \). De esta forma
nuestra notación para el conjunto de Boole y sus elementos será: \(\setB_4 =
\set{\mathfrak{0},\mathfrak{a},\mathfrak{b},\mathfrak{1}}\).



\begin{equation}
  \begin{matrix}
   a + b & {} & \vline & b & {} \\
   {} & {} & \vline & \mathfrak{0} & \mathfrak{a} & \mathfrak{b} & \mathfrak{1} \\
    \hline
   {a} & \mathfrak{0} & \vline  &  \mathfrak{0} & \mathfrak{a} & \mathfrak{b} & \mathfrak{1} \\
   {}  & \mathfrak{a} & \vline  &  \mathfrak{a} & \mathfrak{a} & \mathfrak{1} & \mathfrak{1} \\
   {}  & \mathfrak{b} & \vline  &  \mathfrak{b} & \mathfrak{1} & \mathfrak{b} & \mathfrak{1} \\
   {}  & \mathfrak{1} & \vline  &  \mathfrak{1} & \mathfrak{1} & \mathfrak{1} & \mathfrak{1} \\
  \end{matrix}
  \qquad
  \begin{matrix}
   a \cdot b & {}  & \vline & b & {} \\
   {} & {} &         \vline & \mathfrak{0} & \mathfrak{a} & \mathfrak{b} & \mathfrak{1} \\
   \hline
   {a} & \mathfrak{0} & \vline  &  \mathfrak{0} & \mathfrak{0} & \mathfrak{0} & \mathfrak{0} \\
   {}  & \mathfrak{a} & \vline  &  \mathfrak{0} & \mathfrak{a} & \mathfrak{0} & \mathfrak{a} \\
   {}  & \mathfrak{b} & \vline  &  \mathfrak{0} & \mathfrak{0} & \mathfrak{b} & \mathfrak{b} \\
   {}  & \mathfrak{1} & \vline  &  \mathfrak{0} & \mathfrak{a} & \mathfrak{b} & \mathfrak{1} \\
  \end{matrix}
  \qquad
  \begin{matrix}
    {}  & {} & \vline & {\overline a} \\
    {}  & {} & \vline & {} \\
    \hline
    {a} & \mathfrak{0} & \vline &  \mathfrak{1} \\
    {}  & \mathfrak{a} & \vline &  \mathfrak{b} \\
    {}  & \mathfrak{b} & \vline &  \mathfrak{a} \\
    {}  & \mathfrak{1} & \vline &  \mathfrak{0} \\
  \end{matrix}
\end{equation}



	Podréis comprobar fácilmente que se cumplen todos los postulados de
Huntington tomando como álgebra la definida en las anteriores tablas. Esta
álgebra está es la primera que tiene más elementos que el \(\set{0,1}\).



	Como diagrama de Hasse (dónde se hace patente un orden parcial en el
álgebra de Boole, y por primera vez es parcial) tenemos el siguiente:



\begin{center}
\begin{tikzpicture}
    \node (top) at (0,0) {$ \mathfrak 1 $};
    \node [below left of=top](left)  {$ \mathfrak a $};
    \node [below right of=top](right)  {$ \mathfrak b $};
    \node [below right of=left](bot)  {$ \mathfrak 0 $};
    \draw [black,  thick, -{stealth}] (top) -- (left);
    \draw [black,  thick, -{stealth}] (top) -- (right);
    \draw [red,  thick, -{stealth}] (left) -- (bot);
    \draw [red,  thick, -{stealth}] (right) -- (bot);
\end{tikzpicture}
\end{center}



\subsection{Ejemplo: el álgebra sobre conjuntos de cardinalidad 3.}



	Esta es el siguiente álgebra de Boole: 



\begin{equation}
\setB \mathdef {\setB}_{8} = \left\lbrace
\set{a,b,c},\set{a,b},\set{a,c},\set{b,c},
\set{a},\set{b},\set{c},\set{} \right\rbrace
\end{equation}



	Las operaciones, las mostramos en tablas y proceden de la definición del
álgebra de las partes de un conjunto.



	Llamamos \( \setB_{8} \mathdef \) \large\( \mathscr{P} \) \normalsize
\( \mathsf{U}_3 \), dónde \( \mathsf{U}_3 = \set{a,b,c} \). De esta forma
nos queda



\begin{flalign}
\setB_8 = \mbox{\large\({\mathscr{P}} \)} \mbox{\normalsize \(\mathsf{U}_3
=\)}&
\left\lbrace\mathsf{U}_3,\right. \\ \nonumber
\qquad &\set{a,b},\set{a,c},\set{b,c}, \\ \nonumber
\qquad &\set{a},\set{b},\set{c},\\ \nonumber
\qquad &\left.\emptyset\right\rbrace \\ \nonumber
\end{flalign}



	Nombraremos, por comodidad, \( \mathfrak{a} \mathdef \set{a} \), \(
\mathfrak{b} \mathdef \set{b} \)  y  \( \mathfrak{c} \mathdef \set{c} \).
También nombraremos, \( \mathfrak{A} \mathdef \set{a,b} \), \( \mathfrak{B}
\mathdef \set{a,c} \) y \( \mathfrak{C} \mathdef \set{b,c} \). Este álgebra
no tiene más capas. Claramente \(\mathfrak{1} \mathdef 1_{\setB_{8}} \mathdef
\mathscr{P}\mathsf{U}_3 \) y \(\mathfrak{0} \mathdef 0_{\setB_{8}} \mathdef
\emptyset \). De esta forma nuestra notación para el conjunto de Boole y sus
elementos será:




\begin{equation}
\setB_{8} =
\set{
\mathfrak{0},
\mathfrak{a},\mathfrak{b},\mathfrak{c},
\mathfrak{A},\mathfrak{B},\mathfrak{C},
\mathfrak{1}
}
\end{equation}



	La tabla para la suma, a continuación.



\begin{equation}
  \begin{array}{@{}*{22}c@{}}
      {\textsf{a+b}}&%%c1 operacion
      \big\vert     &%%l1 vertical line
      {}            &%%c2 linea de ¿?
      \big\vert     &%%l2 vertical line
      {}            &%% 1
      {}            &%% 2
      {}            &%% 3
      {}            &%% 4
      {\textsf{b}}  &%% 5 variable b
      {}            &%% 6
      {}            &%% 7
      {}            &%% 8
      \big\vert      %%l3
      \\ \hline \nonumber
      {}            &%%c1
      \big\vert     &%%l1
      {}            &%%c2
      \big\vert     &%%l2
      {\mathfrak{1}}&%% 1 {a,b,c}
      {\mathfrak{A}}&%% 2 {a,b}
      {\mathfrak{B}}&%% 3 {a,c}
      {\mathfrak{C}}&%% 4 {b,c}
      {\mathfrak{a}}&%% 5 {a}
      {\mathfrak{b}}&%% 6 {b}
      {\mathfrak{c}}&%% 7 {c}
      {\mathfrak{0}}&%% 8 {}
      \big\vert      %%l3
      \\ \hline \nonumber
      {}            &%%c1
      \big\vert     &%%l1
      {\mathfrak{1}}&%%c2 {a,b,c} lleno
      \big\vert     &%%l2
      {\mathfrak{1}}&%% 1 {a,b,c}
      {\mathfrak{1}}&%% 2 {a,b}
      {\mathfrak{1}}&%% 3 {a,c}
      {\mathfrak{1}}&%% 4 {b,c}
      {\mathfrak{1}}&%% 5 {a}
      {\mathfrak{1}}&%% 6 {b}
      {\mathfrak{1}}&%% 7 {c}
      {\mathfrak{1}}&%% 8 {}
      \big\vert      %%l3
      \\ \nonumber
      {}            &%%c1
      \big\vert     &%%l1
      {\mathfrak{A}}&%%c2 {a,b}
      \big\vert     &%%l2
      {\mathfrak{1}}&%% 1 {a,b,c}
      {\mathfrak{A}}&%% 2 {a,b}
      {\mathfrak{1}}&%% 3 {a,c}
      {\mathfrak{1}}&%% 4 {b,c}
      {\mathfrak{A}}&%% 5 {a}
      {\mathfrak{A}}&%% 6 {b}
      {\mathfrak{1}}&%% 7 {c}
      {\mathfrak{A}}&%% 8 {}
      \big\vert      %%l3
      \\ \nonumber
      {}            &%%c1
      \big\vert     &%%l1
      {\mathfrak{B}}&%%c2 {a,c}
      \big\vert     &%%l2
      {\mathfrak{1}}&%% 1 {a,b,c}
      {\mathfrak{1}}&%% 2 {a,b}
      {\mathfrak{B}}&%% 3 {a,c}
      {\mathfrak{1}}&%% 4 {b,c}
      {\mathfrak{B}}&%% 5 {a}
      {\mathfrak{1}}&%% 6 {b}
      {\mathfrak{B}}&%% 7 {c}
      {\mathfrak{B}}&%% 8 {}
      \big\vert      %%l3
      \\ \nonumber
      {}             &%%c1
      \big\vert      &%%l1
      {\mathfrak{C}} &%%c2 {b,c}
      \big\vert      &%%l2
      {\mathfrak{1}} &%% 1 {a,b,c}
      {\mathfrak{1}} &%% 2 {a,b}
      {\mathfrak{1}} &%% 3 {a,c}
      {\mathfrak{C}} &%% 4 {b,c}
      {\mathfrak{1}} &%% 5 {a}
      {\mathfrak{C}} &%% 6 {b}
      {\mathfrak{C}} &%% 7 {c}
      {\mathfrak{C}} &%% 8 {}
      \big\vert       %%l3
      \\ \nonumber
      {\textsf{a}}   &%%c1 variable a
      \big\vert      &%%l1
      {\mathfrak{a}} &%%c2 {a}
      \big\vert      &%%l2
      {\mathfrak{1}} &%% 1 {a,b,c}
      {\mathfrak{A}} &%% 2 {a,b}
      {\mathfrak{B}} &%% 3 {a,c}
      {\mathfrak{1}} &%% 4 {b,c}
      {\mathfrak{a}} &%% 5 {a}
      {\mathfrak{A}} &%% 6 {b}
      {\mathfrak{B}} &%% 7 {c}
      {\mathfrak{a}} &%% 8 {}
      \big\vert       %%l3
      \\ \nonumber
      {}             &%%c1
      \big\vert      &%%l1
      {\mathfrak{b}} &%%c2 {b}
      \big\vert      &%%l2
      {\mathfrak{1}} &%% 1 {a,b,c}
      {\mathfrak{A}} &%% 2 {a,b}
      {\mathfrak{1}} &%% 3 {a,c}
      {\mathfrak{C}} &%% 4 {b,c}
      {\mathfrak{A}} &%% 5 {a}
      {\mathfrak{b}} &%% 6 {b}
      {\mathfrak{C}} &%% 7 {c}
      {\mathfrak{b}} &%% 8 {}
      \big\vert       %%l3
      \\ \nonumber
      {}             &%%c1
      \big\vert      &%%l1
      {\mathfrak{c}} &%%c2 {c}
      \big\vert      &%%l2
      {\mathfrak{1}} &%% 1 {a,b,c}
      {\mathfrak{1}} &%% 2 {a,b}
      {\mathfrak{B}} &%% 3 {a,c}
      {\mathfrak{C}} &%% 4 {b,c}
      {\mathfrak{B}} &%% 5 {a}
      {\mathfrak{C}} &%% 6 {b}
      {\mathfrak{c}} &%% 7 {c}
      {\mathfrak{c}} &%% 8 {}
      \big\vert       %%l3
      \\ \nonumber
      {}             &%%c1
      \big\vert      &%%l1
      {\mathfrak{0}} &%%c2 {}
      \big\vert      &%%l2
      {\mathfrak{1}} &%% 1 {a,b,c}
      {\mathfrak{A}} &%% 2 {a,b}
      {\mathfrak{B}} &%% 3 {a,c}
      {\mathfrak{C}} &%% 4 {b,c}
      {\mathfrak{a}} &%% 5 {a}
      {\mathfrak{b}} &%% 6 {b}
      {\mathfrak{c}} &%% 7 {c}
      {\mathfrak{0}} &%% 8 {}
      \big\vert       %%l3
      \\ \nonumber
    \end{array}
\end{equation}





	La tabla para el producto, a continuación.




\begin{equation}
  \begin{array}{@{}*{22}c@{}}
      {\textsf{a\(\cdot\)b}}&%%c1 operacion
      \big\vert     &%%l1 vertical line
      {}            &%%c2 linea de ¿?
      \big\vert     &%%l2 vertical line
      {}            &%% 1
      {}            &%% 2
      {}            &%% 3
      {}            &%% 4
      {\textsf{b}}  &%% 5 variable b
      {}            &%% 6
      {}            &%% 7
      {}            &%% 8
      \big\vert      %%l3
      \\ \hline \nonumber
      {}            &%%c1
      \big\vert     &%%l1
      {}            &%%c2
      \big\vert     &%%l2
      {\mathfrak{1}}&%% 1 {a,b,c}
      {\mathfrak{A}}&%% 2 {a,b}
      {\mathfrak{B}}&%% 3 {a,c}
      {\mathfrak{C}}&%% 4 {b,c}
      {\mathfrak{a}}&%% 5 {a}
      {\mathfrak{b}}&%% 6 {b}
      {\mathfrak{c}}&%% 7 {c}
      {\mathfrak{0}}&%% 8 {}
      \big\vert      %%l3
      \\ \hline \nonumber
      {}            &%%c1
      \big\vert     &%%l1
      {\mathfrak{1}}&%%c2 {a,b,c} lleno
      \big\vert     &%%l2
      {\mathfrak{1}}&%% 1 {a,b,c}
      {\mathfrak{A}}&%% 2 {a,b}
      {\mathfrak{B}}&%% 3 {a,c}
      {\mathfrak{C}}&%% 4 {b,c}
      {\mathfrak{a}}&%% 5 {a}
      {\mathfrak{b}}&%% 6 {b}
      {\mathfrak{c}}&%% 7 {c}
      {\mathfrak{0}}&%% 8 {}
      \big\vert      %%l3
      \\ \nonumber
      {}            &%%c1
      \big\vert     &%%l1
      {\mathfrak{A}}&%%c2 {a,b}
      \big\vert     &%%l2
      {\mathfrak{A}}&%% 1 {a,b,c}
      {\mathfrak{A}}&%% 2 {a,b}
      {\mathfrak{a}}&%% 3 {a,c}
      {\mathfrak{b}}&%% 4 {b,c}
      {\mathfrak{a}}&%% 5 {a}
      {\mathfrak{b}}&%% 6 {b}
      {\mathfrak{0}}&%% 7 {c}
      {\mathfrak{0}}&%% 8 {}
      \big\vert      %%l3
      \\ \nonumber
      {}            &%%c1
      \big\vert     &%%l1
      {\mathfrak{B}}&%%c2 {a,c}
      \big\vert     &%%l2
      {\mathfrak{B}}&%% 1 {a,b,c}
      {\mathfrak{a}}&%% 2 {a,b}
      {\mathfrak{B}}&%% 3 {a,c}
      {\mathfrak{c}}&%% 4 {b,c}
      {\mathfrak{a}}&%% 5 {a}
      {\mathfrak{0}}&%% 6 {b}
      {\mathfrak{c}}&%% 7 {c}
      {\mathfrak{0}}&%% 8 {}
      \big\vert      %%l3
      \\ \nonumber
      {}             &%%c1
      \big\vert      &%%l1
      {\mathfrak{C}} &%%c2 {b,c}
      \big\vert      &%%l2
      {\mathfrak{C}} &%% 1 {a,b,c}
      {\mathfrak{b}} &%% 2 {a,b}
      {\mathfrak{c}} &%% 3 {a,c}
      {\mathfrak{C}} &%% 4 {b,c}
      {\mathfrak{0}} &%% 5 {a}
      {\mathfrak{b}} &%% 6 {b}
      {\mathfrak{c}} &%% 7 {c}
      {\mathfrak{0}} &%% 8 {}
      \big\vert       %%l3
      \\ \nonumber
      {\textsf{a}}   &%%c1 variable a
      \big\vert      &%%l1
      {\mathfrak{a}} &%%c2 {a}
      \big\vert      &%%l2
      {\mathfrak{a}} &%% 1 {a,b,c}
      {\mathfrak{a}} &%% 2 {a,b}
      {\mathfrak{a}} &%% 3 {a,c}
      {\mathfrak{0}} &%% 4 {b,c}
      {\mathfrak{a}} &%% 5 {a}
      {\mathfrak{0}} &%% 6 {b}
      {\mathfrak{0}} &%% 7 {c}
      {\mathfrak{0}} &%% 8 {}
      \big\vert       %%l3
      \\ \nonumber
      {}             &%%c1
      \big\vert      &%%l1
      {\mathfrak{b}} &%%c2 {b}
      \big\vert      &%%l2
      {\mathfrak{b}} &%% 1 {a,b,c}
      {\mathfrak{b}} &%% 2 {a,b}
      {\mathfrak{0}} &%% 3 {a,c}
      {\mathfrak{b}} &%% 4 {b,c}
      {\mathfrak{0}} &%% 5 {a}
      {\mathfrak{b}} &%% 6 {b}
      {\mathfrak{0}} &%% 7 {c}
      {\mathfrak{0}} &%% 8 {}
      \big\vert       %%l3
      \\ \nonumber
      {}             &%%c1
      \big\vert      &%%l1
      {\mathfrak{c}} &%%c2 {c}
      \big\vert      &%%l2
      {\mathfrak{c}} &%% 1 {a,b,c}
      {\mathfrak{0}} &%% 2 {a,b}
      {\mathfrak{c}} &%% 3 {a,c}
      {\mathfrak{c}} &%% 4 {b,c}
      {\mathfrak{0}} &%% 5 {a}
      {\mathfrak{0}} &%% 6 {b}
      {\mathfrak{c}} &%% 7 {c}
      {\mathfrak{0}} &%% 8 {}
      \big\vert       %%l3
      \\ \nonumber
      {}             &%%c1
      \big\vert      &%%l1
      {\mathfrak{0}} &%%c2 {}
      \big\vert      &%%l2
      {\mathfrak{0}} &%% 1 {a,b,c}
      {\mathfrak{0}} &%% 2 {a,b}
      {\mathfrak{0}} &%% 3 {a,c}
      {\mathfrak{0}} &%% 4 {b,c}
      {\mathfrak{0}} &%% 5 {a}
      {\mathfrak{0}} &%% 6 {b}
      {\mathfrak{0}} &%% 7 {c}
      {\mathfrak{0}} &%% 8 {}
      \big\vert       %%l3
      \\ \nonumber
    \end{array}
\end{equation}



	Por último la tabla de los complementos.



\begin{equation}
  \begin{array}{@{}*{22}c@{}}
      {\textsf{a}}                 &%%c2
      \big\vert                    &%%l2
      {\mathfrak{1}}               &%% 1 {a,b,c}
      {\mathfrak{A}}               &%% 2 {a,b}
      {\mathfrak{B}}               &%% 3 {a,c}
      {\mathfrak{C}}               &%% 4 {b,c}
      {\mathfrak{a}}               &%% 5 {a}
      {\mathfrak{b}}               &%% 6 {b}
      {\mathfrak{c}}               &%% 7 {c}
      {\mathfrak{0}}               &%% 8 {}
      \big\vert                     %%l3
      \\ \hline \nonumber
      {\neg\textsf{a}}             &%%c2
      \big\vert                    &%%l2
      {\mathfrak{0}}               &%% 1 {a,b,c}
      {\mathfrak{c}}               &%% 2 {a,b}
      {\mathfrak{b}}               &%% 3 {a,c}
      {\mathfrak{a}}               &%% 4 {b,c}
      {\mathfrak{C}}               &%% 5 {a}
      {\mathfrak{B}}               &%% 6 {b}
      {\mathfrak{A}}               &%% 7 {c}
      {\mathfrak{1}}               &%% 8 {}
      \big\vert                     %%l3
      \\ \nonumber
    \end{array}
\end{equation}



	El diagrama de orden (parcial) de Hasse se hace un poco enrevesado, pero merece la pena verlo para hacerse una idea.



\begin{center}
\begin{tikzpicture}
    \node (top) at ( 2,0) {$ \mathfrak 1 $};
    \node (ndA) at ( 0,1) {$ \mathfrak A $};
    \node (ndB) at ( 2,1) {$ \mathfrak B $};
    \node (ndC) at ( 4,1) {$ \mathfrak C $};
    \node (and) at ( 0,3) {$ \mathfrak a $};
    \node (bnd) at ( 2,3) {$ \mathfrak b $};
    \node (cnd) at ( 4,3) {$ \mathfrak c $};
    \node (bot) at ( 2,4) {$ \mathfrak 0 $};
    \draw [green,  thick, {stealth}-] (top) -- (ndA);%%{a,b,c}
    \draw [green,  thick, {stealth}-] (top) -- (ndB);%%{a,b,d}
    \draw [green,  thick, {stealth}-] (top) -- (ndC);%%{a,c,d}
    \draw [red,  thick, {stealth}-] (ndA) -- (and);%%{a,b,c} <- {a,b} 2
    \draw [red,  thick, {stealth}-] (ndA) -- (bnd);%%{a,b,c} <- {a,c} 3
    \draw [red,  thick, {stealth}-] (ndB) -- (and);%%{a,b,d} <- {a,b} 2
    \draw [red,  thick, {stealth}-] (ndB) -- (cnd);%%{a,b,d} <- {a,d} 4
    \draw [red,  thick, {stealth}-] (ndC) -- (bnd);%%{a,c,d} <- {a,c} 3
    \draw [red,  thick, {stealth}-] (ndC) -- (cnd);%%{a,c,d} <- {a,d} 4
    \draw [blue,  thick, {stealth}-] (and) -- (bot);%%{a,b} <- {a}
    \draw [blue,  thick, {stealth}-] (bnd) -- (bot);%%{a,b} <- {b}
    \draw [blue,  thick, {stealth}-] (cnd) -- (bot);%%{a,c} <- {a}
\end{tikzpicture}
\end{center}



\subsection{Ejemplo: el álgebra de conjuntos de cardinalidad 4.}



	Esta es el siguiente álgebra de Boole: 


\begin{multline}
\setB \mathdef {\setB}_{16} = \\ \nonumber
\qquad\left\lbrace   \set{a,b,c,d},\right. \\ \nonumber
\qquad\:             \set{a,b,c},\set{a,b,d},\set{a,c,d},\set{b,c,d},\\ \nonumber
\qquad\:             \set{a,b},\set{a,c},\set{a,d},\set{b,c},\set{b,d},\\ \nonumber
\qquad\:       \set{c,d},\set{a},\set{b},\set{c},\set{d},\\ \nonumber
\qquad \left.\set{\phantom{v}} \right\rbrace \\ \nonumber
\end{multline}



	Las operaciones, como en los casos anteriores, las concretaremos en
tablas y proceden de la definición del álgebra de las partes de un conjunto
ya explicada anteriormente.



	Llamamos \( \setB_{16} \mathdef \) \large \( \mathscr{P}\! \) \normalsize
\( \mathsf{U}_4 \), dónde \( \mathsf{U}_4 = \set{a,b,c,d} \). De esta forma
nos queda


\begin{flalign*}
\mbox{\large \( \mathscr{P} \)} \mbox{\normalsize \(\mathsf{U}_4 =\)}&
\left\lbrace\mathsf{U}_4,\right. \\
\qquad &\set{a,b,c},\set{a,b,d},\set{a,c,d},\set{b,c,d}, \\
\qquad &\set{a,b},\set{a,c},\set{a,d},\set{b,c},\set{b,d},\set{c,d}, \\
\qquad &\set{a},\set{b},\set{c},\set{d},\\
\qquad &\left.\emptyset\right\rbrace \\
\end{flalign*}



Nombraremos, por comodidad, \(
\mathfrak{a} \mathdef \set{a} \), \( \mathfrak{b} \mathdef \set{b} \) ,  \(
\mathfrak{c} \mathdef \set{c} \) y \( \mathfrak{d} \mathdef \set{d} \).
También nombraremos, \( \mathfrak{D} \mathdef \set{a,b,c} \), \( \mathfrak{C}
\mathdef \set{a,b,d} \), \( \mathfrak{B} \mathdef \set{a,c,d} \) y \(
\mathfrak{A} \mathdef \set{b,c,d} \). Para la capas central tomaremos las
siguientes denominaciones, \( \mathfrak{2} \mathdef \set{a,b} \), \(
\mathfrak{3} \mathdef \set{a,c} \), \( \mathfrak{4} \mathdef \set{a,d} \), \(
\mathfrak{5} \mathdef \set{b,c} \), \( \mathfrak{6} \mathdef \set{b,d} \), \(
\mathfrak{7} \mathdef \set{c,d} \).Claramente \(\mathfrak{1} \mathdef
1_{\setB_{16}} \mathdef \mathsf{U}_4 \) y \(\mathfrak{0} \mathdef
0_{\setB_{16}} \mathdef \emptyset \). De esta forma nuestra notación para el
conjunto de Boole y sus elementos será:



\begin{equation}
\setB_{16} =
\set{
\mathfrak{0},
\mathfrak{A},\mathfrak{B},\mathfrak{C},\mathfrak{D},
\mathfrak{2},\mathfrak{3},\mathfrak{4},
\mathfrak{5},\mathfrak{6},\mathfrak{7},
\mathfrak{a},\mathfrak{b},\mathfrak{c},\mathfrak{d},
\mathfrak{1}
}
\end{equation}



	La tabla para la suma, a continuación.



\begin{equation}
  \begin{array}{@{}*{22}c@{}}
      {\textsf{a+b}}&%%c1 operacion
      \big\vert       &%%l1 vertical line
      {}            &%%c2 linea de ¿?
      \big\vert       &%%l2 vertical line
      {}            &%% 1
      {}            &%% 2
      {}            &%% 3
      {}            &%% 4
      {}            &%% 5
      {}            &%% 6
      {\textsf{b}}  &%% 7 variable b
      {}            &%% 8
      {}            &%% 9
      {}            &%%10
      {}            &%%11
      {}            &%%12
      {}            &%%13
      {}            &%%14
      {}            &%%15
      {}            &%%16
      \big\vert        %%l3
      \\ \hline \nonumber
      {}            &%%c1
      \big\vert       &%%l1
      {}            &%%c2
      \big\vert       &%%l2
      {\mathfrak{1}}&%% 1 {a,b,c,d}
      {\mathfrak{A}}&%% 2 {a,b,c}
      {\mathfrak{B}}&%% 3 {a,b,d}
      {\mathfrak{C}}&%% 4 {a,c,d}
      {\mathfrak{D}}&%% 5 {b,c,d}
      {\mathfrak{2}}      &%% 6 {a,b}
      {\mathfrak{3}}       &%% 7 {a,c}
      {\mathfrak{4}}      &%% 8 {a,d}
      {\mathfrak{5}}      &%% 9 {b,c}
      {\mathfrak{6}}    &%%10 {b,d}
      {\mathfrak{7}}       &%%11 {c,d}
      {\mathfrak{a}}&%%12 {a}
      {\mathfrak{b}}&%%13 {b}
      {\mathfrak{c}}&%%14 {c}
      {\mathfrak{d}}&%%15 {d}
      {\mathfrak{0}}&%%16 {}
      \big\vert        %%l3
      \\ \hline \nonumber
      {}            &%%c1
      \big\vert       &%%l1
      {\mathfrak{1}}&%%c2 {a,b,c,d} lleno
      \big\vert       &%%l2
      {\mathfrak{1}}&%% 1 {a,b,c,d}
      {\mathfrak{1}}&%% 2 {a,b,c}
      {\mathfrak{1}}&%% 3 {a,b,d}
      {\mathfrak{1}}&%% 4 {a,c,d}
      {\mathfrak{1}}&%% 5 {b,c,d}
      {\mathfrak{1}}&%% 6 {a,b}
      {\mathfrak{1}}&%% 7 {a,c}
      {\mathfrak{1}}&%% 8 {a,d}
      {\mathfrak{1}}&%% 9 {b,c}
      {\mathfrak{1}}&%%10 {b,d}
      {\mathfrak{1}}&%%11 {c,d}
      {\mathfrak{1}}&%%12 {a}
      {\mathfrak{1}}&%%13 {b}
      {\mathfrak{1}}&%%14 {c}
      {\mathfrak{1}}&%%15 {d}
      {\mathfrak{1}}&%%16 {}
      \big\vert        %%l3
      \\ \nonumber
      {}            &%%c1
      \big\vert       &%%l1
      {\mathfrak{A}}&%%c2 {a,b,c} falta la d
      \big\vert       &%%l2
      {\mathfrak{1}}&%% 1 {a,b,c,d}
      {\mathfrak{A}}&%% 2 {a,b,c}
      {\mathfrak{1}}&%% 3 {a,b,d}
      {\mathfrak{1}}&%% 4 {a,c,d}
      {\mathfrak{1}}&%% 5 {b,c,d}
      {\mathfrak{A}}&%% 6 {a,b}
      {\mathfrak{A}}&%% 7 {a,c}
      {\mathfrak{1}}&%% 8 {a,d}
      {\mathfrak{A}}&%% 9 {b,c}
      {\mathfrak{1}}&%%10 {b,d}
      {\mathfrak{1}}&%%11 {c,d}
      {\mathfrak{A}}&%%12 {a}
      {\mathfrak{A}}&%%13 {b}
      {\mathfrak{A}}&%%14 {c}
      {\mathfrak{1}}&%%15 {d}
      {\mathfrak{A}}&%%16 {}
      \big\vert        %%l2
      \\ \nonumber
      {}            &%%c1
      \big\vert       &%%l1
      {\mathfrak{B}}&%%c2 {a,b,d} flata la c
      \big\vert       &%%l2
      {\mathfrak{1}}&%% 1 {a,b,c,d}
      {\mathfrak{1}}&%% 2 {a,b,c}
      {\mathfrak{B}}&%% 3 {a,b,d}
      {\mathfrak{1}}&%% 4 {a,c,d}
      {\mathfrak{1}}&%% 5 {b,c,d}
      {\mathfrak{B}}&%% 6 {a,b}
      {\mathfrak{1}}&%% 7 {a,c}
      {\mathfrak{B}}&%% 8 {a,d}
      {\mathfrak{1}}&%% 9 {b,c}
      {\mathfrak{B}}&%%10 {b,d}
      {\mathfrak{1}}&%%11 {c,d}
      {\mathfrak{B}}&%%12 {a}
      {\mathfrak{B}}&%%13 {b}
      {\mathfrak{1}}&%%14 {c}
      {\mathfrak{B}}&%%15 {d}
      {\mathfrak{B}}&%%16 {}
      \big\vert        %%l3
      \\ \nonumber
      {}             &%%c1
      \big\vert        &%%l1
      {\mathfrak{C}} &%%c2 {a,c,d} falta la b
      \big\vert        &%%l2
      {\mathfrak{1}} &%% 1 {a,b,c,d}
      {\mathfrak{1}} &%% 2 {a,b,c}
      {\mathfrak{1}} &%% 3 {a,b,d}
      {\mathfrak{C}} &%% 4 {a,c,d}
      {\mathfrak{1}} &%% 5 {b,c,d}
      {\mathfrak{1}} &%% 6 {a,b}
      {\mathfrak{C}} &%% 7 {a,c}
      {\mathfrak{C}} &%% 8 {a,d}
      {\mathfrak{1}} &%% 9 {b,c}
      {\mathfrak{1}} &%%10 {b,d}
      {\mathfrak{C}} &%%11 {c,d}
      {\mathfrak{C}} &%%12 {a}
      {\mathfrak{1}} &%%13 {b}
      {\mathfrak{C}} &%%14 {c}
      {\mathfrak{C}} &%%15 {d}
      {\mathfrak{C}} &%%16 {}
      \big\vert         %%l3
      \\ \nonumber
      {}             &%%c1
      \big\vert        &%%l1
      {\mathfrak{D}} &%%c2 {b,c,d} falta la a
      \big\vert        &%%l2
      {\mathfrak{1}} &%% 1 {a,b,c,d}
      {\mathfrak{1}} &%% 2 {a,b,c}
      {\mathfrak{1}} &%% 3 {a,b,d}
      {\mathfrak{1}} &%% 4 {a,c,d}
      {\mathfrak{D}} &%% 5 {b,c,d}
      {\mathfrak{1}} &%% 6 {a,b}
      {\mathfrak{1}} &%% 7 {a,c}
      {\mathfrak{1}} &%% 8 {a,d}
      {\mathfrak{D}} &%% 9 {b,c}
      {\mathfrak{D}} &%%10 {b,d}
      {\mathfrak{D}} &%%11 {c,d}
      {\mathfrak{1}} &%%12 {a}
      {\mathfrak{D}} &%%13 {b}
      {\mathfrak{D}} &%%14 {c}
      {\mathfrak{D}} &%%15 {d}
      {\mathfrak{D}} &%%16 {}
      \big\vert         %%l3
      \\ \nonumber
      {}             &%%c1
      \big\vert        &%%l1
      {\mathfrak{2}}       &%%c2 {a,b} falta c y d
      \big\vert        &%%l2
      {\mathfrak{1}} &%% 1 {a,b,c,d}
      {\mathfrak{A}} &%% 2 {a,b,c}
      {\mathfrak{B}} &%% 3 {a,b,d}
      {\mathfrak{1}} &%% 4 {a,c,d}
      {\mathfrak{1}} &%% 5 {b,c,d}
      {\mathfrak{2}}       &%% 6 {a,b}
      {\mathfrak{A}} &%% 7 {a,c}
      {\mathfrak{B}} &%% 8 {a,d}
      {\mathfrak{A}} &%% 9 {b,c}
      {\mathfrak{B}} &%%10 {b,d}
      {\mathfrak{1}} &%%11 {c,d}
      {\mathfrak{2}}       &%%12 {a}
      {\mathfrak{2}}       &%%13 {b}
      {\mathfrak{A}} &%%14 {c}
      {\mathfrak{B}} &%%15 {d}
      {\mathfrak{2}}       &%%16 {}
      \big\vert         %%l3
      \\ \nonumber
      {}             &%%c1
      \big\vert        &%%l1
      {\mathfrak{3}}        &%%c2 {a,c} falta a y d
      \big\vert        &%%l2
      {\mathfrak{1}} &%% 1 {a,b,c,d}
      {\mathfrak{A}} &%% 2 {a,b,c}
      {\mathfrak{1}} &%% 3 {a,b,d}
      {\mathfrak{C}} &%% 4 {a,c,d}
      {\mathfrak{1}} &%% 5 {b,c,d}
      {\mathfrak{A}}       &%% 6 {a,b}
      {\mathfrak{3}}        &%% 7 {a,c}
      {\mathfrak{C}}       &%% 8 {a,d}
      {\mathfrak{A}}       &%% 9 {b,c}
      {\mathfrak{1}}     &%%10 {b,d}
      {\mathfrak{C}}        &%%11 {c,d}
      {\mathfrak{3}} &%%12 {a}
      {\mathfrak{A}} &%%13 {b}
      {\mathfrak{3}} &%%14 {c}
      {\mathfrak{C}} &%%15 {d}
      {\mathfrak{3}} &%%16 {}
      \big\vert         %%l3
      \\ \nonumber
      {\textsf{a}}   &%%c1
      \big\vert        &%%l1
      {\mathfrak{4}}       &%%c2 {a,d} falta b y c
      \big\vert        &%%l2
      {\mathfrak{1}} &%% 1 {a,b,c,d}
      {\mathfrak{1}} &%% 2 {a,b,c}
      {\mathfrak{B}} &%% 3 {a,b,d}
      {\mathfrak{C}} &%% 4 {a,c,d}
      {\mathfrak{1}} &%% 5 {b,c,d}
      {\mathfrak{B}}       &%% 6 {a,b}
      {\mathfrak{C}}        &%% 7 {a,c}
      {\mathfrak{4}} &%% 8 {a,d}
      {\mathfrak{1}}       &%% 9 {b,c}
      {\mathfrak{B}}     &%%10 {b,d}
      {\mathfrak{C}}        &%%11 {c,d}
      {\mathfrak{4}} &%%12 {a}
      {\mathfrak{B}} &%%13 {b}
      {\mathfrak{C}} &%%14 {c}
      {\mathfrak{4}} &%%15 {d}
      {\mathfrak{4}} &%%16 {}
      \big\vert         %%l3
      \\ \nonumber
      {}             &%%c1
      \big\vert        &%%l1
      {\mathfrak{5}}       &%%c2 {b,c} falta a y d
      \big\vert        &%%l2
      {\mathfrak{1}} &%% 1 {a,b,c,d}
      {\mathfrak{A}} &%% 2 {a,b,c}
      {\mathfrak{1}} &%% 3 {a,b,d}
      {\mathfrak{1}} &%% 4 {a,c,d}
      {\mathfrak{D}} &%% 5 {b,c,d}
      {\mathfrak{A}}       &%% 6 {a,b}
      {\mathfrak{A}}        &%% 7 {a,c}
      {\mathfrak{1}} &%% 8 {a,d}
      {\mathfrak{5}}       &%% 9 {b,c}
      {\mathfrak{D}}     &%%10 {b,d}
      {\mathfrak{D}}        &%%11 {c,d}
      {\mathfrak{A}} &%%12 {a}
      {\mathfrak{5}} &%%13 {b}
      {\mathfrak{5}} &%%14 {c}
      {\mathfrak{D}} &%%15 {d}
      {\mathfrak{5}} &%%16 {}
      \big\vert        %%l2
      \\ \nonumber
      {}             &%%c1
      \big\vert        &%%l1
      {\mathfrak{6}}     &%%c2 {b,d} falta a y c
      \big\vert        &%%l2
      {\mathfrak{1}} &%% 1 {a,b,c,d}
      {\mathfrak{1}} &%% 2 {a,b,c}
      {\mathfrak{B}} &%% 3 {a,b,d}
      {\mathfrak{1}} &%% 4 {a,c,d}
      {\mathfrak{D}} &%% 5 {b,c,d}
      {\mathfrak{B}}       &%% 6 {a,b}
      {\mathfrak{1}}        &%% 7 {a,c}
      {\mathfrak{B}} &%% 8 {a,d}
      {\mathfrak{D}}       &%% 9 {b,c}
      {\mathfrak{6}}     &%%10 {b,d}
      {\mathfrak{D}}        &%%11 {c,d}
      {\mathfrak{B}} &%%12 {a}
      {\mathfrak{6}} &%%13 {b}
      {\mathfrak{D}} &%%14 {c}
      {\mathfrak{6}} &%%15 {d}
      {\mathfrak{6}} &%%16 {}
      \big\vert         %%l3
      \\ \nonumber
      {}             &%%c1
      \big\vert        &%%l1
      {\mathfrak{7}}        &%%c2 {c,d} falta a y b
      \big\vert        &%%l2
      {\mathfrak{1}} &%% 1 {a,b,c,d}
      {\mathfrak{1}} &%% 2 {a,b,c}
      {\mathfrak{1}} &%% 3 {a,b,d}
      {\mathfrak{C}} &%% 4 {a,c,d}
      {\mathfrak{D}} &%% 5 {b,c,d}
      {\mathfrak{1}}       &%% 6 {a,b}
      {\mathfrak{C}}        &%% 7 {a,c}
      {\mathfrak{C}} &%% 8 {a,d}
      {\mathfrak{D}}       &%% 9 {b,c}
      {\mathfrak{D}}     &%%10 {b,d}
      {\mathfrak{7}}        &%%11 {c,d}
      {\mathfrak{C}} &%%12 {a}
      {\mathfrak{D}} &%%13 {b}
      {\mathfrak{7}} &%%14 {c}
      {\mathfrak{7}} &%%15 {d}
      {\mathfrak{7}} &%%16 {}
      \big\vert         %%l3
      \\ \nonumber
      {}             &%%c1
      \big\vert        &%%l1
      {\mathfrak{a}} &%%c2 {a} falta b c y d
      \big\vert        &%%l2
      {\mathfrak{1}} &%% 1 {a,b,c,d}
      {\mathfrak{A}} &%% 2 {a,b,c}
      {\mathfrak{B}} &%% 3 {a,b,d}
      {\mathfrak{C}} &%% 4 {a,c,d}
      {\mathfrak{1}} &%% 5 {b,c,d}
      {\mathfrak{2}}       &%% 6 {a,b}
      {\mathfrak{3}}        &%% 7 {a,c}
      {\mathfrak{4}} &%% 8 {a,d}
      {\mathfrak{A}}       &%% 9 {b,c}
      {\mathfrak{B}}     &%%10 {b,d}
      {\mathfrak{C}}        &%%11 {c,d}
      {\mathfrak{a}} &%%12 {a}
      {\mathfrak{2}} &%%13 {b}
      {\mathfrak{3}} &%%14 {c}
      {\mathfrak{4}} &%%15 {d}
      {\mathfrak{a}} &%%16 {}
      \big\vert         %%l3
      \\ \nonumber
      {}             &%%c1
      \big\vert        &%%l1
      {\mathfrak{b}} &%%c2
      \big\vert        &%%l2 {b} falta a c y d
      {\mathfrak{1}} &%% 1 {a,b,c,d}
      {\mathfrak{A}} &%% 2 {a,b,c}
      {\mathfrak{B}} &%% 3 {a,b,d}
      {\mathfrak{1}} &%% 4 {a,c,d}
      {\mathfrak{D}} &%% 5 {b,c,d}
      {\mathfrak{2}}       &%% 6 {a,b}
      {\mathfrak{A}}        &%% 7 {a,c}
      {\mathfrak{B}} &%% 8 {a,d}
      {\mathfrak{5}}       &%% 9 {b,c}
      {\mathfrak{6}}     &%%10 {b,d}
      {\mathfrak{D}}        &%%11 {c,d}
      {\mathfrak{2}} &%%12 {a}
      {\mathfrak{b}} &%%13 {b}
      {\mathfrak{5}} &%%14 {c}
      {\mathfrak{6}} &%%15 {d}
      {\mathfrak{b}} &%%16 {}
      \big\vert         %%l3
      \\ \nonumber
      {}             &%%c
      \big\vert        &%%l1
      {\mathfrak{c}} &%%c2 {c} falta a b y d
      \big\vert        &%%l2
      {\mathfrak{1}} &%% 1 {a,b,c,d}
      {\mathfrak{A}} &%% 2 {a,b,c}
      {\mathfrak{1}} &%% 3 {a,b,d}
      {\mathfrak{C}} &%% 4 {a,c,d}
      {\mathfrak{D}} &%% 5 {b,c,d}
      {\mathfrak{A}}       &%% 6 {a,b}
      {\mathfrak{3}}        &%% 7 {a,c}
      {\mathfrak{C}} &%% 8 {a,d}
      {\mathfrak{5}}       &%% 9 {b,c}
      {\mathfrak{B}}     &%%10 {b,d}
      {\mathfrak{7}}        &%%11 {c,d}
      {\mathfrak{3}} &%%12 {a}
      {\mathfrak{5}} &%%13 {b}
      {\mathfrak{c}} &%%14 {c}
      {\mathfrak{7}} &%%15 {d}
      {\mathfrak{c}} &%%16 {}
      \big\vert         %%l3
      \\ \nonumber
      {}             &%%c1
      \big\vert        &%%l1
      {\mathfrak{d}} &%%c2 {d} falta a b y c
      \big\vert        &%%l2
      {\mathfrak{1}} &%% 1 {a,b,c,d}
      {\mathfrak{1}} &%% 2 {a,b,c}
      {\mathfrak{B}} &%% 3 {a,b,d}
      {\mathfrak{C}} &%% 4 {a,c,d}
      {\mathfrak{D}} &%% 5 {b,c,d}
      {\mathfrak{B}}       &%% 6 {a,b}
      {\mathfrak{C}}        &%% 7 {a,c}
      {\mathfrak{4}} &%% 8 {a,d}
      {\mathfrak{D}}       &%% 9 {b,c}
      {\mathfrak{6}}     &%%10 {b,d}
      {\mathfrak{7}}        &%%11 {c,d}
      {\mathfrak{4}} &%%12 {a}
      {\mathfrak{6}} &%%13 {b}
      {\mathfrak{7}} &%%14 {c}
      {\mathfrak{d}} &%%15 {d}
      {\mathfrak{d}} &%%16 {}
      \big\vert         %%l3
      \\ \nonumber
      {}             &%%c1
      \big\vert        &%%l1
      {\mathfrak{0}} &%%c2 {} flata todo
      \big\vert        &%%l2
      {\mathfrak{1}} &%% 1 {a,b,c,d}
      {\mathfrak{A}} &%% 2 {a,b,c}
      {\mathfrak{B}} &%% 3 {a,b,d}
      {\mathfrak{C}} &%% 4 {a,c,d}
      {\mathfrak{D}} &%% 5 {b,c,d}
      {\mathfrak{2}}       &%% 6 {a,b}
      {\mathfrak{3}}        &%% 7 {a,c}
      {\mathfrak{4}} &%% 8 {a,d}
      {\mathfrak{5}}       &%% 9 {b,c}
      {\mathfrak{6}}     &%%10 {b,d}
      {\mathfrak{7}}        &%%11 {c,d}
      {\mathfrak{a}} &%%12 {a}
      {\mathfrak{b}} &%%13 {b}
      {\mathfrak{c}} &%%14 {c}
      {\mathfrak{d}} &%%15 {d}
      {\mathfrak{0}} &%%16 {}
      \big\vert         %%l3
      \\ \hline \nonumber
    \end{array}
\end{equation}



	La tabla para el producto, a continuación.


\begin{equation}
  \begin{array}{@{}*{22}c@{}}
      {\textsf{a}*\textsf{b}}&%%c1 operacion
      \big\vert                    &%%l1 vertical line
      {}                           &%%c2 linea de ¿?
      \big\vert                    &%%l2 vertical line
      {}                           &%% 1
      {}                           &%% 2
      {}                           &%% 3
      {}                           &%% 4
      {}                           &%% 5
      {}                           &%% 6
      {\textsf{b}}                 &%% 7 variable b
      {}                           &%% 8
      {}                           &%% 9
      {}                           &%%10
      {}                           &%%11
      {}                           &%%12
      {}                           &%%13
      {}                           &%%14
      {}                           &%%15
      {}                           &%%16
      \big\vert                     %%l3
      \\ \hline \nonumber
      {}                           &%%c1
      \big\vert                    &%%l1
      {}                           &%%c2
      \big\vert                    &%%l2
      {\mathfrak{1}}               &%% 1 {a,b,c,d}
      {\mathfrak{A}}               &%% 2 {a,b,c}
      {\mathfrak{B}}               &%% 3 {a,b,d}
      {\mathfrak{C}}               &%% 4 {a,c,d}
      {\mathfrak{D}}               &%% 5 {b,c,d}
      {\mathfrak{2}}               &%% 6 {a,b}
      {\mathfrak{3}}               &%% 7 {a,c}
      {\mathfrak{4}}               &%% 8 {a,d}
      {\mathfrak{5}}               &%% 9 {b,c}
      {\mathfrak{6}}               &%%10 {b,d}
      {\mathfrak{7}}               &%%11 {c,d}
      {\mathfrak{a}}               &%%12 {a}
      {\mathfrak{b}}               &%%13 {b}
      {\mathfrak{c}}               &%%14 {c}
      {\mathfrak{d}}               &%%15 {d}
      {\mathfrak{0}}               &%%16 {}
      \big\vert                     %%l3
      \\ \hline \nonumber
      {}                           &%%c1
      \big\vert                    &%%l1
      {\mathfrak{1}}               &%%c2 {a,b,c,d} lleno
      \big\vert                    &%%l2
      {\mathfrak{1}}               &%% 1 {a,b,c,d}
      {\mathfrak{A}}               &%% 2 {a,b,c}
      {\mathfrak{B}}               &%% 3 {a,b,d}
      {\mathfrak{C}}               &%% 4 {a,c,d}
      {\mathfrak{D}}               &%% 5 {b,c,d}
      {\mathfrak{2}}               &%% 6 {a,b}
      {\mathfrak{3}}               &%% 7 {a,c}
      {\mathfrak{4}}               &%% 8 {a,d}
      {\mathfrak{5}}               &%% 9 {b,c}
      {\mathfrak{6}}               &%%10 {b,d}
      {\mathfrak{7}}               &%%11 {c,d}
      {\mathfrak{a}}               &%%12 {a}
      {\mathfrak{b}}               &%%13 {b}
      {\mathfrak{c}}               &%%14 {c}
      {\mathfrak{d}}               &%%15 {d}
      {\mathfrak{0}}               &%%16 {}
      \big\vert                     %%l3
      \\ \nonumber
      {}                           &%%c1
      \big\vert                    &%%l1
      {\mathfrak{A}}               &%%c2 {a,b,c} falta la d
      \big\vert                    &%%l2
      {\mathfrak{A}}               &%% 1 {a,b,c,d}
      {\mathfrak{A}}               &%% 2 {a,b,c}
      {\mathfrak{2}}               &%% 3 {a,b,d}
      {\mathfrak{3}}               &%% 4 {a,c,d}
      {\mathfrak{1}}               &%% 5 {b,c,d}
      {\mathfrak{A}}               &%% 6 {a,b}
      {\mathfrak{A}}               &%% 7 {a,c}
      {\mathfrak{a}}               &%% 8 {a,d}
      {\mathfrak{5}}               &%% 9 {b,c}
      {\mathfrak{b}}               &%%10 {b,d}
      {\mathfrak{c}}               &%%11 {c,d}
      {\mathfrak{a}}               &%%12 {a}
      {\mathfrak{b}}               &%%13 {b}
      {\mathfrak{c}}               &%%14 {c}
      {\mathfrak{0}}               &%%15 {d}
      {\mathfrak{0}}               &%%16 {}
      \big\vert                     %%l2
      \\ \nonumber
      {}            &%%c1
      \big\vert     &%%l1
      {\mathfrak{B}}&%%c2 {a,b,d} flata la c
      \big\vert     &%%l2
      {\mathfrak{B}}&%% 1 {a,b,c,d}
      {\mathfrak{2}}&%% 2 {a,b,c}
      {\mathfrak{B}}&%% 3 {a,b,d}
      {\mathfrak{4}}&%% 4 {a,c,d}
      {\mathfrak{6}}&%% 5 {b,c,d}
      {\mathfrak{2}}&%% 6 {a,b}
      {\mathfrak{a}}&%% 7 {a,c}
      {\mathfrak{4}}&%% 8 {a,d}
      {\mathfrak{b}}&%% 9 {b,c}
      {\mathfrak{6}}&%%10 {b,d}
      {\mathfrak{d}}&%%11 {c,d}
      {\mathfrak{a}}&%%12 {a}
      {\mathfrak{b}}&%%13 {b}
      {\mathfrak{0}}&%%14 {c}
      {\mathfrak{d}}&%%15 {d}
      {\mathfrak{0}}&%%16 {}
      \big\vert      %%l3
      \\ \nonumber
      {}             &%%c1
      \big\vert      &%%l1
      {\mathfrak{C}} &%%c2 {a,c,d} falta la b
      \big\vert      &%%l2
      {\mathfrak{C}} &%% 1 {a,b,c,d}
      {\mathfrak{B}} &%% 2 {a,b,c}
      {\mathfrak{4}} &%% 3 {a,b,d}
      {\mathfrak{C}} &%% 4 {a,c,d}
      {\mathfrak{7}} &%% 5 {b,c,d}
      {\mathfrak{a}} &%% 6 {a,b}
      {\mathfrak{3}} &%% 7 {a,c}
      {\mathfrak{4}} &%% 8 {a,d}
      {\mathfrak{c}} &%% 9 {b,c}
      {\mathfrak{d}} &%%10 {b,d}
      {\mathfrak{7}} &%%11 {c,d}
      {\mathfrak{a}} &%%12 {a}
      {\mathfrak{0}} &%%13 {b}
      {\mathfrak{c}} &%%14 {c}
      {\mathfrak{d}} &%%15 {d}
      {\mathfrak{0}} &%%16 {}
      \big\vert       %%l3
      \\ \nonumber
      {}             &%%c1
      \big\vert      &%%l1
      {\mathfrak{D}} &%%c2 {b,c,d} falta la a
      \big\vert      &%%l2
      {\mathfrak{D}} &%% 1 {a,b,c,d}
      {\mathfrak{5}} &%% 2 {a,b,c}
      {\mathfrak{6}} &%% 3 {a,b,d}
      {\mathfrak{7}} &%% 4 {a,c,d}
      {\mathfrak{D}} &%% 5 {b,c,d}
      {\mathfrak{b}} &%% 6 {a,b}
      {\mathfrak{c}} &%% 7 {a,c}
      {\mathfrak{d}} &%% 8 {a,d}
      {\mathfrak{5}} &%% 9 {b,c}
      {\mathfrak{6}} &%%10 {b,d}
      {\mathfrak{7}} &%%11 {c,d}
      {\mathfrak{0}} &%%12 {a}
      {\mathfrak{b}} &%%13 {b}
      {\mathfrak{c}} &%%14 {c}
      {\mathfrak{d}} &%%15 {d}
      {\mathfrak{0}} &%%16 {}
      \big\vert       %%l3
      \\ \nonumber
      {}             &%%c1
      \big\vert      &%%l1
      {\mathfrak{2}} &%%c2 {a,b} falta c y d
      \big\vert      &%%l2
      {\mathfrak{2}} &%% 1 {a,b,c,d}
      {\mathfrak{2}} &%% 2 {a,b,c}
      {\mathfrak{2}} &%% 3 {a,b,d}
      {\mathfrak{a}} &%% 4 {a,c,d}
      {\mathfrak{b}} &%% 5 {b,c,d}
      {\mathfrak{2}} &%% 6 {a,b}
      {\mathfrak{a}} &%% 7 {a,c}
      {\mathfrak{a}} &%% 8 {a,d}
      {\mathfrak{b}} &%% 9 {b,c}
      {\mathfrak{b}} &%%10 {b,d}
      {\mathfrak{0}} &%%11 {c,d}
      {\mathfrak{a}} &%%12 {a}
      {\mathfrak{b}} &%%13 {b}
      {\mathfrak{0}} &%%14 {c}
      {\mathfrak{0}} &%%15 {d}
      {\mathfrak{0}} &%%16 {}
      \big\vert       %%l3
      \\ \nonumber
      {}             &%%c1
      \big\vert      &%%l1
      {\mathfrak{3}} &%%c2 {a,c} falta a y d
      \big\vert      &%%l2
      {\mathfrak{3}} &%% 1 {a,b,c,d}
      {\mathfrak{3}} &%% 2 {a,b,c}
      {\mathfrak{a}} &%% 3 {a,b,d}
      {\mathfrak{3}} &%% 4 {a,c,d}
      {\mathfrak{c}} &%% 5 {b,c,d}
      {\mathfrak{a}} &%% 6 {a,b}
      {\mathfrak{3}} &%% 7 {a,c}
      {\mathfrak{a}} &%% 8 {a,d}
      {\mathfrak{c}} &%% 9 {b,c}
      {\mathfrak{0}} &%%10 {b,d}
      {\mathfrak{c}} &%%11 {c,d}
      {\mathfrak{a}} &%%12 {a}
      {\mathfrak{0}} &%%13 {b}
      {\mathfrak{c}} &%%14 {c}
      {\mathfrak{0}} &%%15 {d}
      {\mathfrak{0}} &%%16 {}
      \big\vert       %%l3
      \\ \nonumber
      {\textsf{a}}   &%%c1
      \big\vert      &%%l1
      {\mathfrak{4}} &%%c2 {a,d} falta b y c
      \big\vert      &%%l2
      {\mathfrak{4}} &%% 1 {a,b,c,d}
      {\mathfrak{a}} &%% 2 {a,b,c}
      {\mathfrak{4}} &%% 3 {a,b,d}
      {\mathfrak{4}} &%% 4 {a,c,d}
      {\mathfrak{d}} &%% 5 {b,c,d}
      {\mathfrak{a}} &%% 6 {a,b}
      {\mathfrak{a}} &%% 7 {a,c}
      {\mathfrak{4}} &%% 8 {a,d}
      {\mathfrak{0}} &%% 9 {b,c}
      {\mathfrak{d}} &%%10 {b,d}
      {\mathfrak{d}} &%%11 {c,d}
      {\mathfrak{a}} &%%12 {a}
      {\mathfrak{0}} &%%13 {b}
      {\mathfrak{0}} &%%14 {c}
      {\mathfrak{d}} &%%15 {d}
      {\mathfrak{0}} &%%16 {}
      \big\vert       %%l3
      \\ \nonumber
      {}             &%%c1
      \big\vert      &%%l1
      {\mathfrak{5}} &%%c2 {b,c} falta a y d
      \big\vert      &%%l2
      {\mathfrak{5}} &%% 1 {a,b,c,d}
      {\mathfrak{5}} &%% 2 {a,b,c}
      {\mathfrak{b}} &%% 3 {a,b,d}
      {\mathfrak{c}} &%% 4 {a,c,d}
      {\mathfrak{5}} &%% 5 {b,c,d}
      {\mathfrak{b}} &%% 6 {a,b}
      {\mathfrak{c}} &%% 7 {a,c}
      {\mathfrak{0}} &%% 8 {a,d}
      {\mathfrak{5}} &%% 9 {b,c}
      {\mathfrak{b}} &%%10 {b,d}
      {\mathfrak{c}} &%%11 {c,d}
      {\mathfrak{0}} &%%12 {a}
      {\mathfrak{b}} &%%13 {b}
      {\mathfrak{c}} &%%14 {c}
      {\mathfrak{0}} &%%15 {d}
      {\mathfrak{0}} &%%16 {}
      \big\vert       %%l2
      \\ \nonumber
      {}             &%%c1
      \big\vert      &%%l1
      {\mathfrak{6}} &%%c2 {b,d} falta a y c
      \big\vert      &%%l2
      {\mathfrak{6}} &%% 1 {a,b,c,d}
      {\mathfrak{b}} &%% 2 {a,b,c}
      {\mathfrak{6}} &%% 3 {a,b,d}
      {\mathfrak{d}} &%% 4 {a,c,d}
      {\mathfrak{6}} &%% 5 {b,c,d}
      {\mathfrak{b}} &%% 6 {a,b}
      {\mathfrak{0}} &%% 7 {a,c}
      {\mathfrak{d}} &%% 8 {a,d}
      {\mathfrak{b}} &%% 9 {b,c}
      {\mathfrak{6}} &%%10 {b,d}
      {\mathfrak{d}} &%%11 {c,d}
      {\mathfrak{0}} &%%12 {a}
      {\mathfrak{b}} &%%13 {b}
      {\mathfrak{0}} &%%14 {c}
      {\mathfrak{d}} &%%15 {d}
      {\mathfrak{0}} &%%16 {}
      \big\vert       %%l3
      \\ \nonumber
      {}             &%%c1
      \big\vert      &%%l1
      {\mathfrak{7}} &%%c2 {c,d} falta a y b
      \big\vert      &%%l2
      {\mathfrak{7}} &%% 1 {a,b,c,d}
      {\mathfrak{c}} &%% 2 {a,b,c}
      {\mathfrak{d}} &%% 3 {a,b,d}
      {\mathfrak{7}} &%% 4 {a,c,d}
      {\mathfrak{7}} &%% 5 {b,c,d}
      {\mathfrak{0}} &%% 6 {a,b}
      {\mathfrak{c}} &%% 7 {a,c}
      {\mathfrak{d}} &%% 8 {a,d}
      {\mathfrak{c}} &%% 9 {b,c}
      {\mathfrak{d}} &%%10 {b,d}
      {\mathfrak{7}} &%%11 {c,d}
      {\mathfrak{0}} &%%12 {a}
      {\mathfrak{0}} &%%13 {b}
      {\mathfrak{c}} &%%14 {c}
      {\mathfrak{d}} &%%15 {d}
      {\mathfrak{0}} &%%16 {}
      \big\vert       %%l3
      \\ \nonumber
      {}             &%%c1
      \big\vert      &%%l1
      {\mathfrak{a}} &%%c2 {a} falta b c y d
      \big\vert      &%%l2
      {\mathfrak{a}} &%% 1 {a,b,c,d}
      {\mathfrak{a}} &%% 2 {a,b,c}
      {\mathfrak{a}} &%% 3 {a,b,d}
      {\mathfrak{a}} &%% 4 {a,c,d}
      {\mathfrak{0}} &%% 5 {b,c,d}
      {\mathfrak{a}} &%% 6 {a,b}
      {\mathfrak{a}} &%% 7 {a,c}
      {\mathfrak{a}} &%% 8 {a,d}
      {\mathfrak{0}} &%% 9 {b,c}
      {\mathfrak{0}} &%%10 {b,d}
      {\mathfrak{0}} &%%11 {c,d}
      {\mathfrak{a}} &%%12 {a}
      {\mathfrak{0}} &%%13 {b}
      {\mathfrak{0}} &%%14 {c}
      {\mathfrak{0}} &%%15 {d}
      {\mathfrak{0}} &%%16 {}
      \big\vert       %%l3
      \\ \nonumber
      {}             &%%c1
      \big\vert      &%%l1
      {\mathfrak{b}} &%%c2
      \big\vert      &%%l2 {b} falta a c y d
      {\mathfrak{b}} &%% 1 {a,b,c,d}
      {\mathfrak{b}} &%% 2 {a,b,c}
      {\mathfrak{b}} &%% 3 {a,b,d}
      {\mathfrak{0}} &%% 4 {a,c,d}
      {\mathfrak{b}} &%% 5 {b,c,d}
      {\mathfrak{b}} &%% 6 {a,b}
      {\mathfrak{0}} &%% 7 {a,c}
      {\mathfrak{0}} &%% 8 {a,d}
      {\mathfrak{b}} &%% 9 {b,c}
      {\mathfrak{b}} &%%10 {b,d}
      {\mathfrak{0}} &%%11 {c,d}
      {\mathfrak{0}} &%%12 {a}
      {\mathfrak{b}} &%%13 {b}
      {\mathfrak{0}} &%%14 {c}
      {\mathfrak{0}} &%%15 {d}
      {\mathfrak{0}} &%%16 {}
      \big\vert       %%l3
      \\ \nonumber
      {}             &%%c
      \big\vert      &%%l1
      {\mathfrak{c}} &%%c2 {c} falta a b y d
      \big\vert      &%%l2
      {\mathfrak{c}} &%% 1 {a,b,c,d}
      {\mathfrak{c}} &%% 2 {a,b,c}
      {\mathfrak{0}} &%% 3 {a,b,d}
      {\mathfrak{c}} &%% 4 {a,c,d}
      {\mathfrak{c}} &%% 5 {b,c,d}
      {\mathfrak{0}} &%% 6 {a,b}
      {\mathfrak{c}} &%% 7 {a,c}
      {\mathfrak{0}} &%% 8 {a,d}
      {\mathfrak{c}} &%% 9 {b,c}
      {\mathfrak{0}} &%%10 {b,d}
      {\mathfrak{c}} &%%11 {c,d}
      {\mathfrak{0}} &%%12 {a}
      {\mathfrak{0}} &%%13 {b}
      {\mathfrak{c}} &%%14 {c}
      {\mathfrak{0}} &%%15 {d}
      {\mathfrak{0}} &%%16 {}
      \big\vert       %%l3
      \\ \nonumber
      {}             &%%c1
      \big\vert      &%%l1
      {\mathfrak{d}} &%%c2 {d} falta a b y c
      \big\vert      &%%l2
      {\mathfrak{d}} &%% 1 {a,b,c,d}
      {\mathfrak{0}} &%% 2 {a,b,c}
      {\mathfrak{d}} &%% 3 {a,b,d}
      {\mathfrak{d}} &%% 4 {a,c,d}
      {\mathfrak{d}} &%% 5 {b,c,d}
      {\mathfrak{0}} &%% 6 {a,b}
      {\mathfrak{0}} &%% 7 {a,c}
      {\mathfrak{d}} &%% 8 {a,d}
      {\mathfrak{0}} &%% 9 {b,c}
      {\mathfrak{d}} &%%10 {b,d}
      {\mathfrak{d}} &%%11 {c,d}
      {\mathfrak{0}} &%%12 {a}
      {\mathfrak{0}} &%%13 {b}
      {\mathfrak{0}} &%%14 {c}
      {\mathfrak{d}} &%%15 {d}
      {\mathfrak{0}} &%%16 {}
      \big\vert       %%l3
      \\ \nonumber
      {}             &%%c1
      \big\vert      &%%l1
      {\mathfrak{0}} &%%c2 {} flata todo
      \big\vert      &%%l2
      {\mathfrak{0}} &%% 1 {a,b,c,d}
      {\mathfrak{0}} &%% 2 {a,b,c}
      {\mathfrak{0}} &%% 3 {a,b,d}
      {\mathfrak{0}} &%% 4 {a,c,d}
      {\mathfrak{0}} &%% 5 {b,c,d}
      {\mathfrak{0}} &%% 6 {a,b}
      {\mathfrak{0}} &%% 7 {a,c}
      {\mathfrak{0}} &%% 8 {a,d}
      {\mathfrak{0}} &%% 9 {b,c}
      {\mathfrak{0}} &%%10 {b,d}
      {\mathfrak{0}} &%%11 {c,d}
      {\mathfrak{0}} &%%12 {a}
      {\mathfrak{0}} &%%13 {b}
      {\mathfrak{0}} &%%14 {c}
      {\mathfrak{0}} &%%15 {d}
      {\mathfrak{0}} &%%16 {}
      \big\vert       %%l3
      \\ \hline \nonumber
    \end{array}
\end{equation}



	Por último la tabla de los complementos.



\begin{equation}
  \begin{array}{@{}*{22}c@{}}
      {\textsf{a}}                 &%%c2
      \big\vert                    &%%l2
      {\mathfrak{1}}               &%% 1 {a,b,c,d}
      {\mathfrak{A}}               &%% 2 {a,b,c}
      {\mathfrak{B}}               &%% 3 {a,b,d}
      {\mathfrak{C}}               &%% 4 {a,c,d}
      {\mathfrak{D}}               &%% 5 {b,c,d}
      {\mathfrak{2}}               &%% 6 {a,b}
      {\mathfrak{3}}               &%% 7 {a,c}
      {\mathfrak{4}}               &%% 8 {a,d}
      {\mathfrak{5}}               &%% 9 {b,c}
      {\mathfrak{6}}               &%%10 {b,d}
      {\mathfrak{7}}               &%%11 {c,d}
      {\mathfrak{a}}               &%%12 {a}
      {\mathfrak{b}}               &%%13 {b}
      {\mathfrak{c}}               &%%14 {c}
      {\mathfrak{d}}               &%%15 {d}
      {\mathfrak{0}}               &%%16 {}
      \big\vert                     %%l3
      \\ \hline \nonumber
      {\neg\textsf{a}}             &%%c2
      \big\vert                    &%%l2
      {\mathfrak{0}}               &%% 1 {a,b,c,d}
      {\mathfrak{d}}               &%% 2 {a,b,c}
      {\mathfrak{c}}               &%% 3 {a,b,d}
      {\mathfrak{b}}               &%% 4 {a,c,d}
      {\mathfrak{a}}               &%% 5 {b,c,d}
      {\mathfrak{7}}               &%% 6 {a,b}
      {\mathfrak{6}}               &%% 7 {a,c}
      {\mathfrak{5}}               &%% 8 {a,d}
      {\mathfrak{4}}               &%% 9 {b,c}
      {\mathfrak{3}}               &%%10 {b,d}
      {\mathfrak{2}}               &%%11 {c,d}
      {\mathfrak{D}}               &%%12 {a}
      {\mathfrak{C}}               &%%13 {b}
      {\mathfrak{B}}               &%%14 {c}
      {\mathfrak{A}}               &%%15 {d}
      {\mathfrak{1}}               &%%16 {}
      \big\vert                     %%l3
      \\ \nonumber
    \end{array}
\end{equation}



	El diagrama de orden (parcial) de Hasse se hace un poco enrevesado, pero
merece la pena verlo para hacerse una idea.

\begin{center}
\begin{tikzpicture}
    \node (top) at ( 6,1) {$ \mathfrak 1 $};
    \node (ndA) at ( 3,2) {$ \mathfrak A $};
    \node (ndB) at ( 5,2) {$ \mathfrak B $};
    \node (ndC) at ( 7,2) {$ \mathfrak C $};
    \node (ndD) at ( 9,2) {$ \mathfrak D $};
    \node (nd2) at ( 1,4) {$ \mathfrak 2 $};
    \node (nd3) at ( 3,4) {$ \mathfrak 3 $};
    \node (nd4) at ( 5,4) {$ \mathfrak 4 $};
    \node (nd5) at ( 7,4) {$ \mathfrak 5 $};
    \node (nd6) at ( 9,4) {$ \mathfrak 6 $};
    \node (nd7) at (11,4) {$ \mathfrak 7 $};
    \node (and) at ( 3,6) {$ \mathfrak a $};
    \node (bnd) at ( 5,6) {$ \mathfrak b $};
    \node (cnd) at ( 7,6) {$ \mathfrak c $};
    \node (dnd) at ( 9,6) {$ \mathfrak d $};
    \node (bot) at ( 6,7) {$ \mathfrak 0 $};
    \draw [green,  thick, {stealth}-] (top) -- (ndA);%%{a,b,c}
    \draw [green,  thick, {stealth}-] (top) -- (ndB);%%{a,b,d}
    \draw [green,  thick, {stealth}-] (top) -- (ndC);%%{a,c,d}
    \draw [green,  thick, {stealth}-] (top) -- (ndD);%%{b,c,d}
    \draw [red,  thick, {stealth}-] (ndA) -- (nd2);%%{a,b,c} <- {a,b} 2
    \draw [red,  thick, {stealth}-] (ndA) -- (nd3);%%{a,b,c} <- {a,c} 3
    \draw [red,  thick, {stealth}-] (ndA) -- (nd5);%%{a,b,c} <- {b,c} 5
    \draw [red,  thick, {stealth}-] (ndB) -- (nd2);%%{a,b,d} <- {a,b} 2
    \draw [red,  thick, {stealth}-] (ndB) -- (nd4);%%{a,b,d} <- {a,d} 4
    \draw [red,  thick, {stealth}-] (ndB) -- (nd6);%%{a,b,d} <- {b,d} 6
    \draw [red,  thick, {stealth}-] (ndC) -- (nd3);%%{a,c,d} <- {a,c} 3
    \draw [red,  thick, {stealth}-] (ndC) -- (nd4);%%{a,c,d} <- {a,d} 4
    \draw [red,  thick, {stealth}-] (ndC) -- (nd7);%%{a,c,d} <- {c,d} 7
    \draw [red,  thick, {stealth}-] (ndD) -- (nd5);%%{b,c,d} <- {b,c} 5
    \draw [red,  thick, {stealth}-] (ndD) -- (nd6);%%{b,c,d} <- {b,d} 6
    \draw [red,  thick, {stealth}-] (ndD) -- (nd7);%%{b,c,d} <- {c,d} 7
    \draw [blue,  thick, {stealth}-] (nd2) -- (and);%%{a,b} <- {a}
    \draw [blue,  thick, {stealth}-] (nd2) -- (bnd);%%{a,b} <- {b}
    \draw [blue,  thick, {stealth}-] (nd3) -- (and);%%{a,c} <- {a}
    \draw [blue,  thick, {stealth}-] (nd3) -- (cnd);%%{a,c} <- {c}
    \draw [blue,  thick, {stealth}-] (nd4) -- (and);%%{a,d} <- {a}
    \draw [blue,  thick, {stealth}-] (nd4) -- (dnd);%%{a,d} <- {d}
    \draw [blue,  thick, {stealth}-] (nd5) -- (bnd);%%{b,c} <- {b}
    \draw [blue,  thick, {stealth}-] (nd5) -- (cnd);%%{b,c} <- {c}
    \draw [blue,  thick, {stealth}-] (nd6) -- (bnd);%%{b,d} <- {b}
    \draw [blue,  thick, {stealth}-] (nd6) -- (dnd);%%{b,d} <- {d}
    \draw [blue,  thick, {stealth}-] (nd7) -- (cnd);%%{c,d} <- {c}
    \draw [blue,  thick, {stealth}-] (nd7) -- (dnd);%%{c,d} <- {d}
    \draw [orange,  thick, {stealth}-] (and) -- (bot);
    \draw [orange,  thick, {stealth}-] (bnd) -- (bot);
    \draw [orange,  thick, {stealth}-] (cnd) -- (bot);
    \draw [orange,  thick, {stealth}-] (dnd) -- (bot);
\end{tikzpicture}
\end{center}



\subsection{Álgebra de las uniones finitas de intervalos disjuntos en los
racionales.}



	Presentamos el álgebra de Boole de los conjuntos que se pueden expresar
como unión disjunta finita de sub-intervalos genéricos de
\(\left[0,1\right]\cap\setQ\) . Definimos, por conveniencia, \(I \mathdef
{\left[ {0,1} \right]}_{\setQ} \). Para esto haremos abstracción de cualquier
conjunto finito de puntos de \(I_{\setQ}\) , esto es, consideraremos que dos
conjuntos son equivalentes si su diferencia simétrica (la unión de las
diferencias, los elementos que no son comunes de ambos conjuntos) es vacía o
es un conjunto finito de puntos. Esta álgebra de Boole tiene un cardinal
infinito numerable (como el cardinal de los números racionales es igual que
el de los naturales). Lo más importante es que no puede desarrollarse de
manera semejante a como desarrollamos el álgebra de las partes de un
conjunto. Lo formalizaremos del siguiente modo:



\newcommand {\powset}[1]{
	\ensuremath{
		\text{
			\large \(\mathscr{P}\) \normalsize 
				\( #1 \)
			}
		}
	}
\newcommand {\IQ} {\ensuremath{{\mathsf{I}}_{\setQ}}}
\newcommand {\IIQ} {\ensuremath{{\mathsf{I}\mathsf{I}}_{\setQ}}}
\newcommand {\IIIQ} {\ensuremath{{\mathsf{I}\mathsf{I}\mathsf{I}}_{\setQ}}}
\newcommand {\IVQ} {\ensuremath{{\mathsf{I}\mathsf{V}}_{\setQ}}}
\newcommand {\FinIQ} {\ensuremath{{\mathsf{Fin}}_{\IQ}}}
\newcommand {\FinDiffRel} {\ensuremath{\!\!\!\!\mod\FinIQ}}
\newcommand {\ClEqFinDiff}[1]{
	\ensuremath{
		{
			\left[\!\left[
				{ #1 }
			\right]\!\right]
		}_{
			\FinDiffRel
		}
	}
}
\newcommand {\QuotSetEqRel}[1]{\ensuremath{#1 \big{/} {\!\!\FinDiffRel}}}



\begin{enumerate} {
  \item {Definimos el intervalo universo \(\IQ\) y los intervalos internos
  \(\left[ \cdot , \cdot \right]_{\setQ}\).
  
  
  
  \begin{equation}
  \IQ \quad \mathdef \quad \setCompr{x \in \setQ}{0 \leq x \leq 1}
  \end{equation}
  
  
  
  \begin{equation}
  \forall a, b \in \IQ, a \leq b \then { \left[ {a,b} \right] }_{\setQ}
  \quad \mathdef \quad { \left[ {a,b} \right] } \cap \IQ = \setCompr{x \in
  \IQ} {a \le x \le b}
  \end{equation}
  }}
  
  
  
  \item { Si escribimos \( { \left[ { a , b } \right] }_{ \setQ } \), la
  expresión tendrá sentido solo si \( 0 \leq a \leq b \leq 1 \), y \(a \in
  \setQ \) y \( b \in \setQ \).}
  \item {Ahora definimos \( \IIQ \mathdef \setCompr{ { \left[ { a , b }
  \right] }_{ \setQ } } { a \in \IQ \ly b \in\IQ \ly a \leq b } \), que
  representa el conjunto de los intervalos en \(\IQ\).}
  
  
  
  \item {Nos hará falta definir los conjuntos finitos de \(\IQ\), \( \FinIQ
  \).
  \begin{equation}
  \FinIQ \mathdef \setCompr{A \in \powset{\IQ} } { \natural A \in \setN
  \cup \set{0} }
  \end{equation}
  }
  
  
  
  \item { Para evitar problemas con las uniones o intersecciones de
  intervalos en sus límites (cuando estos coincidan) etc. vamos crear una
  relación de equivalencia con la que emparentar cada conjunto de \(\setQ\)
  con todos los que se diferencien en a lo sumo un conjunto finito de
  valores de \(\IQ\), y que llamaremos \(\mod \FinIQ\). Previamente
  definimos la diferencia simétrica de conjuntos:
  
  
  
  \begin{equation}
  A \text{ is a } \mathsf{set}, B \text{ is a } \mathsf{set} \;\then\; A \triangle B \mathdef \left( A \setminus B \right) \cup \left( B \setminus A \right)
  \end{equation}
  
  
  
  Ahora podemos definir nuestro objetivo:
  
  
  
  \begin{equation}
  A , B \in \powset{\left( \IQ \right)} \;\then\; A \cong B \FinDiffRel 
  \mathdef  \natural \left( A \triangle B \right) \in \setN \cup \set{0}
  \end{equation}
  
  
  
  Veremos que la relación \( \FinDiffRel \) entre subconjuntos de \(\IQ\) 
  es de equivalencia.
  \begin{enumerate}
  \item { La relación es \( \FinDiffRel \) \textsf{reflexiva} de forma
  trivial:
  \begin{align}
  \forall \;A\; &\in \IIQ \\ 
  A \triangle A = \left( A \setminus A \right) &\cup \left( A \setminus A
  \right) = \emptyset \\ 
  \natural \left(A \triangle A \right) = \natural \emptyset = 0 \in \setN
  \cup \set{0} &\;\;\Leftrightarrow\;\; A \cong A \FinDiffRel
  \end{align}
  }
  
  
  
  \item { La relación es \textsf{simétrica}:
  
  
  
  \begin{equation}
  A \in \IIQ, B \in \IIQ \; A \triangle B = \left( A \setminus B \right)
  \cup \left( B \setminus A \right) =  \left( B \setminus A \right) \cup
  \left( A \setminus B \right) = B \triangle A
  \end{equation}
  
  
  
  	ya que la unión es conmutativa.
  }
  
  
  
  \item { La relación es \textsf{transitiva}. Primero estableceremos una 
      propiedad de las uniones \( \left( A \triangle B \right) \cup 
      \left( B \triangle C \right) \). La propiedad la podemos enunciar: 
      dados \( A \in \IIQ \) y \( C \in \IIQ \), \(\forall B \in \IIQ A 
      \triangle C \subseteq \left( A \triangle B \right) \cup \left( B 
      \triangle C \right) \). Se parece a la desigualdad triangular de 
      las distancias métricas.
  \begin{align}
  &\forall E \in \IIQ \;\; \overline{E} \;\mathdef\; \IQ \setminus E\\
  &E \;=\; E \cap \IQ \;=\; E \cap \left( F \cup \overline{F} \right) \\
  &\overline{G} \;=\; \overline{G} \cap \left( F \cup \overline{F} 
  \right) \\
  &E \cap \overline{G} \;=\; E \cap \left( F \cup \overline{F} \right) 
  \cap \overline{G} \\ \nonumber
  &\phantom{E \cap \overline{G} }\;=\; \left( \left( E \cap F \right) 
  \cup \left( E \cap \overline{F} \right) \right) \cap \overline{G}\\ 
  \nonumber
  &\phantom{E \cap \overline{G} }\;=\; \left( E \cap F \cap \overline{G} 
  \right) \cup \left( E \cap \overline{F} \cap \overline{G} \right) \\  
  \nonumber
  &\phantom{E \cap \overline{G} }\;\subseteq\; \left( F \cap \overline{G}
  \right) \cup \left( E \cap \overline{F} \right) \\
  &E \setminus G \;\subseteq\; \left( F \setminus G \right) \cup \left( E
  \setminus F \right) \\
  &G \setminus E \;\subseteq\; \left( F \setminus E \right) \cup \left( G
  \setminus F \right) \\
  &\left( E \setminus G\right) \cup \left( G \setminus E \right)
  \;\subseteq\; \left( \left( F \setminus G \right) \cup \left( G
  \setminus F \right) \right) \cup \left( \left( E \setminus F \right)
  \cup \left( F \setminus E \right) \right) \\
  &\forall F \in \IIQ \;\;\; E \triangle G \subseteq \left( E \triangle F
  \right) \cup \left( F \triangle G \right)
  \end{align}
  
  
  
  Ahora, la transitividad de la relación \(\FinDiffRel\) queda clara, \(
  \natural \left( A \triangle C \right) \leq \natural \left( A \triangle
  B \right) + \natural \left( B \triangle C \right) = n_{AB} + n_{BC} \in
  \setN \cup \set{0}\), con lo que ya sabemos que \( A \triangle C \in
  \FinIQ \) y llegamos a la conclusión:
  
  
  
  \begin{align}
  &\forall A, B, C \in \IIQ\\ \nonumber
  &\left( A \cong B \FinDiffRel \right) \land \left( B \cong C
  \FinDiffRel \right) \;\then\; A \cong C \FinDiffRel
  \end{align}
  }
  \end{enumerate}
  
  
  
  }
  
  
  
  \item {Definimos ahora un nuevo conjunto, el de las uniones finitas de
  conjuntos en \(\IIQ\) (intervalos):
  
  
  
  \begin{equation}
  \IIIQ \mathdef \setCompr
  { A \in \powset{\left( \IIQ
  \right) } }  
  { A = {{\bigcup}_{\lambda\in\Lambda}} {\mathsf{I}}_{\lambda}, 
  \forall \lambda \in \Lambda, 
  {{\mathsf{I}}_{\lambda}} \in \IIQ,  
  \natural{ \Lambda } \in \left( \setN \cup \set{0} \right) }
  \end{equation}
  \begin{equation}
  \forall J \in \IIIQ \quad \ClEqFinDiff{J} \;\;\mathdef\;\;
  \setCompr{K\in\IIIQ}{K \cong J \FinDiffRel}
  \end{equation}
  }
  
  
  
  \item {Ahora ya podemos definir nuestro conjunto definitivo como 
      subyacente al álgebra de Boole, el de los elementos en \(\IIIQ\) 
      (uniones finitas de intervalos de \(\IIQ\)) cociente la relación de 
      equivalencia \(\FinDiffRel\):
  
  
  
  \begin{equation}
  \IVQ \mathdef \QuotSetEqRel{\IIIQ} = \setCompr{\ClEqFinDiff{J}}{J \in 
  \IIIQ}
  \end{equation}
  
  
  
  }
\end{enumerate} 



	Pero, ¿cuál es nuestro álgebra de Boole \( \mathscr{B} \)? Primero
tenemos el conjunto base, \( \setB \mathdef \IVQ \). Como operaciones
binarias buleanas escogemos, dados \(\alpha\in\IVQ\) y \(\beta\in\IVQ\),
dónde \(\exists A \in \IIIQ \;\; \alpha = \ClEqFinDiff{A}\) y \(\exists B \in
\IIIQ \;\; \beta = \ClEqFinDiff{B}\), 



\begin{equation}
\setB \mathdef \IVQ
\end{equation}



\begin{equation} 
\alpha + \beta \mathdef \ClEqFinDiff{A \cup B} 
\end{equation}



\begin{equation} 
\alpha \cdot \beta \mathdef \ClEqFinDiff{A \cap B} 
\end{equation} 



	Tomando ahora \( 0_{\setB} \mathdef \ClEqFinDiff{ \emptyset } \) que será
elemento neutro de la suma y \( 1_{\setB} \mathdef \ClEqFinDiff{ \IQ } \) que
será el neutro del producto. Dado una unión finita de intervalos de \(\IQ\),
digamos \(A\), el complementario de \( \overline{\ClEqFinDiff{ A }} =
\ClEqFinDiff{ \IQ \setminus A } \).



	Lo primero es ver que las operaciones están bien definidas. 



	Sean \(A,B,P,Q \in \IIIQ\), de forma que se cumple \(A \cong P
\FinDiffRel\) y \(B \cong Q \FinDiffRel\). Se trata de ver que
\(\ClEqFinDiff{A}+\ClEqFinDiff{B}=\ClEqFinDiff{P}+\ClEqFinDiff{Q}\). Las
trataremos como conjuntos sin detenernos por ahora en si la unión de clases
de intervalos disjuntos es una clase de unión de intervalos disjuntos. Esto
lo dejaremos para un momento posterior. Por la definición:



\begin{align}
&\ClEqFinDiff{A}+\ClEqFinDiff{B} = \ClEqFinDiff{A\cup B}\\
&\ClEqFinDiff{P}+\ClEqFinDiff{Q} = \ClEqFinDiff{P\cup Q}
\end{align}



\begin{align}
&\left(A \cup B\right)\setminus\left(P \cup Q\right) =  \\
&\qquad\qquad = \left(A \cup B\right)\cap\overline{\left(P \cup Q\right)} =
\\
&\qquad\qquad = \left(A \cup B\right)\cap {\left(\overline{P} \cap
\overline{Q}\right)} = \\
&\qquad\qquad = \left( A \cap { \left( \overline{P} \cap \overline{Q}
\right)} \right) \cup \left( B \cap { \left( \overline{P} \cap \overline{Q}
\right) } \right) = \\
&\qquad\qquad = \left( \left( A \cap \overline{P} \right) \cap \overline{Q}
\right) \cup \left( B \cap \left( \overline{Q} \cap \overline{P} \right)
\right) \subseteq \\
&\qquad\qquad = \left( \left( A \cap \overline{P} \right) \cap \overline{Q}
\right) \cup \left( \left( B \cap \overline{Q} \right) \cap \overline{P}
\right) \subseteq \\
&\qquad\qquad \subseteq \left( A \cap \overline{P} \right) \cup \left( B \cap
\overline{Q} \right) = \\
&\qquad\qquad \subseteq \left( A \setminus P \right) \cup \left( B \setminus
Q \right) \subseteq \\
&\qquad\qquad \subseteq \left( A \triangle P \right) \cup \left( B \triangle
Q \right)
\end{align}



\begin{align}
&\qquad \text{ Con idéntico razonamiento, por simetría:} \\
& \left(P \cup Q\right) \setminus \left(A \cup B\right) \subseteq \left( A
\triangle P \right) \cup \left( B \triangle Q \right)\\
&\qquad \text{ Uniendo las dos aserciones de contención:} \\
& \left(P \cup Q\right) \triangle \left(A \cup B\right) \subseteq \left( A
\triangle P \right) \cup \left( B \triangle Q \right) \in \FinIQ
\end{align}



\begin{align}
&\qquad\qquad A \cup B \cong P \cup Q \FinDiffRel \\
&\qquad\qquad \ClEqFinDiff{A\cup B} = \ClEqFinDiff{P\cup Q}
\end{align}


	Como conclusión: \emph{la suma no depende de los representantes de la clase elegidos}.


	Sean \(A,B,P,Q \in \IIIQ\), de forma que se cumple \(A \cong P
\FinDiffRel\) y \(B \cong Q \FinDiffRel\). Ahora se trata de ver que
\(\ClEqFinDiff{A}\cdot\ClEqFinDiff{B}=\ClEqFinDiff{P}\cdot\ClEqFinDiff{Q}\).
Las trataremos como conjuntos sin detenernos por ahora en si la intersección
de clases de uniones de intervalos disjuntos es una clase de uniones de
intervalos disjuntos. Esto lo dejaremos para un momento posterior. Por la
definición:



\begin{align}
&\ClEqFinDiff{A}\cdot\ClEqFinDiff{B} = \ClEqFinDiff{A\cap B}\\
&\ClEqFinDiff{P}\cdot\ClEqFinDiff{Q} = \ClEqFinDiff{P\cap Q}
\end{align}
\begin{align}
&\left(A \cap B\right)\setminus\left(P \cap Q\right) =  \\
&\qquad\qquad = \left(A \cap B\right)\cap\overline{\left(P \cap Q\right)} =
\\
&\qquad\qquad = \left(A \cap B\right)\cap {\left(\overline{P} \cup
\overline{Q}\right)} = \\
&\qquad\qquad = \left( \left( A \cap B \right) \cap \overline{P} \right) \cup
\left( \left( A \cap B \right) \cap \overline{Q} \right) = \\
&\qquad\qquad = \left( \left( B \cap A \right) \cap \overline{P} \right) \cup
\left( \left( A \cap B \right) \cap \overline{Q} \right) = \\
&\qquad\qquad = \left( B \cap \left( A \cap \overline{P} \right) \right) \cup
\left( A \cap \left( B \cap \overline{Q} \right) \right) \subseteq \\
&\qquad\qquad \subseteq \left( A \cap \overline{P} \right) \cup \left( B \cap
\overline{Q} \right) \subseteq \\
&\qquad\qquad \subseteq \left( A \setminus P \right) \cup \left( B \setminus
Q \right) \subseteq \\
&\qquad\qquad \subseteq \left( A \triangle P \right) \cup \left( B \triangle
Q \right)
\end{align}
\begin{align}
&\qquad \text{ Con idéntico razonamiento, por simetría:} \\
& \left(P \cap Q\right) \setminus \left(A \cap B\right) \subseteq \left( A
\triangle P \right) \cup \left( B \triangle Q \right)\\
&\qquad \text{ Uniendo las dos aserciones de contención:} \\
& \left(P \cap Q\right) \triangle \left(A \cap B\right) \subseteq \left( A
\triangle P \right) \cup \left( B \triangle Q \right) \in \FinIQ\\
\end{align}



\begin{align}
&\qquad\qquad A \cap B \cong P \cap Q \FinDiffRel \\
&\qquad\qquad \ClEqFinDiff{A\cap B} = \ClEqFinDiff{P\cap Q}
\end{align}


	Como conclusión: \emph{el producto no depende de los representantes elegidos.}




	Aún queda ver si en el sentido restringido de las pruebas de buena
definición anteriores, dados \(A,P \in \IIIQ\), ver que si se cumple \(A
\cong P \FinDiffRel\), también se cumple que
\(\overline{\ClEqFinDiff{A}}=\overline{\ClEqFinDiff{P}}\). Dicho más
desmenuzadamente, que \(\IQ\setminus A \triangle \IQ\setminus P \in \FinIQ\).
Los trataremos como conjuntos sin detenernos por ahora en si el
complementario de clases de uniones de intervalos disjuntos es una clase de
uniones de intervalos disjuntos. Como hasta ahora lo dejaremos para el
siguiente paso. Por la definición:



\begin{align}
&\overline{\ClEqFinDiff{A}} = \ClEqFinDiff{\IQ \setminus A}\\
&\overline{\ClEqFinDiff{P}} = \ClEqFinDiff{\IQ \setminus P}\\
&\ClEqFinDiff{A} = \ClEqFinDiff{P}\\
\end{align}



\begin{align}
&\left( \IQ \setminus A \right) \setminus \left( \IQ \setminus P \right) = \\
&\qquad\qquad = \overline{A} \setminus \overline{P} = \\
&\qquad\qquad = \setCompr{ x \in \IQ } { x \notin A \land x \in P } = \\
&\qquad\qquad = \setCompr{ x \in P } { x \notin A} = \\
&\qquad\qquad = P \setminus A \subseteq \FinIQ
\end{align}




\begin{align}
&\qquad \text{ Con idéntico razonamiento, por simetría:} \\
&\left( \IQ \setminus P \right) \setminus \left( \IQ \setminus A \right)
\subseteq  A \setminus P \subseteq \FinIQ
\end{align}

\begin{align}
&\qquad \text{ Uniendo las dos desigualdades anteriores (contenciones):} \\
\overline{A}\triangle\overline{P} &= \\
\qquad &= \left( \overline{P} \setminus \overline{A} \right) \cup \left(
\overline{A} \setminus \overline{P} \right) = \\
\qquad &= \left( \left( \IQ \setminus P \right) \setminus \left( \IQ
\setminus A \right) \right) \cup \left( \left( \IQ \setminus A \right)
\setminus \left( \IQ \setminus P \right) \right) \subseteq \\
\qquad &\subseteq \left( A \setminus P \right) \cup \left( P \setminus A
\right) = A \triangle P \subseteq \FinIQ
\end{align}

\begin{align}
&\qquad\qquad \overline{A} \cong \overline{P} \FinDiffRel \\
&\qquad\qquad \overline{ \ClEqFinDiff{ A } } = \overline{ \ClEqFinDiff{ P } }
\end{align}


	Conclusión: \emph{el elemento complementario, tal como se ha definido, no depende del representante de la clase elegido.}




	Lo que nos falta para ver que esas tres funciones están bien definidas es
ver que de manera efectiva son elementos del conjunto de clases de
equivalencias de uniones finitas de intervalos. Para eso no hemos definido
aún una forma adecuada de representar sucesiones de intervalos disjuntos en
\(\IQ\). Hasta ahora hemos tratado el símbolo \(\ClEqFinDiff{a,b}\) de una
determinada manera. Variaremos la definición para que se adecúe bien a
nuestros propósitos de representación explícita.



	Sean \(a_0,a_1,\ldots,a_{2\cdot n-1}\) una lista de elementos de \(\IQ\)
dónde \(n\in\setN\) y se cumple que \(\forall \iota \in
\set{0,1,\ldots,2\cdot n-1}\) se cumple que \( 0 \leq a_{\iota} \lneqq
a_{\iota+1} \leq 1 \). Los elementos de \(\IIIQ\) los representaremos como
\(\left[{a_0, a_1, \ldots , a_{2\cdot n - 1}}\right]_{\IIIQ}\) queriendo
decir (o más globalmente):



\(A\in\IIIQ \then A\) tendrá una de las siguientes representaciones:



Representación de tipo \(Repr\;I\) si \(A\) es un intervalo vacío:
\begin{equation}
\left[{}\right]_{\IIIQ} \mathdef \set{} = {\left[{}\right]}_{\IIQ} =
\emptyset
\end{equation}



	Representación de tipo \(Repr\;II\) is \(A\) es una unión de intervalos
de un solo elemento, o, dicho de otra forma, es un conjunto finito de valores
de \(\IQ\):


\begin{eqnarray}
\nonumber
{}&\left[{a_0;a_1;\ldots;a_{n-1}}\right]_{\IIIQ} \mathdef
{\set{a_0,a_1,\ldots,a_{n-1}}}_{\IQ} = \\ \nonumber
=&{\left[{a_0,a_0}\right]}_{\IIQ}\cup{\left[{a_1,a_1}\right]}_{\IIQ}\cup
\ldots\cup{\left[{a_{n-1},a_{n-1}}\right]}_{\IIQ}=\\ \nonumber
=&{\bigcup}_{0\leq{\iota}\leq{{\iota}_{\max}}}
{\left[{a_{\iota},a_{\iota}}\right]}_{\IIQ} =
{\bigcup}_{0\leq{\iota}\leq{{\iota}_{\max}}}
{\set{a_{\iota}}}_{\IQ}
\end{eqnarray}



Representación de tipo \(Repr\;III\) de \(A\) si es una unión de intervalos
de un solo elemento más el intervalo vacío, o, dicho de otra forma, es un
conjunto finito de valores de \(IQ\) que incluye el conjunto vacío:



\begin{align}
\nonumber
{}&\left[{a_0;a_1;\ldots;a_{n-1}};\right]_{\IIIQ} \mathdef
{\set{a_0,a_1,\ldots,a_{n-1}}}_{\IQ} \cup \set{\set{}} = \\ \nonumber
=&{\left[{a_0,a_0}\right]}_{\IIQ}\cup{\left[{a_1,a_1}\right]}_{\IIQ}
\cup\ldots\cup{\left[{a_{n-1},a_{n-1}}\right]}_{\IIQ}\cup
{\left[\right]}_{\IIQ}=\\ \nonumber
=& {\left[\right]}_{\IIQ} \cup {\bigcup}_{0\leq{\iota}\leq{{\iota}_{\max}}}
{\left[{a_{\iota},a_{\iota}}\right]}_{\IIQ} =
\set{\emptyset}\cup{\bigcup}_{0\leq{\iota}\leq{{\iota}_{\max}}}
{\set{a_{\iota}}}_{\IQ}
\end{align}



	Representación de tipo \(Repr\;IV\) de \(A\) si es una unión de
intervalos de propios de racionales en \({\IQ}\) sin contar con otro conjunto
de elemento, por ejemplo el vacío los conjuntos finitos:



\begin{align}
\nonumber
{}&\left[{a_0,a_1,\ldots,a_{2\cdot{{{\iota}_{\max}+1}}}}\right]_{\IIIQ}
\mathdef {\left[{a_0,a_1}\right]}_{\IIQ}\cup{\left[{a_2,a_3}\right]}_{\IIQ}
\cup\ldots\cup{\left[a_{2\cdot{{\iota}_{\max}}},a_{2\cdot{{\iota}_{\max}} 
+1}\right]}_{\IIQ}=\\ \nonumber
=&{\bigcup}_{0\leq{\iota}\lneq{{\iota}_{\max} \text{ impar}}}
{\left[{a_{\iota,par},a_{\iota+1}}\right]}_{\IIQ}
\end{align}



	La representación general de tipo \(Repr\;V\) del conjunto \(A\) de
\(\IIIQ\) si es una unión de intervalos de propios de racionales en \({\IQ}\)
con un conjunto de elementos que es finito, y además el conjunto vacío:



\begin{align}
\nonumber
{}&\left[{a_0,\ldots,a_{2\cdot{{\iota}_{\max}}+1}};b_0;\ldots;b_{n-1};;
\right]_{\IIIQ} \mathdef&
\left[{a_0,\ldots,a_{2\cdot{{\iota}_{\max}}+1}}\right]_{\IIIQ} \cup
{\left[b_0;\ldots;b_{n-1};;\right]}_{\IIIQ}
\end{align}



	Uniones de las representaciones anteriores agrupadas en 
representación de intervalos, representación de números de \(\IQ\) en
cantidad finita y unión de la representación del vacío.



	Ya solo nos queda ver como se representa un conjunto de \(\IVQ\) que
estarán muy simplificados:



\newcommand {\ClEqFinDiffInter}[1]{
	\ensuremath{ 
		{
		\left[\!\left[ 
			{#1} 
		\right]\!\right]
		}_{
		\FinDiffRel
		} 
	}
}



\begin{align}
\ClEqFinDiffInter{a_0,\ldots,a_{2\cdot{{\iota}_{\max}}+1}} &\mathdef&
{\bigcup}_{0\leq \iota \leq {\iota}_{\max}}\ClEqFinDiff{ { a_{2\cdot\iota} ,
a_{2\cdot\iota+1} } }\\
\ClEqFinDiffInter{\;} &\mathdef& \ClEqFinDiff{\emptyset} =
\ClEqFinDiff{\set{a\in\IQ}} =\\
&=&\ClEqFinDiff{\setCompr{a_\iota \in \IQ}{\exists n\in \setN\big{.}\;\;0\leq
\iota \leq n}}
\end{align}



	Bien. Ya estamos preparados para tratar con la suma buleana para los
elementos \(A\in\IVQ\) y \(B\in\IVQ\), tales que \(A =
\ClEqFinDiffInter{a_0,\ldots,a_{2\cdot n+1}}\) y  \(B =
\ClEqFinDiffInter{b_0,\ldots,a_{2\cdot m+1}}\). ¿Cómo representaríamos
\(A+B\)? Probemos desde los casos más fáciles:



	Primero, la forma canónica de la unión de intervalos disjuntos es en una
sucesión de intervalos ordenada, de forma que el límite superior de un
intervalo es estrictamente menor que el límite inferior del intervalo
nombrado a continuación en la expresión de la unión disjunta de intervalos.
Que esto podemos hacerlo así es fácil de ver: \(\left[ a , b \right[ \uplus
\left] b , c \right] \) tiene la misma clase de equivalencia (módulo un
conjunto finito de puntos) que \( \left[ a , c \right] \) ya que \( \left(
\left[ a , b \right[ \uplus \left] b , c \right] \right) \triangle \left[ a ,
c \right] = \set{b}\) y \( \natural\set{b} = 1 \in \setN \). Luego podemos
tomar para el trabajo con intervalos la condición mencionada. Por eso en vez
de trabajar con \( A \in \IVQ \) y con \( B \in \IVQ \), relajaremos el
lenguaje para no tener que hacer continuamente referencia a que son clases de
equivalencia, y trataremos con uniones de intervalos disjuntos en \( \IIIQ \)
dónde los intervalos en la serie de uniones disjuntas están ya ordenados. De
hecho solamente tendremos en mente que cualquier intervalo es de \( \IIQ \),
esto es \( A \subseteq \IQ \) y \( B \subseteq \IQ \). Si que tendremos que
tener cuidado de tratar con el intervalo vacío, \( \left[{\;}\right] \), que
es un valor que pueden tomar \(A\) y \(B\), y que nunca aparece cuando
ponemos \(A\) o \(B\) como una unión de intervalos disjuntos, so pena de
inventar un índice especial para ese intervalo, como \(-1\), pero es
farragoso y no lo trataremos así, sino como un caso especial. Además,
trataremos como otro caso especial cuando \(A\) o \(B\) tomen valor \(\IQ\).
Además de estos casos trataremos como casos base cuando \(A\) o \(B\) sean un
solo intervalo no nulo. De esta manera tendremos construidos nuestros casos
base. El caso general de cualquier operación interna lo definiremos por
recursión, teniendo en cuenta que el número de intervalos disjuntos, tanto en
\(A\) como en \(B\) es finito. También veremos que cuando la unión no es
finita no podemos mantener que la operación sea interna. Por último
utilizaremos el signo de operación binaria \(A \uplus B\) como \( \sup A
\lneqq \inf B \) esto implica \( A \cap B = \emptyset = \left[\;\right] \), y
además, en tanto que conjuntos, equivale a \( A \uplus B \equiv A \cup B \).



	Teniendo lo anterior en cuenta definimos, para el operando izquierdo
\(A\):



\begin{equation}
\exists n \in \setN \;\big{|}\; A \mathdef
\displaystyle{{\biguplus}_{\kappa=0}^{n}} \left[ a_{ 2 \cdot \kappa}, a_{ {
\left( 2 \cdot \kappa \right) } + 1 } \right] \; \bf{ \big{ | } } \; \left[
\; \right]
\end{equation}
\begin{equation}
\textsf{Dónde:    }\quad \forall \imath,\jmath \in \setN \quad \left( 0 \leq
\imath \lneqq \jmath \leq n \right) \then \left( 0 \leq a_{\imath} \lneqq
a_{\jmath} \leq 1\;  \land \; a_{\imath}, a_{\jmath}\in \setQ \right)
\end{equation}



	Igualmente definimos, para el operando derecho \(B\):



\begin{equation}
\exists m \in \setN \;\big{|}\; B \mathdef
\displaystyle{{\biguplus}_{\nu=0}^{m}} \left[
b_{2 \cdot \nu}, b_{{\left( 2 \cdot \nu \right)} + 1} \right] \; \bf{ \big{|}
} \; \left[ \; \right]
\end{equation}
\begin{equation}
\textsf{Dónde:    }\quad \forall \imath,\jmath \in \setN \quad \left( 0 \leq
\imath \lneqq \jmath \leq m \right) \then \left( 0 \leq b_{\imath} \lneqq
b_{\jmath} \leq 1\;  \land \; b_{\imath}, b_{\jmath}\in \setQ \right)
\end{equation}



	Comencemos a construir la suma como una función en un lenguaje funcional
que permita recursividad (Haskell, por ejemplo):



\begin{eqnarray}
\nonumber
&&\text{Casos de elemento neutro de la suma y del producto}\\
&&A + B : \IIIQ \times \IIIQ \fromto \IIIQ\\
&&      \left[\,\right] + \left[\,\right] = \left[\,\right]\\
&&      \left[\,\right] + B = B\\
&&      A + \left[\,\right] = A\\
&&      \IQ + \IQ = \IQ\\
&&      \IQ + B   = \IQ\\
&&      A   + \IQ = \IQ\\ \nonumber
&&      \text{1ª parte.}\;\;\text{Continúa}\ldots
\end{eqnarray}



	Ahora introducimos los casos en que \(A\) toma como representante un solo
intervalo \( A \mathdef \left[ a_0, a_1 \right] \) y \(B\) también \(B
\mathdef \left[ b_0, b_1 \right] \). Habrá varios casos: \(a_0\) es el menor
de los límites, o \(b_0\) es el menor de ellos:



\begin{eqnarray}
\nonumber
\ldots&&\text{Continúa de la 1ª parte.}\\
\text{Caso}\quad n = 1 \;\; m = 1 &&\\
A + B &:& \IIIQ \times \IIIQ \fromto \IIIQ\\
      &&\left[a_0,a_1\right] + \left[b_0,b_1\right] = \qquad\\
      &|& a_0 \lneqq a_1 \lneqq b_0 \lneqq b_1\qquad = A \uplus B\\ 
      &|& a_0 \lneqq b_0 \leqq a_1 \lneqq b_1\qquad = \left[a_0,b_1\right]\\ 
      &|& a_0 \lneqq b_0 \lneqq b_1 \lneqq a_1\qquad = A\\ 
      &|& b_0 \lneqq b_1 \lneqq a_0 \lneqq a_1\qquad = B \uplus A\\ 
      &|& b_0 \lneqq a_0 \leqq b_1 \lneqq a_1\qquad = \left[b_0,a_1\right]\\ 
      &|& b_0 \lneqq a_0 \lneqq a_1 \lneqq b_1\qquad = B\\ \nonumber
      &&\text{2ª parte.}\;\text{Continúa}\ldots
\end{eqnarray}



	Ahora abordaremos el caso general en que \(B\) es una unión finita
cualquiera de intervalos disjuntos y ordenados. Se hace algo más laboriosa la
prueba porque procuramos que la unión disjunta se mantenga ordenada, para así
hacer resaltar que la suma es disjunta. Empezamos con los casos en que
\(n=1\) ó \(m=1\).



\begin{eqnarray}
\nonumber
\ldots&&\text{Continúa de la 2ª parte.}\\ \nonumber
\text{Caso}\quad n = 1 &\quad& \forall B \in \IIIQ \\ 
&\big{|}&\exists k \leq m \;A = \left[a_0,a_1\right] \subseteq \left[b_{2
\cdot k},b_{(2 \cdot k)+1}\right] \\
&\phantom{\big{|}}&\; = B\\
&\big{|}&\exists k \leq m \;A = \left[a_0,a_1\right] \subseteq \left]b_{(2
\cdot k)+1},b_{(2 \cdot k)+2}\right]\\
&\phantom{\big{|}}&\; = {\biguplus}_{\kappa=0}^{k} \left[b_{2 \cdot
\kappa},b_{(2 \cdot \kappa)+1}\right] \uplus \left[a_0,a_1\right] \uplus
{\biguplus}_{\kappa=k+1}^{m} \left[b_{2 \cdot \kappa},b_{(2 \cdot
\kappa)+1}\right]\\ \nonumber
\text{Caso}\quad m = 1 &\quad& \forall A \in \IIIQ \\ 
&\big{|}&\exists k \leq n \;\; B = \left[b_0,b_1\right] \subseteq \left[a_{2
\cdot k},a_{(2 \cdot k)+1}\right] \\
&\phantom{\big{|}}&\; = A\\
&\big{|}&\exists k \leq n \;\;  B = \left[b_0,b_1\right] \subseteq
\left]a_{(2 \cdot k)+1},a_{(2 \cdot k)+2}\right] \\
&\phantom{\big{|}}&\; = {\biguplus}_{\kappa=0}^{k} \left[a_{2 \cdot
\kappa},a_{(2 \cdot \kappa)+1}\right] \uplus \left[b_0,b_1\right] \uplus
{\biguplus}_{\kappa=k+1}^{n} \left[a_{2 \cdot \kappa},a_{(2 \cdot
\kappa)+1}\right]\\ \nonumber
&&\text{3ª parte.}\;\text{Continúa}\ldots
\end{eqnarray}



	Seguimos ahora con el caso \(m \in \setN\setminus\set{1}\). La longitud
de \(A\) sigue siendo \(n=1\), pero ahora vemos el caso en que \(A\) ocupa a
lo sumo un intervalo y un intersticio (un intervalo entre dos intervalos
consecutivos de \(B\)). Utilizaremos la notación \( \left| \left] \cdot ,
\cdot \right| \right[ \) para denotar un intervalo de \(\IIQ\) que \emph{está
excluido} de la unión (aunque esté en una cadena de uniones) pero que es
disjunto y ordenado con el resto de la unión de intervalos incluidos y
excluidos. Además lo consideraremos abierto, como se sugiere en la notación.
\(B_b\) es el comienzo de \(B\) hasta la parte mencionada explícitamente.
\(B_l\) es el final de \(B\) desde lo último mencionado en la expresión
explícitamente.



\begin{eqnarray}
\nonumber
\ldots&&\text{Continúa de la 3ª parte.}\\ \nonumber
\exists k \leq m &\quad&
B = B_b                     \uplus 
    \left[b_{2k},b_{2k+1}\right]       \uplus %% I_1
    B_l\\
&& A = \left[ a_0 , a_1 \right] \\
\text{Caso}\quad&\quad& \left( a_0 \in \left] b_{2k-1} , b_{2k} \right[
\right) \land \left( a_1 \in \left[ b_{2k} , b_{2k+1} \right] \right) \\
A + B & = & B_b \uplus \left[ a_0 , a_1 \right] \uplus B_l \\
%% a_0 in E_1 y a_1 en I_1 => A+B = B_b ≠<U [a_0,b_{2k+1}] ≠<U B_l
\text{Caso}\quad&\quad& \left( a_0 \in \left] b_{2k-1} , b_{2k} \right[
\right) \land \left( a_1 \in \left] b_{2k+1} , b_{2k+2} \right[ \right) \\
A + B & = & B_b \uplus \left[ a_0 , b_{2k+1} \right] \uplus B_l\\
\text{Caso}\quad&\quad& \left( a_0 \in \left[ b_{2k} , b_{2k+1} \right]
\right) \land \left( a_1 \in \left] b_{2k+1} , b_{2k+2} \right[ \right) \\
A + B & = & B_b \uplus \left[ b_{2k} , a_1 \right] \uplus B_l\\ \nonumber
&&\text{4ª Parte.}\;\text{Continúa}\ldots
\end{eqnarray}



	Ahora pasamos al caso general dónde \(B\) tiene \(m > 1\) y \(A\) sigue
siendo con \(n = 1\), aunque ahora no ponemos restricciones sobre la longitud
del intervalo \(A\).



\begin{eqnarray}
\nonumber
&&\ldots\text{Continúa de la 4ª parte.}\\ \nonumber
&&\exists s \leq m\\
&&\exists r \leq s\\
&&\;\;    0 \leq r\\
&&B = B_b                                \uplus 
    \left[b_{2k},b_{2k+1}\right]       \uplus %% I_1
    \uplus \ldots \uplus
    \left[b_{2s},b_{2s+1}\right]       \uplus %% I_2
    B_l\\ \nonumber
&&\text{Dónde }\\
&&B_b \mathdef 
{{\biguplus}_{\kappa=0}^{k-1}}{\left[{b_{2\kappa},b_{2\kappa+1}}\right]}\\
&&B_l \mathdef 
{{\biguplus}_{\kappa=s+1}^{m}}{\left[{b_{2\kappa},b_{2\kappa+1}}\right]}\\
&&A = \left[ a_0 , a_1 \right]\\
&&\text{Tenemos que  }\\
&&B \subseteq B_b \uplus E_{r} \uplus I_{r} \uplus
E_{r+1} \uplus \ldots \uplus E_{s} \uplus I_{s} \uplus E_{s+1} \uplus B_l\\
&&\text{Dónde }\\
&&E_{r} \mathdef \left]b_{2r-1},b_{2r}\right[\\
&&I_{r} \mathdef \left[b_{2r},b_{2r+1}\right]\\
&&E_{r+1} \mathdef \left]b_{2r+1},b_{2r+2}\right[\\
&&E_{s} \mathdef \left]b_{2s-1},b_{2s}\right[\\
&&I_{s} \mathdef \left[b_{2s},b_{2s+1}\right]\\
&&E_{s+1} \mathdef \left]b_{2s+1},b_{2s+2}\right[\\
&&A+B = \left\lbrace \begin{matrix}
B_b \uplus [a_0,a_1] \uplus I_{s} \uplus B_l &\Leftarrow& a_0 \in E_{r-1} \wedge a_1 \in E_{s-1}\\
B_b \uplus [a_0,b_{2s+1}] \uplus B_l &\Leftarrow& a_0 \in E_{r-1} \wedge a_1 \in I_{s}\\
B_b \uplus [a_0,a_1] \uplus B_l &\Leftarrow& a_0 \in E_{r-1} \wedge a_1 \in E_{s+1}\\
B_b \uplus [b_{2r},a_1] \uplus I_{2s} \uplus B_l &\Leftarrow& a_0 \in I_{r}  \wedge a_1 \in E_{s-1}\\
B_b \uplus [b_{2r},b_{2r+1}] \uplus B_l &\Leftarrow& a_0 \in I_{r} \wedge a_1 \in I_{s}\\
B_b \uplus [b_{2s},a_1] \uplus B_l &\Leftarrow& a_0 \in I_{r}   \wedge a_1 \in E_{s+1}\\
B_b \uplus I_{r} \uplus [a_0,a_1] \uplus I_{s} \uplus B_l &\Leftarrow& a_0 \in E_{r+1} \wedge a_1 \in E_{s-1}\\
B_b \uplus I_{r} \uplus [a_0,b_{s+1}] \uplus B_l &\Leftarrow& a_0 \in E_{r+1} \wedge a_1 \in I_{s}\\
B_b \uplus I_{r} \uplus [a_0,a_1] \uplus B_l &\Leftarrow& a_0 \in E_{r+1} \wedge a_1 \in E_{s+1}\\ 
\end{matrix}\right.\\ \nonumber
&&\text{5ª Parte.}\;\text{Continúa}\ldots
\end{eqnarray}



	Ahora ya podemos terminar de definir la suma \( A + B \) por recursión sobre cada uno de los intervalos de \(A\).



\begin{eqnarray}
\nonumber
&&\ldots\text{Continúa de la 5ª parte.}\\ \nonumber
&&\text{Sea  } n \in \setN \;\; \exists \set{a_0,\ldots,a_{2n+1}} \subset \IQ\\
&&\qquad 0 \leqq a_0 \lneqq a_1  \lneqq \ldots \lneqq a_{2n} \lneqq a_{2n+1} \leqq 1\\
&& A = {\biguplus}_{\kappa=0}^{n} I_{\kappa}^A = {\biguplus}_{\kappa=0}^{n} [a_{2\kappa},a_{2\kappa+1}] \vee A = []\\
&&\text{Sea  } m \in \setN \;\; \exists \set{b_0,\ldots,b_{2m+1}} \subset \IQ\\
&&\qquad 0 \leqq b_0 \lneqq b_1 \lneqq \ldots \lneqq b_{2m} \lneqq b_{2m+1} \leqq 1\\
&& B = {\biguplus}_{\nu=0}^{m} I_{\nu}^B = {\biguplus}_{\nu=0}^{m} [b_{2\nu},b_{2\nu+1}] \vee B = []\\
&&A + B \quad=\quad {\biguplus}_{\kappa=0}^{n} I_{\kappa}^A + B =\\
&& \qquad\qquad =\quad {\biguplus}_{\kappa=0}^{n-1} I_{\kappa}^A + \left( I_{n} + B \right) =\\
&&\qquad\qquad=\quad {\biguplus}_{\kappa=0}^{n-2} I_{\kappa}^A + \left( I_{n-1}^A + \left( I_{n}^A + B \right) \right) =\\
&&\qquad\qquad=\quad {\biguplus}_{\kappa=0}^{n-3} I_{\kappa}^A + \left( I_{n-2}^A + \left( I_{n-1}^A + \left( I_{n}^A + B \right) \right) \right) = \ldots =\\
&&\qquad\qquad=\quad I_{0}^A + \left( I_{1}^A + \left( I_{2}^A +  \ldots + \left( I_{n-2}^A + \left( I_{n-1}^A + \left( I_{n}^A + B \right) \right) \right) \right) \right)\\ \nonumber
&&\text{6ª Parte.}\qquad\text{Fin de la definición.}
\end{eqnarray}



	Ahora es ya fácil ver que para el caso completamente general, solo hay
que ir sumando un único intervalo de \(A\) cada vez hasta terminar de sumar
todos los términos (intervalos simples). Por lo tanto es una operación
binaria interna bien definida. Como vemos en la propia definición, tiene
elemento neutro, es conmutativa y asociativa por recaer enteramente la
operación en uniones que son a su vez conmutativas y asociativas.



	Para la multiplicación buleana daremos parecidos pasos, pero será más
sencillo. Primero probaremos que \(1_{\IIIQ} \mathdef \IQ\) hace de elemento
neutro del producto (universo). Al ser el producto una intersección, la
intersección con el universo, que es nuestro uno, dejará siempre todo igual.
Por las mismas razones que en la definición de la suma, la definición del
producto será también conmutativa y asociativa al recaer en la intersección
de conjuntos.



	Como en el caso de la suma, primero supondremos que \(A\) y \(B\) toman
algunos valores particulares. Procederemos por casos:



\begin{eqnarray}
\nonumber
&&\text{Casos de elemento neutro de la suma y del producto}\\
&&A \cdot B : \IIIQ \times \IIIQ \fromto \IIIQ\\
&&      \left[\,\right] \cdot \left[\,\right] = \left[\,\right]\\
&&      \left[\,\right] \cdot B = \left[\,\right]\\
&&      A \cdot \left[\,\right] = \left[\,\right]\\
&&      \IQ \cdot \IQ = \IQ\\
&&      \IQ \cdot B   = B\\
&&      A   \cdot \IQ = A\\ \nonumber
&&      \text{1ª parte.}\;\;\text{Continúa}\ldots
\end{eqnarray}



	Ahora trataremos con los casos en que \(A\) y \(B\) consisten cada uno en
un intervalo, \(n=1\) y \(m=1\).



\begin{eqnarray}
\nonumber
\ldots&&\text{Continúa de la 1ª parte.}\\
\text{Caso}\quad n = 1 \;\; m = 1 &&\\
A \cdot B &:& \IIIQ \times \IIIQ \fromto \IIIQ\\
      &&\left[a_0,a_1\right] \cdot \left[b_0,b_1\right] = \qquad\\
      &|& a_0 \lneqq a_1 \lneqq b_0 \lneqq b_1\qquad = \left[\;\right]\\ 
      &|& a_0 \lneqq b_0 \leqq a_1 \lneqq b_1\qquad = \left[b_0,a_1\right]\\ 
      &|& a_0 \lneqq b_0 \lneqq b_1 \lneqq a_1\qquad = B\\ 
      &|& b_0 \lneqq b_1 \lneqq a_0 \lneqq a_1\qquad = \left[b_1,a_0\right]\\ 
      &|& b_0 \lneqq a_0 \leqq b_1 \lneqq a_1\qquad = \left[a_0,b_1\right]\\ 
      &|& b_0 \lneqq a_0 \lneqq a_1 \lneqq b_1\qquad = A\\ \nonumber
      &&\text{2ª parte.}\;\text{Continúa}\ldots
\end{eqnarray}



	Es importante darse cuenta que el producto de dos intervalos es siempre,
a lo sumo, un solo intervalo. En ocasiones puede dar un caso extraño: cuando
coincidan \(a_1 = b_0\) ó \(a_0 = b_1\) en que el producto
\(\left[a_1,b_0\right]\) ó \(\left[a_0,b_1\right]\) consisten en conjuntos de
un solo punto. Para estos casos hemos introducido el cociente por un conjunto
finito de puntos, de forma que estos casos son de la misma clase de
equivalencia de \(\left[\,\right]\).



	Como consecuencia directa, si \(B={\biguplus}_{\kappa=0}^{m}
\left[b_{2\cdot\kappa},b_{2\cdot\kappa+1}\right]\) y
\(A=\left[a_0,a_1\right]\), \(A\cdot B\) será la unión disjunta y ordenada de
a lo sumo \(m\) intervalos. Investigamos este caso, en los casos especiales
en que \(\exists r \in \setN\) con \(r \leq m\), tal que bien \(A \subseteq
\left[b_{2\cdot r},b_{2\cdot r+1}\right]\) o bien  \(\left[b_{2\cdot
r},b_{2\cdot r+1}\right] \subset A\). También trataremos los casos que la
intersección con \(A\) es vacío, esto es, \(\forall r \in \setN\quad r \leq
m\), \(A \cap \left[b_{2\cdot r},b_{2\cdot r+1}\right] = \emptyset\) o bien
\(A \cap \left[b_{2\cdot r-1},b_{2\cdot r}\right] = \set{b_{2\cdot
r-1},b_{2\cdot r}}\) para algún caso \(r\in\setN\) o incluso, \(A \cap
\left[b_{2\cdot r},b_{2\cdot r+1}\right] = \set{b_{2\cdot r}}\) para algún
caso \(r\in\setN\).



\begin{eqnarray}
\nonumber
&&\ldots\text{Continúa de la 2ª parte.}\\
&&\text{Caso}\quad n = 1 \;\; m \in \setN\\
&&B \mathdef {\biguplus}_{\kappa=0}^{m} {\left[b_{2\cdot \kappa},b_{2\cdot
\kappa + 1}\right]}\\
&&\exists r \in \setN \quad r \leq m\\
\qquad&& B_b \mathdef {\biguplus}_{\kappa=0}^{r-1} {\left[b_{2\cdot
\kappa},b_{2\cdot \kappa + 1}\right]}\\
\qquad&& B_l \mathdef {\biguplus}_{\kappa=r+1}^{m} {\left[b_{2\cdot
\kappa},b_{2\cdot \kappa + 1}\right]}\\
\qquad&& I_r^B \mathdef {\left[b_{2\cdot r},b_{2\cdot r + 1}\right]}\\
&&B = B_b \uplus I_r^B \uplus B_l\\
&&A = \left[a_0,a_1\right]\\
A \cdot B &:& \IIIQ \times \IIIQ \fromto \IIIQ\\
      &&\left[a_0,a_1\right] \cdot B =\\
      &|& b_{2\cdot r-1} \lneqq a_0 \lneqq a_1 \lneqq b_{2\cdot r} \lneqq
      b_{2\cdot r+1}\\
      \qquad&&= \left[\;\right]\\
      &|& b_{2\cdot r-1} = a_0 \lneqq a_1 \lneqq b_{2\cdot r} \lneqq
      b_{2\cdot r+1}\\
      \qquad&&= \set{a_0} {\cong}_{\mathbf{Fin}_{\IQ}} \left[\;\right] \\
      &|& b_{2\cdot r-1} = a_0 \lneqq a_1 = b_{2\cdot r} \lneqq
      b_{2\cdot r+1}\\
      \qquad&&= \set{a_0,a_1} {\cong}_{\mathbf{Fin}_{\IQ}} \left[\;\right] \\
      &|& b_{2\cdot r-1} \lneqq a_0 \lneqq a_1 = b_{2\cdot r} \lneqq
      b_{2\cdot r+1}\\
      \qquad&&= \set{a_1} {\cong}_{\mathbf{Fin}_{\IQ}} \left[\;\right] \\ 
      &|& b_{2\cdot r-1} \lneqq a_0 \leqq b_{2\cdot r} \lneqq a_1 \leqq
      b_{2\cdot r+1}\\
      \qquad&&= \left[b_{2\cdot r},a_1\right]\\
      &|& b_{2\cdot r-1} = a_0 \lneqq b_{2\cdot r} \lneqq a_1 \leqq
      b_{2\cdot r+1}\\
      \qquad&&= \set{a_0} \uplus \left[b_{2\cdot r},a_1\right]
      {\cong}_{\mathbf{Fin}_{\IQ}} \left[b_{2\cdot r},a_1\right] \\
      &|& b_{2\cdot r} \leqq a_0 \lneqq a_1 \leqq
      b_{2\cdot r+1}\\
      \qquad&&= A\\
      &|& b_{2\cdot r} \leqq a_0 \leqq b_{2\cdot r+1} \leqq a_1 \lneqq
      b_{2\cdot r+2}\\
      \qquad&&= \left[b_{2\cdot r+1},a_1\right]\\
      &|& b_{2\cdot r} \leqq a_0 \leqq b_{2\cdot r+1} \leqq a_1 =
      b_{2\cdot r+2}\\
      \qquad&&= \left[b_{2\cdot r+1},a_1\right] \uplus \set{b_{2\cdot
      r+2}}{\cong}_{\mathbf{Fin}_{\IQ}} \left[b_{2\cdot r+1},a_1\right]\\
      &&\text{3ª parte.}\;\text{Continúa}\ldots
\end{eqnarray}



	Ahora miramos el caso en que \(A\mathdef\left[a_0,a_1\right]\) tiene
longitud superior, esto es, existen \(r\) y \(s\), \(r \lneqq s \leqq m\),
tales que \(b_{2\cdot r-1} \leqq a_0 \leqq b_{2\cdot r+1}\) y \(b_{2\cdot
s-1} \leqq a_1 \leqq b_{2\cdot s+1}\).  



\begin{eqnarray}
\nonumber
&&\ldots\text{Continúa de la 3ª parte.}\\
&&\text{Caso}\quad n = 1 \;\; m \in \setN\\
&&B \mathdef {\biguplus}_{\kappa=0}^{m} {\left[b_{2\cdot \kappa},b_{2\cdot
\kappa + 1}\right]}\\
&&\exists s \in \setN \quad s \leq m\\
&&\exists r \in \setN \quad r \lneq r\\
\qquad&& B_b \mathdef {\biguplus}_{\kappa=0}^{r-1} {\left[b_{2\cdot
\kappa},b_{2\cdot \kappa + 1}\right]}\\
\qquad&& I_r^B \mathdef {\left[b_{2\cdot r},b_{2\cdot r + 1}\right]}\\
\qquad&& B_l \mathdef {\biguplus}_{\kappa=s+1}^{m} {\left[b_{2\cdot
\kappa},b_{2\cdot \kappa + 1}\right]}\\
\qquad&& I_s^B \mathdef {\left[b_{2\cdot s},b_{2\cdot s + 1}\right]}\\
\qquad&& B_i \mathdef {\biguplus}_{\kappa=r+1}^{s-1} {\left[b_{2\cdot
\kappa},b_{2\cdot \kappa + 1}\right]}\\
&&B = B_b \uplus I_r^B \uplus B_i \uplus I_s^B \uplus B_l\\
&&A = \left[a_0,a_1\right]\\  \nonumber
&&\text{4ª parte.}\;\text{Continúa}\ldots
\end{eqnarray}



\begin{eqnarray}
A \cdot B &:& \IIIQ \times \IIIQ \fromto \IIIQ\\
&&\left[a_0,a_1\right] \cdot B =\\
&|&{\left(b_{2\cdot r-1} = a_0 \leqq b_{2\cdot
r}\right)}\land{\left(b_{2\cdot s-1} = a_1 \leqq b_{2\cdot
s}\right)}\\
\qquad&&= \set{a_0} \uplus I_r^B \uplus B_i \uplus \set{a_1}
{\cong}_{\mathbf{Fin}_{\IQ}} I_r^B \uplus B_i\\
&|&{\left(b_{2\cdot r-1} \lneqq a_0 \leqq b_{2\cdot
r}\right)}\land{\left(b_{2\cdot s-1} = a_1 \leqq b_{2\cdot
s}\right)}\\
\qquad&&= I_r^B \uplus B_i \uplus \set{a_1}
{\cong}_{\mathbf{Fin}_{\IQ}} I_r^B \uplus B_i\\
&|&{\left(b_{2\cdot r-1} \lneqq a_0 \leqq b_{2\cdot
r}\right)}\land{\left(b_{2\cdot s-1} \lneqq a_1 \leqq b_{2\cdot
s}\right)}\\
\qquad&&= I_r^B \uplus B_i\\
&|&{\left(b_{2\cdot r-1} = a_0 \leqq b_{2\cdot
r}\right)}\land{\left(b_{2\cdot s-1} \lneqq a_1 \leqq b_{2\cdot
s}\right)}\\
\qquad&&= \set{a_0} \uplus I_r^B \uplus B_i {\cong}_{\mathbf{Fin}_{\IQ}}
I_r^B \uplus B_i \\  \nonumber
&&\text{4ª parte.}\;\text{Continúa}\ldots
\end{eqnarray}



\begin{eqnarray}
A \cdot B &:& \IIIQ \times \IIIQ \fromto \IIIQ\\
&&\left[a_0,a_1\right] \cdot B =\\
&|&{\left(b_{2\cdot r} \leqq a_0 \leqq b_{2\cdot
r + 1}\right)}\land{\left(b_{2\cdot s-1} \lneqq a_1 \leqq b_{2\cdot
s}\right)}\\
\qquad&&= \left[a_0,b_{2\cdot r + 1}\right] \uplus B_i
{\cong}_{\mathbf{Fin}_{\IQ}}
\left[a_0,b_{2\cdot r + 1}\right] \uplus B_i\\
&|&{\left(b_{2\cdot r} \leqq a_0 \leqq b_{2\cdot
r + 1}\right)}\land{\left(b_{2\cdot s-1} = a_1 \leqq b_{2\cdot
s}\right)}\\
\qquad&&= \left[a_0,b_{2\cdot r + 1}\right] \uplus B_i \uplus \set{b_{2\cdot
s-1}} {\cong}_{\mathbf{Fin}_{\IQ}} \left[a_0,b_{2\cdot r + 1}\right] \uplus
B_i\\
&|&{\left(b_{2\cdot r + 1} = a_0 \lneqq b_{2\cdot
r + 2}\right)}\land{\left(b_{2\cdot s-1} = a_1 \leqq b_{2\cdot
s}\right)}\\
\qquad&&= \set{a_0} \uplus B_i \uplus \set{a_1} {\cong}_{\mathbf{Fin}_{\IQ}}
B_i \\ \nonumber
&&\text{4ª parte.}\;\text{Continúa}\ldots
\end{eqnarray}



\begin{eqnarray}
A \cdot B &:& \IIIQ \times \IIIQ \fromto \IIIQ\\
&&\left[a_0,a_1\right] \cdot B =\\
&|&{\left(b_{2\cdot r + 1} \lneqq a_0 \lneqq b_{2\cdot
r + 2}\right)}\land{\left(b_{2\cdot s-1} = a_1 \leqq b_{2\cdot
s}\right)}\\
\qquad&&= B_i \uplus \set{a_1} {\cong}_{\mathbf{Fin}_{\IQ}}
B_i\\
&|&{\left(b_{2\cdot r + 1} \lneqq a_0 \lneqq b_{2\cdot
r + 2}\right)}\land{\left(b_{2\cdot s-1} \lneqq a_1 \leqq b_{2\cdot
s}\right)}\\
\qquad&&= B_i\\
&|&{\left(b_{2\cdot r + 1} = a_0 \lneqq b_{2\cdot
r + 2}\right)}\land{\left(b_{2\cdot s-1} \lneqq a_1 \leqq b_{2\cdot
s}\right)}\\
\qquad&&= \set{a_0} \uplus B_i {\cong}_{\mathbf{Fin}_{\IQ}} B_i\\
&|&{\left(b_{2\cdot r-1} = a_0 \leqq b_{2\cdot
r}\right)}\land{\left(b_{2\cdot s} \leqq a_1 \leqq b_{2\cdot
s + 1}\right)}\\
\qquad&&= \set{a_0} \uplus I_r^B \uplus B_i \uplus \left[b_{2\cdot 
s},a_1\right]
{\cong}_{\mathbf{Fin}_{\IQ}} I_r^B \uplus B_i \uplus \left[b_{2\cdot 
s},a_1\right]\\ \nonumber
&&\text{4ª parte.}\;\text{Continúa}\ldots
\end{eqnarray}



\begin{eqnarray}
A \cdot B &:& \IIIQ \times \IIIQ \fromto \IIIQ\\
&&\left[a_0,a_1\right] \cdot B =\\
&|&{\left(b_{2\cdot r-1} \lneqq a_0 \leqq b_{2\cdot
r}\right)}\land{\left(b_{2\cdot s} \leqq a_1 \leqq b_{2\cdot
s + 1}\right)}\\
\qquad&&= I_r^B \uplus B_i \uplus \left[b_{2\cdot 
s},a_1\right]\\
&|&{\left(b_{2\cdot r} \lneqq a_0 \lneqq b_{2\cdot
r + 1}\right)}\land{\left(b_{2\cdot s} \leqq a_1 \leqq b_{2\cdot
s + 1}\right)}\\
\qquad&&= \left[a_0,b_{{2\cdot r}+1}\right] \uplus B_i \uplus \left[b_{2\cdot 
s},a_1\right]\\
&|&{\left(b_{2\cdot r} \lneqq a_0 = b_{2\cdot
r + 1}\right)}\land{\left(b_{2\cdot s} \leqq a_1 \leqq b_{2\cdot
s + 1}\right)}\\
\qquad&&= \set{b_{{2\cdot r}+1}} \uplus B_i \uplus \left[b_{2\cdot 
s},a_1\right] {\cong}_{\mathbf{Fin}_{\IQ}} B_i \uplus \left[b_{2\cdot 
s},a_1\right]\\ \nonumber
&&\text{4ª parte.}\;\text{Continúa}\ldots
\end{eqnarray}



\begin{eqnarray}
A \cdot B &:& \IIIQ \times \IIIQ \fromto \IIIQ\\
&&\left[a_0,a_1\right] \cdot B =\\
&|&{ \left( b_{{2 \cdot r} + 1} = a_0 \leqq b_{ {2 \cdot
r} + 2} \right) } \land {\left( b_{2\cdot s} \lneqq a_1 \leqq b_{ 2 \cdot
s + 1} \right) } \\
\qquad&&= \set{a_0} \uplus B_i \uplus \left[b_{2\cdot 
s},a_1\right] {\cong}_{\mathbf{Fin}_{\IQ}} B_i \uplus \left[b_{2\cdot 
s},a_1\right]\\
&|&{ \left( b_{{2 \cdot r} + 1} \lneqq a_0 \leqq b_{ {2 \cdot
r} + 2} \right) } \land {\left( b_{2\cdot s} \lneqq a_1 \leqq b_{ 2 \cdot
s + 1} \right) } \\
\qquad&&= B_i \uplus \left[b_{2\cdot s},a_1\right]\\
&|&{ \left( b_{{2 \cdot r} + 1} = a_0 \leqq b_{ {2 \cdot
r} + 2} \right) } \land {\left( b_{2\cdot s} = a_1 \leqq b_{ 2 \cdot
s + 1} \right) } \\
\qquad&&= \set{a_0} \uplus B_i \uplus \set{a_1} {\cong}_{\mathbf{Fin}_{\IQ}}
B_i\\
&|&{ \left( b_{{2 \cdot r} + 1} \lneqq a_0 = b_{ {2 \cdot
r} + 2} \right) } \land {\left( b_{2\cdot s} = a_1 \leqq b_{ 2 \cdot
s + 1} \right) } \\
\qquad&&= B_i \uplus \set{a_1} {\cong}_{\mathbf{Fin}_{\IQ}}
B_i\\ \nonumber
&&\text{4ª parte.}\;\text{Continúa}\ldots
\end{eqnarray}



\begin{eqnarray}
A \cdot B &:& \IIIQ \times \IIIQ \fromto \IIIQ\\
&&\left[a_0,a_1\right] \cdot B =\\
&|&{ \left( b_{{2 \cdot r} - 1} = a_0 \leqq b_{ {2 \cdot
r} } \right) } \land {\left( b_{{2\cdot s} + 1} = a_1 \leqq b_{ {2
\cdot s} + 2 } \right) } \\
\qquad&&= \set{a_0} \uplus I_r^B \uplus B_i \uplus I_s^B \uplus \set{a_1}
{\cong}_{\mathbf{Fin}_{\IQ}} I_r^B \uplus B_i \uplus I_s^B\\
&|&{ \left( b_{{2 \cdot r} - 1} \lneqq a_0 \leqq b_{ {2 \cdot
r} } \right) } \land {\left( b_{{2\cdot s} + 1} = a_1 \leqq b_{ {2
\cdot s} + 2 } \right) } \\
\qquad&&= I_r^B \uplus B_i \uplus I_s^B \uplus \set{a_1} 
{\cong}_{\mathbf{Fin}_{\IQ}} I_r^B \uplus B_i \uplus I_s^B\\
&|&{ \left( b_{{2 \cdot r} - 1} = a_0 \leqq b_{ {2 \cdot
r} } \right) } \land {\left( b_{{2\cdot s} + 1} \lneqq a_1 \leqq b_{ {2
\cdot s} + 2 } \right) } \\
\qquad&&= \set{a_0} \uplus I_r^B \uplus B_i \uplus I_s^B
{\cong}_{\mathbf{Fin}_{\IQ}}
I_r^B \uplus B_i \uplus I_s^B\\
&|&{ \left( b_{{2 \cdot r} - 1} \lneqq a_0 \leqq b_{ {2 \cdot
r} } \right) } \land {\left( b_{{2\cdot s} + 1} = a_1 \lneqq b_{ {2
\cdot s} + 2 } \right) } \\
\qquad&&= I_r^B \uplus B_i \uplus I_s^B\\ \nonumber
&&\text{4ª parte.}\;\text{Continúa}\ldots
\end{eqnarray}




\begin{eqnarray}
A \cdot B &:& \IIIQ \times \IIIQ \fromto \IIIQ\\
&&\left[a_0,a_1\right] \cdot B =\\
&|&{ \left( b_{{2 \cdot r}} = a_0 \leqq b_{ {2 \cdot
r} + 1 } \right) } \land {\left( b_{{2\cdot s} + 1} = a_1 \leqq b_{ {2
\cdot s} + 2 } \right) } \\
\qquad&&= I_r^B \uplus B_i \uplus I_s^B \uplus \set{a_1} 
{\cong}_{\mathbf{Fin}_{\IQ}} I_r^B \uplus B_i \uplus I_s^B\\
&|&{ \left( b_{ {2 \cdot r} } \lneqq a_0 \leqq b_{ {2 \cdot
r} + 1 } \right) } \land {\left( b_{{2\cdot s} + 1} \leqq a_1 \leqq b_{ {2
\cdot s} + 2 } \right) } \\
\qquad&&= \left[a_0,b_{{2 \cdot r} + 1}\right] \uplus B_i \uplus I_s^B\\
&|&{ \left( b_{ { 2 \cdot r } } = a_0 \leqq b_{ { 2 \cdot
r } + 1 } \right) } \land {\left( b_{{2\cdot s} + 1} \lneqq a_1 \leqq b_{ {2
\cdot s} + 2 } \right) } \\
\qquad&&= I_r^B \uplus B_i \uplus \left[ {a_1,b_{ { 2 \cdot s } + 1 } }
\right] \\
&|&{ \left( b_{ {2 \cdot r} } \lneqq a_0 \leqq b_{ { 2 \cdot
r } + 1 } \right) } \land {\left( b_{{2\cdot s} + 1} \lneqq a_1 \leqq b_{ {2
\cdot s} + 2 } \right) } \\
\qquad&&= \left[ {a_0 , b_{ { 2 \cdot r } + 1 } } \right] \uplus B_i \uplus
\left[ {a_1,b_{ { 2 \cdot s } + 1 } } \right] \\ \nonumber
&&\text{4ª parte.}\;\text{Continúa}\ldots
\end{eqnarray}




\begin{eqnarray}
A \cdot B &:& \IIIQ \times \IIIQ \fromto \IIIQ\\
&&\left[a_0,a_1\right] \cdot B =\\
&|&{ \left( b_{ {2 \cdot r} + 1 } = a_0 \leqq b_{ {2 \cdot
r } + 2 } \right) } \land { \left( b_{ { 2 \cdot s } + 1 } \leqq a_1 \leqq
b_{ { 2 \cdot s } + 2 } \right) } \\
\qquad&&= \set{a_0} \uplus B_i \uplus I_s^B \uplus \set{a_1} 
{\cong}_{\mathbf{Fin}_{\IQ}} B_i \uplus I_s^B\\
&|&{ \left( b_{ { 2 \cdot r } + 1 } \lneqq a_0 \leqq b_{ {2 \cdot
r } + 2 } \right) } \land { \left( b_{ { 2 \cdot s } + 1 } \leqq a_1 \leqq
b_{ {2 \cdot s} + 2 } \right) } \\
\qquad&&= B_i \uplus I_s^B \\ \nonumber
&&\text{4ª parte.}\;\text{Continúa}\ldots
\end{eqnarray}



	En la definición del producto como intersección de uniones disjuntas (y
ordenadas) de intervalos de \(\IQ\) ha quedado clara la necesidad de
convertir los intervalos a sus clases de equivalencia módulo un conjunto
finito de puntos de \(\IQ\). A estos intervalos los llamaremos intervalos
genéricos, ya que no requieren de la especificación de cerrados o abiertos o
semicerrados por un extremo u otro. Mejor nombre que \(\IVQ\) sería
\({\mathbf{III}}_{\setQ}^{\mathbf{Gen}}\).



Aún nos queda definir el caso general, por recursión sobre los intervalos
genéricos de \(A\). Procedemos como en el caso de la suma.



\begin{eqnarray} \nonumber
&&\ldots\text{Continúa de la 4ª parte.}\\
&&\exists n \in \setN \qquad
\exists \set{a_0,a_1,\ldots,a_{2\cdot n} , a_{{2 \cdot n} + 1}}
\subsetneqq \IQ \\
&&\text{       tal que  }\quad { 0 \leqq a_0 \lneqq a_1 \lneqq \ldots \lneqq
a_{2\cdot n} \lneqq a_{{2\cdot n} + 1} \leqq 1 }\\
&& A \mathdef {\biguplus}_{\kappa = 0}^{n} \left[a_{2\cdot\kappa},a_{{2\cdot
\kappa}+1}\right]\\
&&\exists m \in \setN \qquad
\exists \set{b_0,b_1,\ldots,b_{2\cdot m} , b_{{2 \cdot m} + 1}}
\subsetneqq \IQ \\
&&\text{       tal que  }\quad { 0 \leqq b_0 \lneqq b_1 \lneqq \ldots \lneqq
b_{2\cdot m} \lneqq b_{{2\cdot m} + 1} \leqq 1 }\\
&& B \mathdef {\biguplus}_{\kappa = 0}^{m} \left[b_{2\cdot\kappa},b_{{2\cdot
\kappa}+1}\right]\\ \nonumber
&& \textsf{		Definimos	}\\
&& A \cdot B = \left( { \biguplus }_{ \kappa = 0 }^{ n } \left[ a_{ 2 \cdot
\kappa } , a_{ { 2 \cdot \kappa } + 1 } \right] \right) \cdot B \mathdef \\
&&\qquad \mathdef \left( { \biguplus }_{ \kappa = 0 }^{ n-1 } \left[ a_{ 2
\cdot \kappa } , a_{ { 2 \cdot \kappa } + 1 } \right] \right) \cdot \left(
\left[ a_{ 2 \cdot n } , a_{ { 2 \cdot n } + 1 } \right] \cdot B \right) \\
\nonumber
&&\textsf{5ª Parte. Fin de la definición de producto buleano.}
\end{eqnarray}



	Nos queda solo ver como definimos el complementario de una unión finita
de intervalos disjuntos (y ordenados). Esto podremos hacerlo con una definición sencilla:



\begin{eqnarray}\nonumber
&&\textsf{Definición de complementario de un intervalo genérico.}\\
&&\exists n \in \setN \qquad \exists \set{a_0,a_1,\ldots,a_{2 \cdot n},a_{ \cdot n} + 1} \subsetneqq \IQ\\
&&\textsf{ tal que } \qquad 0 \leqq a_0 \lneqq a_1 \lneqq \ldots \lneqq a_{2 \cdot n} \lneqq a_{{2 \cdot n} + 1} \leqq 1\\
&&\forall k \in \setN \qquad k \leqq n \then I_k^A \mathdef \left[a_{2 \cdot k},a_{{2 \cdot k} + 1}\right]\\
&& A \mathdef {\biguplus}_{\kappa = 0}^{n} I_{\kappa}^A\\ \nonumber
&&\textsf{Ahora definimos los intervalos en los intersticios }\\
&&\forall k \in \setN \qquad k \leqq n \then E_k \mathdef \left]a_{2 \cdot k - 1},a_{2 \cdot k}\right[\\
&& \neg A \mathdef \left[0,a_0\right[\; \uplus\; {\biguplus}_{\kappa = 0}^{n} E_{\kappa}\; \uplus \; \left]a_{2 \cdot n + 1},1\right]
\end{eqnarray}



	Ahora queda claro que el complementario de una unión finita disjunta de
intervalos genéricos en \(\IQ\) es de nuevo una unión finita disjunta de
intervalos genéricos. Cumple claramente con los postulados de Huntington: \(A
+ \neg A {\cong}_{{\mathsf{Fin}}_{\IQ}} \IQ\) y \(A \cdot \neg A
{\cong}_{{\mathsf{Fin}}_{\IQ}} \left[\;\right]\).



	Lo único que queda por ver es la propiedad distributiva de la suma sobre
el producto y del producto sobre la suma. Estas propiedades se cumplen al ser
heredadas de la distribución del unión sobre la intersección y de la
intersección sobre la unión.





	El postulado seis de Huntington se cumple desde que
\(\ClEqFinDiffInter{\IQ} \neq \ClEqFinDiffInter{\;}\).



	Ahora bien, es importante darse cuenta que la cardinalidad de este
álgebra de Boole \( {\mathscr{B}}_{\mathsf{I_\setQ}^\mathsf{Gen}} = \langle
\setB = \IVQ , + \simeq \cup , \cdot \simeq \cap \rangle \) es como el de los
números naturales \( \natural\setN = {\aleph}_0 \), pues cada elemento viene
representado totalmente por un conjunto finito de números racionales entre el
\(0\) y el \(1\). En los próximos capítulos veremos que esto constituye una
excepcionalidad de este álgebra de Boole. Por lo pronto vemos que no puede
ser un álgebra de Boole de un álgebra de las potencias de un conjunto
(prácticamente todas serán el álgebra de las potencias de un conjunto).
Además no tendremos elementos atómicos desde los que construir los demás
elementos del álgebra. Tampoco podremos hacer uso de uniones o intersecciones
infinitas, que si están bien definidas en casi todas las álgebras de Boole
que veremos (y sean infinitas).




\chapter{Teoremas de las álgebras de Boole}



\section{Unicidad del complementario y doble complemento.}



\begin{thm}[Unicidad del complementario]\label{thm:TUC}
Dado \(x \in \setB \), existe \(\widehat{x}\subseteq\setB\) tal que \(
\widehat{x} \neq \emptyset \). Sean \(y\in\widehat{x}\) y
\(z\in\widehat{x}\), entonces \(y = z\).  
\begin{equation}
\forall x \qquad x \in \setB \then \left( \set{y,z} \subseteq \widehat{x}
\then y = z \right)
\end{equation}
\end{thm}



\begin{formalproof}
  \begin{flalign}
  x \in \setB \quad &\then \quad \exists \set{y,z} \subseteq \widehat{x}\\
  \\ \nonumber
  y &= y \cdot 1 =\\
  {} &= y \cdot ( x + z ) =\\
  {} &= ( y \cdot x ) + ( y \cdot z ) =\\
  {} &= 0 + ( y \cdot z ) =\\
  {} &= ( y \cdot z ) + 0 =\\
  {} &= y \cdot z\\
  \\ \nonumber
  z &= z \cdot 1 =\\
  {} &= z \cdot ( x + y ) =\\
  {} &= ( x \cdot x ) + ( z \cdot y ) =\\
  {} &= 0 + ( z \cdot y ) =\\
  {} &= ( z \cdot y ) + 0 =\\
  {} &= z \cdot y =\\
  {} &= y \cdot z\\
  \\ \nonumber
  y &= y \cdot z =\\
  {}&= z   
  \end{flalign}
\end{formalproof}



\begin{defn}[Elemento complementario]\label{def:ElemCompl}
Para cualquier \(x\) elemento del álgebra de Boole hay definido un único
elemento complementario perteneciente a \(\widehat{x}\) \(=\) \(\left\lbrace y \in \setB \right.\) \(\;\big\vert\;\) \( \left. \left( x + y = 1
\right) \right. \) \(\ly\) \( \left. \left( x \cdot y = 0 \right) \right\rbrace \). Denotaremos a ese único elemento complementario por \(
\overline{x} \).
\begin{equation}
\forall x \in \setB \quad \exists ! y \in \widehat{x} \quad \overline {x} \mathdef y
\end{equation}
\end{defn}


\begin{thm}[Doble complementación es identidad]\label{thm:TDCID}
Dado \(x \in \setB \), entonces \(\overline{\overline{x}}=x\).  
\begin{equation}
\forall x \qquad x \in \setB \then x = \overline{\overline{x}}
\end{equation}
\end{thm}



\begin{formalproof}
  \begin{flalign}
  x \in \setB \quad &{}\\
  \\ \nonumber
  1 &= x + \overline{x}\\
  0 &= x \cdot \overline{x}\\
  \\ \nonumber
  1 &= \overline{x} + x\\
  0 &= \overline{x} \cdot x\\
  \\ \nonumber
  x &\in \widehat{\overline{x}}\\
  \\ \nonumber
  1 &= \overline{x} + \overline{\overline{x}}\\
  0 &= \overline{x} \cdot \overline{\overline{x}}\\
  \\ \nonumber
  \overline{\overline{x}} &\in \widehat{\overline{x}}\\
  \\ \nonumber
  x &= \overline{\overline{x}}
  \end{flalign}
\end{formalproof}


\section{Unicidad de elementos neutros.}


\begin{thm}[El elementos neutro de la suma es único]\label{thm:TUENS}
Si un elemento del álgebra de Boole \(e\) es tal que para todo elemento del álgebra \(x\) resulta \(x + e = x\), entonces \(e = 0\).
\begin{equation}
e \in \setB \quad \left( \forall x \quad \left( x \in \setB \; \ly \; x + e = x \right) \right) \Longrightarrow \quad e = 0
\end{equation}
\end{thm}



\begin{formalproof}
  \begin{flalign}
  \forall x \in \setB \qquad \quad & x + e = x\\
  0 \in \setB \qquad \quad & 0 + e = 0\\
  {} & e + 0 = 0 \\
  {} & e + 0 = e \\
  {} & e = 0
  \end{flalign}
\end{formalproof}


\begin{thm}[El elemento neutro del producto es único]\label{thm:TUENP}
Si un elemento del álgebra de Boole \(u\) es tal que para todo elemento del
álgebra \(x\) resulta \(x \cdot u = x\), entonces \(u = 1\).
\begin{equation}
u \in \setB \quad \left( \forall x \quad \left( x \in \setB \; \ly \; x \cdot
u = x \right) \right) \Longrightarrow \quad u = 1
\end{equation}
\end{thm}



\begin{formalproof}
  \begin{flalign}
  \forall x \in \setB \qquad \quad & x \cdot u = x\\
  1 \in \setB \qquad \quad & 1 \cdot u = 1\\
  {} & u \cdot 1 = 1 \\
  {} & u \cdot 1 = u \\
  {} & u = 1
  \end{flalign}
\end{formalproof}



\section{Idempotencia de las operaciones binarias internas.}



\begin{thm}[Idempotencia de la suma]\label{thm:TIS}
La suma de todo elemento del álgebra de Boole consigo mismo tiene como imagen
el propio elemento. En ecuación, \( x + x = x\) para cualquier \(x \in
\setB\).
\begin{equation}
\forall x \quad \quad x \in \setB \quad \Longrightarrow \quad x + x = x
\end{equation}
\end{thm}



\begin{formalproof}
  \begin{flalign}
  \forall x \in \setB \qquad & \; x = x + 0 = \\
  {} \qquad & \quad  = x + ( x \cdot \overline{x} ) = \\
  {} & \quad  = ( x + x ) \cdot ( x + \overline{x} ) = \\
  {} & \quad  = ( x + x ) \cdot 1 = \\
  {} & \quad  = x + x
  \end{flalign}
\end{formalproof}



\begin{thm}[Idempotencia del producto]\label{thm:TIP}
El producto de todo elemento del álgebra de Boole consigo mismo tiene como
imagen el propio elemento. En ecuación, \( x \cdot x = x\) para cualquier \(x
\in \setB\).
\begin{equation}
\forall x \quad \quad x \in \setB \quad \Longrightarrow \quad x \cdot x = x
\end{equation}
\end{thm}



\begin{formalproof}
  \begin{flalign}
  \forall x \in \setB \qquad\; & \; x = x \cdot 1 = \\
  {} \qquad & \quad  = x \cdot ( x + \overline{x} ) = \\
  {} & \quad  = ( x \cdot x ) + ( x \cdot \overline{x} ) = \\
  {} & \quad  = ( x \cdot x ) + 0 = \\
  {} & \quad  = x \cdot x
  \end{flalign}
\end{formalproof}



\section{Elementos absorbentes de las operaciones internas.}



\begin{thm}[Elemento absorbente de la suma]\label{thm:TEAS}
La suma de todo elemento del álgebra de Boole con el neutro del producto \(1\) es \(1\). En ecuación, \( x + 1 = 1\) para cualquier \(x \in
\setB\).
\begin{equation}
\forall x \quad \quad x \in \setB \quad \then \quad x + 1 = 1
\end{equation}
\end{thm}



\begin{formalproof}
  \begin{flalign}
  \forall x \in \setB \quad & \quad 1 = x + \overline{x}
  = \\
  {} & \quad  \;\;= x + ( \overline{x} \cdot 1 ) = \\
  {} & \quad  \;\;= ( x + \overline{x} ) \cdot ( x + 1 ) = \\
  {} & \quad  \;\;= 1 \cdot ( x + 1 ) = \\
  {} & \quad  \;\;= ( x + 1 ) \cdot 1 = \\
  {} & \quad  \;\;= x + 1
  \end{flalign}
\end{formalproof}



\begin{thm}[Elemento absorbente del producto]\label{thm:TEAP}
El producto de todo elemento del álgebra de Boole con el neutro de la suma
\(0\) es \(0\). En ecuación, \( x \cdot 0 = 0\) para cualquier \(x \in
\setB\).
\begin{equation}
\forall x \quad \quad x \in \setB \quad \Longrightarrow \quad x \cdot 0 = 0
\end{equation}
\end{thm}



\begin{formalproof}
  \begin{flalign}
  \forall x \in \setB \quad & \quad 0 = x \cdot \overline{x}  = \\
  {} & \quad  \;\;= x \cdot ( \overline{x} + 0 ) = \\
  {} & \quad  \;\;= ( x \cdot \overline{x} ) + ( x \cdot 0 ) = \\
  {} & \quad  \;\;= 0 + ( x \cdot 0 ) = \\
  {} & \quad  \;\;= ( x \cdot 0 ) + 0 = \\
  {} & \quad  \;\;= x \cdot 0
  \end{flalign}
\end{formalproof}



\section{Álgebras de Boole triviales. Condiciones.}


\begin{thm} [Si el álgebra de Boole no es trivial entonces el neutro de la
suma es distinto del neutro del producto] \label{thm:If0Eq1ABT}
Si \( 0 = 1 \), entonces \( \forall x \in \setB \then \setB=\set{x} \).
\begin{equation}
0 = 1 \quad \then \quad \forall x \quad \left( \quad x \in \setB \quad \then \quad \setB = \set{x} \quad \right)
\end{equation}
\end{thm}



\begin{formalproof}
  \begin{flalign}
  \textsf{Hipótesis}: \quad 0 = 1&{}\\
  {\forall x \in \setB} & \quad  \;\; x = x \cdot 1 =\\
  {} & \quad  \;\;   = x \cdot \left( x + \overline{x} \right) = \\
  {} & \quad  \;\;   = \left( x \cdot x \right) + \left( x \cdot \overline{x}
  \right) = \\
  {} & \quad  \;\;= x + 0 = \\
  \textsf{Hipótesis} & \quad  \;\;= x + 1 = \\
  \textsf{Hipótesis} & \quad  \;\;= 1 = 0 \\
  0 = 1 \then& \left(\; \forall x \; \left(\; x \in \setB \then x = 0 = 1 \; \right)\; \right)
  \end{flalign}
\end{formalproof}



\begin{thm} [Si un álgebra de Boole no es trivial entonces todo elemento es
distinto de su complementario] \label{thm:ABTIfxEqnx}
Si \( \exists x \in \setB \) tal que \( x = \overline{x}\), entonces \(
\forall x\) \( x \in \setB \then \setB = \set{x} \).
\begin{equation}
\exists x \qquad x \in \setB \ly x = \overline{x} \quad \then  \quad \forall
y \quad \left( \quad y \in \setB \quad \then \quad \setB = \set{y} \quad
\right)
\end{equation}
\end{thm}



\begin{formalproof}
  \begin{flalign}
  \textsf{Hipótesis}: \quad \exists x \in \setB \quad x = \overline{x}&{}\\
  \\ \nonumber
  {} & \quad  \;\; x = x \cdot 1 =\\
  {} & \quad  \;\;   = x \cdot \left( x + \overline{x} \right) = \\
  {} & \quad  \;\;   = \left( x \cdot x \right) + \left( x \cdot \overline{x}
  \right) = \\
  \textsf{Hipótesis} & \quad  \;\;   = \left( x \cdot \overline{x} \right) +
  \left( x \cdot \overline{x} \right) = \\
  {} & \quad  \;\;= 0 + 0 = 0\\
  \\ \nonumber
  \textsf{Hipótesis} & \quad  \;\;    x = \overline{x}\\
  {} & \quad  \;\;    x = 0\\
  {} & \quad  \;\; x \;\; = \;\; \overline{0} \;\; =   \;\; 1\\
  \\ \nonumber
  {} & \quad  \;\; 0 \;\; = \;\; x \;\; = \;\; \overline{x} \;\; = \;\; 1 \\
  \\ \nonumber
  \exists x \in \setB \quad x = \overline{x} &\then \;\; \forall y \;\;
  \left( y \in \setB \then y = 0 = 1 \right)\\ 
  {}&\;\; por \;\; \ref{thm:If0Eq1ABT} 
  \end{flalign}
\end{formalproof}


\begin{defn}
Llamamos a un álgebra de Boole \(\mathscr{B}\) trivial si su conjunto base
\(\setB\) tiene elementos neutros de la suma \(0\) y del producto \(1\)
diferentes, o, dicho de otra manera, si cada elemento del conjunto \(\forall
x\) \(x\in\setB\) cumple \(x \neq \overline{x}\). Dicho de otra forma
\(\natural \setB > 1\).
\begin{equation}
\mathsf{Triv} \mathscr{B} \mathdef \natural\setB = 1
\end{equation}
\begin{equation}
\mathsf{NoTriv} \mathscr{B} \mathdef \natural\setB \geq 2
\end{equation}



\paragraph*{}
Con base en los teoremas anteriores tenemos algunas caracterizaciones
completas de las álgebras de Boole triviales y no triviales.



\begin{equation}
\mathsf{Triv} \mathscr{B} \; \iff \; 0 = 1 \; \iff \; \exists x \in
\setB \; x = \overline{x} \; \iff \; \natural\setB = 1
\end{equation}
\begin{equation}
\mathsf{NoTriv} \mathscr{B} \; \iff \; 0 \neq 1 \; \iff \; \forall x \in
\setB \; x \neq \overline{x} \; \iff \; \natural\setB \geq 2
\end{equation}
\end{defn}



\section{Teoremas sobre los elementos neutros.}



\begin{thm}[Complementos de los neutros]\label{thm:compl01}
Si \(\mathscr{B}\) es un álgebra de Boole, \(\overline{0} = 1\) y
\(\overline{1} = 0\). 
\end{thm}



\begin{proof}
\begin{flalign}
\text{por H2s} \quad& 1 + 0 = 1\\
\text{por elemento absorbente} \quad& 1 \cdot 0 = 0\\
\text{por unicidad del complementario} \quad& \overline{1} = 0\\
\text{por unicidad del complementario} \quad& \overline{\overline{1}} =
\overline{0}\\
\text{por doble negación es identidad} \quad& 1 = \overline{0}
\end{flalign}
\end{proof}



\begin{thm}[Suma de los neutros]
Si \(\mathscr{B}\) es un álgebra de Boole, \(0 + 0 = 0, 0 + 1 = 1, 1 + 0 = 1,
1 + 1 = 1\). 
\end{thm}



\begin{proof}
\begin{flalign}
\text{por H2s}\quad & 0 + 0 = 0\\
\text{por H2s}\quad & 1 + 0 = 1\\
\text{por H3s}\quad & 0 + 1 = 1 + 0 = 1\\
\text{por elemento absorbente de la suma}\quad & 1 + 1 = 1
\end{flalign}
\end{proof}



\begin{thm}[Producto de los neutros]
Si \(\mathscr{B}\) es un álgebra de Boole, \(0 \cdot 0 = 0, 0 \cdot 1 = 0, 1
\cdot 0 = 0, 1 \cdot 1 = 1\). 
\end{thm}



\begin{proof}
\begin{flalign}
\text{por H2p}\quad & 0 \cdot 1 = 0\\
\text{por H2p}\quad & 1 \cdot 1 = 1\\
\text{por H3p}\quad & 1 \cdot 0 = 0 \cdot 1 = 0\\
\text{por elemento absorbente del producto}\quad & 0 \cdot 0 = 0
\end{flalign}
\end{proof}



\section{Simplificaciones.}


\begin{thm}[Simplificación por absorción: suma]\label{thm:TSimplPorAbsS}
Dados dos elementos del conjunto base \(\setB\) de un álgebra de Boole \(
\mathscr{B} \), \( x \) e \( y \), \( x + ( x \cdot y ) = x \).
\end{thm}



\begin{proof}
  \begin{flalign}
\forall x \; x \in \setB \;\;& \forall y \; y \in \setB {}&{}\\
x + \left( x \cdot y \right) &\;=\; \left( {x \cdot 1} \right) +  \left( {x
\cdot y} \right) \;=\\
{} &\;=\; x \cdot \left( 1 + y \right) \;=\\
{} &\;=\; x \cdot \left( y + 1 \right) \;=\\
{} &\;=\; x \cdot 1 \;=\\
{} &\;=\; x \\
  \end{flalign}
\end{proof}



\begin{thm}[Simplificación por absorción: producto]\label{thm:TSimplPorAbsP}
Dados dos elementos del conjunto base \(\setB\) de un álgebra de Boole \(
\mathscr{B} \), \( x \) e \( y \), \( x \cdot ( x + y ) = x \).
\end{thm}



\begin{proof}
  \begin{flalign}
\forall x \; x \in \setB \;\;& \forall y \; y \in \setB {}&{}\\
x \cdot \left( x + y \right) &\;=\; \left( {x + 0} \right) \cdot  \left( {x
+ y} \right) \;=\\
{} &\;=\; x + \left( 0 \cdot y \right) \;=\\
{} &\;=\; x + \left( y \cdot 0 \right) \;=\\
{} &\;=\; x + 0 \;=\\
{} &\;=\; x \\
  \end{flalign}
\end{proof}



\begin{thm}[Expansión de Shannon: suma]\label{thm:TExpShaS}
Dados dos elementos del conjunto base \(\setB\) de un álgebra de Boole \(
\mathscr{B} \), \( x \) e \( y \), \( x = ( x + y ) \cdot ( x + \overline{y}
) \).
\end{thm}



\begin{proof}
  \begin{flalign}
\forall x \; x \in \setB \;\;& \forall y \; y \in \setB {}&{}\\
x \;\; &\;=\; x \cdot 1 \;=\\
{} &\;=\; x \cdot \left( y \cdot \overline{y} \right) \;=\\
{} &\;=\; \left( x \cdot y \right) + \left( x \cdot \overline{y} \right)
\;\\
  \end{flalign}
\end{proof}




\begin{thm}[Expansión de Shannon: producto]\label{thm:TExpShaP}
Dados dos elementos del conjunto base \(\setB\) de un álgebra de Boole \(
\mathscr{B} \), \( x \) e \( y \), \( x = ( x \cdot y ) + ( x \cdot
\overline{y} ) \).
\end{thm}



\begin{proof}
  \begin{flalign}
\forall x \; x \in \setB \;\;& \forall y \; y \in \setB {}&{}\\
x \;\; &\;=\; x + 0 \;=\\
{} &\;=\; x + \left( y \cdot \overline{y} \right) \;=\\
{} &\;=\; \left( x + y \right) \cdot \left( x + \overline{y} \right)
\;\\
  \end{flalign}
\end{proof}



\begin{thm}[Simplificación elemento y negado: suma]\label{thm:TSimplOtraS}
Dados dos elementos del conjunto base \(\setB\) de un álgebra de Boole \(
\mathscr{B} \), \( x \) e \( y \), \( x + y = x + ( \overline{x} \cdot
y ) \).
\end{thm}



\begin{proof}
  \begin{flalign}
\forall x \; x \in \setB \;\;& \forall y \; y \in \setB {}&{}\\
x + \left( \overline{x} \cdot y \right)\;\; &\;=\; \left(x + \overline{x}\right) \cdot \left(x + y\right) \;=\\
{} &\;=\; 1 \cdot \left(x + y\right) \;=\\
{} &\;=\; \left(x + y\right) \cdot 1 \;=\\
{} &\;=\; x + y
  \end{flalign}
\end{proof}



\begin{thm}[Simplificación elemento y negado: producto] 
\label{thm:TSimplOtraP}
Dados dos elementos del conjunto base \(\setB\) de un álgebra de Boole \(
\mathscr{B} \), \( x \) e \( y \), \( x \cdot y = x \cdot ( \overline{x} +
y ) \).
\end{thm}



\begin{proof}
  \begin{flalign}
\forall x \; x \in \setB \;\;& \forall y \; y \in \setB {}&{}\\
x \cdot \left( \overline{x} + y \right) \;\; &\;=\;\\ 
{} &\;=\; \left(x \cdot \overline{x}\right) + \left( x \cdot y \right) \;=\\
{} &\;=\; 0 + \left( x \cdot y \right) \;=\\
{} &\;=\; \left( x \cdot y \right) + 0 \;=\\
{} &\;=\; x \cdot y
  \end{flalign}
\end{proof}



\section{Acercamiento a los retículos.}


\begin{thm}[Equivalencias en retículos]\label{thm:TEquivRet}
Dados dos elementos cualquiera del conjunto base \(\setB\) de un álgebra de
Boole \( \mathscr{B} \), \( x \) e \( y \), \( x \cdot y = x \iff x + y = y
\).
\end{thm}



\begin{proof}
  \begin{flalign}
  \forall x \; x \in \setB \;\;& \forall y \; y \in \setB\\
  \Longleftarrow &\\ \nonumber
  x \cdot y \;\; &\;=\; y\\ 
  x + (x \cdot y) &\;=\; x + y\\
  x &\;=\; x + y\\
  \Longrightarrow &\\ \nonumber
  x + y \;\; &\;=\; x\\ 
  y \cdot (x + y) &\;=\; y \cdot x\\
  y \cdot (y + x) &\;=\; y \cdot x\\
  y \cdot (y + x) &\;=\; x \cdot y\\
  y &\;=\; x \cdot y\\
  \end{flalign}
\end{proof}



\begin{thm}[Cancelación en retículos]\label{thm:TCancelRet}
Dados dos elementos cualquiera del conjunto base \(\setB\) de un álgebra de
Boole \( \mathscr{B} \), \( x \) e \( y \), si \( x + y = x \cdot y \;\then\;
x = y \).
\end{thm}



\begin{proof}
  \begin{flalign}
    \text{Dados los elementos }\quad& x,y \in \setB\\
    \\ \nonumber
    \text{Hipótesis} \quad& x + y = x \cdot y\\
    \text{Multiplicamos por la derecha} \quad& x \cdot ( x + y ) = x \cdot (
    x \cdot y )\\
    \text{Simplificación por Absorción} \quad& x = x \cdot ( x \cdot y )\\
    \\ \nonumber
    \text{Sumamos por la izquierda} \quad&  x + y \cdot y =  ( x \cdot ( x \cdot y ) ) + y\\
    \text{Distributiva} \quad&  x + y = (x + y) \cdot ( ( x \cdot y ) + y )\\
    \text{Simplificación por Absorción} \quad&  x + y = (x + y) \cdot y\\
    \text{Simplificación por Absorción} \quad&  x + y = y\\
    \text{Hipótesis} \quad& x = x \cdot ( x + y )\\
    \text{} \quad& x = x \cdot y = x + y  = y\\
  \end{flalign}
\end{proof}







\begin{thm}[Cancelación en sumas]\label{thm:TCancelS}
Dados dos elementos cualquiera del conjunto base \(\setB\) de un álgebra de
Boole \( \mathscr{B} \), \( x \) e \( y \), si existe \( z \) en \( \setB \)
\( ( x + z = y + z ) \ly ( x + \overline{z} = y + \overline{z} ) \then x = y
\).
\end{thm}



\begin{proof}
  \begin{flalign}
    \text{Dados los elementos }\quad& x,y \in \setB\\
    \text{Hipótesis 1} \quad& x + z = y + z\\
    \text{Hipótesis 2} \quad& x + \overline{z} = y + \overline{z}\\
    {} &\;\; \left(x + z\right)\cdot\left(x + \overline{z}\right) = \left(y +
    z\right)\cdot\left(y + \overline{z}\right)\\ 
    \text{Simplificación de Shannon} &\;\; x = y
  \end{flalign}
\end{proof}



\begin{thm}[Cancelación en productos]\label{thm:TCancelP}
Dados dos elementos cualquiera del conjunto base \(\setB\) de un álgebra de
Boole \( \mathscr{B} \), \( x \) e \( y \), si existe \( z \) en \( \setB \)
\( ( x \cdot z = y \cdot z ) \ly ( x \cdot \overline{z} = y \cdot
\overline{z} ) \then x = y
\).
\end{thm}



\begin{proof}
  \begin{flalign}
    \text{Dados los elementos }\quad& x,y \in \setB \\
    \text{Hipótesis 1} \quad& x \cdot z = y \cdot z \\
    \text{Hipótesis 2} \quad& x \cdot \overline{z} = y \cdot \overline{z} \\
    {} &\;\; \left(x \cdot z\right) + \left(x \cdot \overline{z}\right) =
    \left( y \cdot z \right) + \left( y \cdot \overline{z} \right) \\ 
    \text{Simplificación de Shannon} &\;\; x = y
  \end{flalign}
\end{proof}



\section{Pre-asociatividad y leyes de De Morgan.}



\begin{lem}[Pre-asociativa en sumas]\label{lem:LPreAsocS}
Dados dos elementos cualquiera del conjunto base \(\setB\) de un álgebra de
Boole \( \mathscr{B} \), \( x \) e \( y \), siempre se cumple \( x + y =
x + ( x + y ) \).
\end{lem}



\begin{proof}
  \begin{flalign}
    \forall x\quad \forall y\quad&\quad x \in \setB \;\ly\; y \in \setB\\
    \\ \nonumber
    \text{Término 1 : T1}\quad&\quad T1 = x  + y \\
    \\ \nonumber
    \text{Término 2 : T2}\quad&\quad T2 = x + ( x  + y ) = x + T1 \\
  \end{flalign}
  \begin{flalign}
    \text{Producto \(x\) a \(T1\)} \quad&\quad x \cdot (x  + y ) = x\\
  \end{flalign}
  \begin{flalign}
    \text{Producto \(\overline{x}\) a \(T1\)} \quad&\quad \overline{x} \cdot (x  + y ) = \overline{x} \cdot y
  \end{flalign}
  \begin{flalign}
    \text{Producto \(x\) a \(T2\)} \quad&\;\quad x \cdot ( x + (x  + y ) ) = \\
    \text{Distribución del producto} \quad&=\quad (x \cdot x) + ( x \cdot (x  + y ) ) =\\
    \text{Idempotencia producto} \quad&=\quad x + ( x \cdot (x  + y ) ) =\\
    \text{Simplificación Absorción producto} \quad&=\quad x \cdot x =\\
    \text{Idempotencia Producto} \quad&=\quad x
  \end{flalign}
  \begin{flalign}
  \text{Multiplico \(x\) a \(T2\)} \quad&\;\quad x \cdot ( x + (x  + y ) ) = x
  \end{flalign}
  \begin{flalign}
    \text{Multiplico \(\overline{x}\) a \(T2\)} \quad&\;\quad \overline{x} \cdot ( x + (x  + y ) ) = \\
    \text{Distribución del producto} \quad&=\quad (\overline{x} \cdot x) + ( \overline{x} \cdot (x  + y ) ) =\\
    \text{Conmutativa producto} \quad&=\quad (x \cdot \overline{x}) + ( \overline{x} \cdot (x  + y ) ) =\\
    \text{Elemento neutro} \quad&=\quad 0 + ( \overline{x} \cdot (x  + y ) ) =\\
    \text{Conmutativa Suma} \quad&=\quad ( \overline{x} \cdot (x  + y ) ) + 0 =\\
    \text{Neutro de la Suma} \quad&=\quad \overline{x} \cdot (x + y ) =\\
    \text{Simplificación otro} \quad&=\quad \overline{x} \cdot y
  \end{flalign}
  \begin{flalign}
  \text{Sumo \(\overline{x}\) a \(T2\)} \quad&\;\quad \overline{x} \cdot ( x + ( x + y ) ) = \overline{x} \cdot y
  \end{flalign}
  \begin{flalign}
  \text{Ecuación 1 : }&\quad x \cdot (x + y ) \;=\; x \cdot ( x + ( x + y ) )\\
  \text{Ecuación 2 : }&\quad \overline{x} \cdot ( x + y ) \;=\; \overline{x} \cdot ( x + (x  + y ) )\\
  \text{Cancelación en Retículos : }&\quad x + y = x + ( x + y )
  \end{flalign}
\end{proof}



\begin{lem}[Pre-asociativa en productos]\label{lem:LPreAsocP}
Dados dos elementos cualquiera del conjunto base \(\setB\) de un álgebra de
Boole \( \mathscr{B} \), \( x \) e \( y \), siempre se cumple \( x \cdot y =
x \cdot ( x \cdot y ) \).
\end{lem}



\begin{proof}
  \begin{flalign}
    \forall x\quad \forall y\quad&\quad x \in \setB \;\ly\; y \in \setB\\
    \\ \nonumber
    \text{Término 1 : T1}\quad&\quad T1 = x  \cdot y \\
    \\ \nonumber
    \text{Término 2 : T2}\quad&\quad T2 = x \cdot ( x  \cdot y ) = x \cdot T1 \\
  \end{flalign}
  \begin{flalign}
    \text{Sumo \(x\) a \(T1\)} \quad&\quad x + (x  \cdot y ) = x\\
  \end{flalign}
  \begin{flalign}
    \text{Sumo \(\overline{x}\) a \(T1\)} \quad&\quad \overline{x} + (x  \cdot y ) = \overline{x} + y
  \end{flalign}
  \begin{flalign}
    \text{Sumo \(x\) a \(T2\)} \quad&\;\quad x + ( x \cdot (x  \cdot y ) ) = \\
    \text{Distribución de la suma} \quad&=\quad (x + x) \cdot ( x + (x  \cdot y ) ) =\\
    \text{Idempotencia suma} \quad&=\quad x \cdot ( x + (x  \cdot y ) ) =\\
    \text{Simplificación Absorción suma} \quad&=\quad x + x =\\
    \text{Idempotencia Suma} \quad&=\quad x
  \end{flalign}
  \begin{flalign}
  \text{Sumo \(x\) a \(T2\)} \quad&\;\quad x + ( x \cdot (x  \cdot y ) ) = x
  \end{flalign}
  \begin{flalign}
    \text{Sumo \(\overline{x}\) a \(T2\)} \quad&\;\quad \overline{x} + ( x \cdot (x  \cdot y ) ) = \\
    \text{Distribución de la suma} \quad&=\quad (\overline{x} + x) \cdot ( \overline{x} + (x  \cdot y ) ) =\\
    \text{Conmutativa suma} \quad&=\quad (x + \overline{x}) \cdot ( \overline{x} + (x  \cdot y ) ) =\\
    \text{Elemento neutro} \quad&=\quad 1 \cdot ( \overline{x} + (x  \cdot y ) ) =\\
    \text{Conmutativa Producto} \quad&=\quad ( \overline{x} + (x  \cdot y ) ) \cdot 1 =\\
    \text{Neutro del Producto} \quad&=\quad \overline{x} + (x  \cdot y ) =\\
    \text{Simplificación otro} \quad&=\quad \overline{x} + y
  \end{flalign}
  \begin{flalign}
  \text{Sumo \(\overline{x}\) a \(T2\)} \quad&\;\quad \overline{x} + ( x \cdot (x  \cdot y ) ) = \overline{x} + y
  \end{flalign}
  \begin{flalign}
  \text{Ecuación 1 : }&\quad x + (x  \cdot y ) \;=\; x + ( x \cdot (x  \cdot y ) )\\
  \text{Ecuación 2 : }&\quad \overline{x} + (x  \cdot y ) \;=\; \overline{x} + ( x \cdot (x  \cdot y ) )\\
  \text{Cancelación en Retículos : }&\quad x  \cdot y = x \cdot (x  \cdot y )
  \end{flalign}
\end{proof}



\begin{lem}[Pre-asociativa 2 de la suma]\label{lem:PreAsoc2S}
Dado un álgebra de Boole \( \mathscr{B} \), con conjunto base \( \setB \), y dos elementos cualquiera de él, \( x \) e \( y \), entonces \( \overline{ x } + \left( { x + y } \right) = 1 \).
\end{lem}



\begin{proof}
\begin{flalign}
\forall x \qquad \forall y \qquad & x \in \setB \quad \ly \quad y \in \setB\\
\text{Término 1 : }&\qquad T_1 \mathdef \overline{x} + \left( x + y \right)\\
\text{Término 2 : }&\qquad T_2 \mathdef 1
\end{flalign}
\begin{flalign}
\text{Multiplico Término 1 por \(x\): }\qquad x \cdot T_1 \quad&=\quad x \cdot \left( \overline{x} + \left( x + y \right) \right) \quad=\quad\\
\text{Distribuyo el producto : }\qquad \quad&=\quad \left( x \cdot \overline{x} \right) + \left( x \cdot \left( x + y \right) \right) \quad=\quad\\
\text{Elemento Complementario : }\qquad \quad&=\quad 0 + \left( x \cdot \left( x + y \right) \right) \quad=\quad\\
\text{Conmutativa Suma : }\qquad \quad&=\quad \left( x \cdot \left( x + y \right) \right) + 0 \quad=\quad\\
\text{Neutro de la Suma : }\qquad \quad&=\quad x \cdot \left( x + y \right) \quad=\quad\\
\text{Simplificación Absorción Suma : }\qquad \quad&=\quad x
\end{flalign}
\begin{flalign}
\text{Multiplico Término 2 por \(x\): }\qquad x \cdot T_1 \quad&=\quad x\\
\end{flalign}
\begin{flalign}
\text{\( x \cdot T_1 \;=\; x \cdot T_2\): }\qquad &\quad x \cdot \left( \overline{x} + \left( x + y \right) \right) \quad=\quad x \cdot 1
\end{flalign}
\begin{flalign}
\text{Multiplico Término 1 por \(\overline{x}\): }\qquad \overline{x} \cdot T_1 \quad&=\quad \overline{x} \cdot \left( \overline{x} + \left( x + y \right) \right) \quad=\quad\\
\text{Distribuyo el producto : }\qquad \quad&=\quad \left( \overline{x} \cdot \overline{x} \right) + \left( \overline{x} \cdot \left( x + y \right) \right) \quad=\quad\\
\text{Idempotencia Producto : }\qquad \quad&=\quad \overline{x} + \left( \overline{x} \cdot \left( x + y \right) \right) \quad=\quad\\
\text{Simplificación absorción Producto : }\qquad \quad&=\quad \overline{x} + \left( \overline{x} \cdot y \right) \quad=\quad\\
\text{Simplificación Absorción Suma : }\qquad \quad&=\quad \overline{x}
\end{flalign}
\begin{flalign}
\text{Multiplico Término 2 por \(\overline{x}\): }\qquad \overline{x} \cdot T_2 \quad&=\quad \overline{x} \cdot 1 \quad=\quad \overline{x} \\
\end{flalign}
\begin{flalign}
\text{\( x \cdot T_1 \;=\; x \cdot T_2\): }\qquad &\quad x \cdot \left( \overline{x} + \left( x + y \right) \right) \quad=\quad x \cdot 1\\
\text{\( \overline{x} \cdot T_1 \;=\; \overline{x} \cdot T_2\): }\qquad &\quad \overline{x} \cdot \left( \overline{x} + \left( x + y \right) \right) \quad=\quad \overline{x} \cdot 1\\
\text{Cancelación en Retículos : }\qquad&\qquad \overline{x} + \left( x + y \right) = 1
\end{flalign}
\end{proof}






\begin{lem}[Pre-asociativa 2 del producto]\label{lem:PreAsoc2P}
Dado un álgebra de Boole \( \mathscr{B} \), con conjunto base \( \setB \), y dos elementos cualquiera de él, \( x \) e \( y \), entonces \( \overline{ x } \cdot \left( { x \cdot y } \right) = 0 \).
\end{lem}



\begin{proof}
\begin{flalign}
\forall x \qquad \forall y \qquad & x \in \setB \quad \ly \quad y \in \setB\\
\text{Término 1 : }&\qquad T_1 \mathdef \overline{x} \cdot \left( x \cdot y \right)\\
\text{Término 2 : }&\qquad T_2 \mathdef 0
\end{flalign}
\begin{flalign}
\text{Sumo a Término 1 \(x\): }\qquad x + T_1 \quad&=\quad x + \left( \overline{x} \cdot \left( x \cdot y \right) \right) \quad=\quad\\
\text{Distribuyo el producto : }\qquad \quad&=\quad \left( x + \overline{x} \right) \cdot \left( x + \left( x \cdot y \right) \right) \quad=\quad\\
\text{Elemento Complementario : }\qquad \quad&=\quad 1 \cdot \left( x + \left( x \cdot y \right) \right) \quad=\quad\\
\text{Conmutativa Producto : }\qquad \quad&=\quad \left( x + \left( x \cdot y \right) \right) \cdot 1 \quad=\quad\\
\text{Neutro del Producto : }\qquad \quad&=\quad x + \left( x \cdot y \right) \quad=\quad\\
\text{Simplificación Absorción Producto : }\qquad \quad&=\quad x
\end{flalign}
\begin{flalign}
\text{Sumo a Término 2 \(x\): }\qquad x + T_2 \quad&=\quad x + 0\\
\end{flalign}
\begin{flalign}
\text{\( x + T_1 \;=\; x + T_2\): }\qquad &\quad x + \left( \overline{x} \cdot \left( x \cdot y \right) \right) \quad=\quad x + 0
\end{flalign}
\begin{flalign}
\text{Sumo a Término 1 \(\overline{x}\): }\qquad \overline{x} + T_1 \quad&=\quad \overline{x} + \left( \overline{x} \cdot \left( x \cdot y \right) \right) \quad=\quad\\
\text{Distribuyo la suma : }\qquad \quad&=\quad \left( \overline{x} + \overline{x} \right) \cdot \left( \overline{x} + \left( x \cdot y \right) \right) \quad=\quad\\
\text{Idempotencia Suma : }\qquad \quad&=\quad \overline{x} \cdot \left( \overline{x} + \left( x \cdot y \right) \right) \quad=\quad\\
\text{Simplificación Absorción Suma : }\qquad \quad&=\quad \overline{x} \cdot \left( \overline{x} + y \right) \quad=\quad\\
\text{Simplificación Absorción Producto : }\qquad \quad&=\quad \overline{x}
\end{flalign}
\begin{flalign}
\text{Sumo a Término 2 \(\overline{x}\): }\qquad \overline{x} + T_2 \quad&=\quad \overline{x} + 0 \quad=\quad \overline{x} \\
\end{flalign}
\begin{flalign}
\text{\( x + T_1 \;=\; x + T_2\): }\qquad &\quad x + \left( \overline{x} \cdot \left( x \cdot y \right) \right) \quad=\quad x + 0\\
\text{\( \overline{x} + T_1 \;=\; \overline{x} + T_2\): }\qquad &\quad \overline{x} + \left( \overline{x} \cdot \left( x \cdot y \right) \right) \quad=\quad \overline{x} + 0\\
\text{Cancelación en Retículos : }\qquad&\qquad \overline{x} \cdot \left( x \cdot y \right) = 0
\end{flalign}
\end{proof}



\begin{thm}[Teorema de Morgan, suma]
Dados dos elementos cualquiera, \(x \) e \(y\), del conjunto subyacente \(\setB\) de un álgebra de Boole \(\mathscr{B}\), \( \overline{x + y} = \overline{x} \cdot \overline{y} \).
\end{thm}



\begin{proof}

\begin{flalign}
\forall x \qquad \forall y \qquad & \qquad x \in \setB \qquad y \in \setB\\
\text{Término 1 : }\qquad T_1 \quad&\mathdef\quad x + y\\
\text{Término 2 : }\qquad T_2 \quad&\mathdef\quad \overline{x} \cdot \overline{y}
\end{flalign}

\begin{align}
\text{Multiplicamos los términos }\quad& \left( { x + y } \right) \cdot \left( \overline{x} \cdot \overline{y} \right) \quad&=\\
\text{Conmutativa del producto }\quad&=\; \left( \overline{x} \cdot \overline{y} \right) \cdot \left( { x + y } \right) \;&=\\
\text{Distribución del producto }\quad&=\; \left( \left( \overline{x} \cdot \overline{y} \right) \cdot x \right) + \left( \left( \overline{x} \cdot \overline{y} \right) \cdot y \right)\;&=\\
\text{Conmutativa del producto }\quad&=\; \left( x \cdot \left( \overline{x} \cdot \overline{y} \right) \right) + \left( \left( \overline{x} \cdot \overline{y} \right) \cdot y \right)\;&=\\
\text{Lemas }\quad&=\; 0 + \left( \left( \overline{x} \cdot \overline{y} \right) \cdot y \right)\;&=\\
\text{Conmutativa de la suma }\quad&=\; \left( \left( \overline{x} \cdot \overline{y} \right) \cdot y \right) + 0 \;&=\\
\text{Neutro de la suma }\quad&=\; \left( \overline{x} \cdot \overline{y} \right) \cdot y \;&=\\
\text{Conmutativa del producto }\quad&=\; \left( \overline{y} \cdot \overline{x} \right) \cdot y \;&=\\
\text{Conmutativa del producto }\quad&=\; y \cdot \left( \overline{y} \cdot \overline{x} \right) \;&=\\
\text{Lemas }\quad&=\; 0 &{}
\end{align}

\begin{align}
\forall x \qquad \forall y \qquad & \qquad x \in \setB \qquad y \in \setB\\
{} & \left( x + y \right) \cdot \left( \overline{x}\cdot \overline{y} \right) = 0
\end{align}


\begin{align}
\text{Sumamos los términos }\quad& \left( { x + y } \right) + \left( \overline{x} \cdot \overline{y} \right) \quad&=\\
\text{Conmutativa de la suma }\quad&=\; \left( \overline{x} \cdot \overline{y} \right) + \left( { x + y } \right) \quad&=\\
\text{Lemas }\quad&=\; 1 + \left( \overline{y} + \left( x + y \right) \right) \;&=\\
\text{Conmutativa de la suma }\quad&=\; 1 + \left( \overline{y} + \left( y + x \right) \right) \;&=\\
\text{Doble negación es identidad }\quad&=\; 1 + \left( \overline{y} + \left( \overline{\overline{y}} + x \right) \right) \;&=\\
\text{Lemas }\quad&=\; 1 + 1 \;&=\\
\text{Elemento absorbente de la suma }\quad&=\; 1
\end{align}



\begin{align}
\forall x \qquad& \forall y \qquad \qquad x \in \setB \qquad& y \in \setB\\
{} & \left( x + y \right) + \left( \overline{x}\cdot \overline{y} \right) \;&=\; 1\\
{} & \left( x + y \right) \cdot \left( \overline{x}\cdot \overline{y} \right) \;&=\; 0\\
\\ \nonumber
\text{Por unicidad del complementario }\quad&\overline{ {x + y} } &\;=\;\overline{x}\cdot\overline{y}
\end{align}




\end{proof}






\begin{thm}[Teorema de Morgan, producto]
Dados dos elementos cualquiera, \(x \) e \(y\), del conjunto subyacente \(\setB\) de un álgebra de Boole \(\mathscr{B}\), \( \overline{x \cdot y} = \overline{x} + \overline{y} \).
\end{thm}



\begin{proof}

\begin{flalign}
\forall x \qquad \forall y \qquad & \qquad x \in \setB \qquad y \in \setB\\
\text{Término 1 : }\qquad T_1 \quad&\mathdef\quad x \cdot y\\
\text{Término 2 : }\qquad T_2 \quad&\mathdef\quad \overline{x} + \overline{y}
\end{flalign}

\begin{align}
\text{Multiplicamos los términos }\quad& \left( { x \cdot y } \right) + \left( \overline{x} + \overline{y} \right) \quad&=\\
\text{Conmutativa de la suma }\quad&=\; \left( \overline{x} + \overline{y} \right) + \left( { x \cdot y } \right) \;&=\\
\text{Distribución de la suma }\quad&=\; \left( \left( \overline{x} \cdot \overline{y} \right) + x \right) \cdot \left( \left( \overline{x} + \overline{y} \right) + y \right)\;&=\\
\text{Conmutativa de la suma }\quad&=\; \left( x + \left( \overline{x} + \overline{y} \right) \right) \cdot \left( \left( \overline{x} + \overline{y} \right) + y \right)\;&=\\
\text{Lemas }\quad&=\; 1 \cdot \left( \left( \overline{x} + \overline{y} \right) + y \right)\;&=\\
\text{Conmutativa del producto }\quad&=\; \left( \left( \overline{x} + \overline{y} \right) + y \right) \cdot 1 \;&=\\
\text{Neutro del producto }\quad&=\; \left( \overline{x} + \overline{y} \right) + y \;&=\\
\text{Conmutativa de la suma }\quad&=\; y + \left( \overline{y} + \overline{x} \right) \;&=\\
\text{Lemas }\quad&=\; 1 &{}
\end{align}

\begin{align}
\forall x \qquad \forall y \qquad & \qquad x \in \setB \qquad y \in \setB\\
{} & \left( x \cdot y \right) + \left( \overline{x} + \overline{y} \right) = 1
\end{align}


\begin{align}
\text{Multiplicamos los términos }\quad& \left( { x \cdot y } \right) \cdot \left( \overline{x} + \overline{y} \right) \quad&=\\
\text{Distribuimos el producto }\quad&=\; \left( \left( x \cdot y \right) \cdot \overline{x} \right) + \left( \left( x \cdot y \right) \cdot \overline{y} \right) \quad&=\\
\text{Conmutativa del producto }\quad&=\; \left( \overline{x} \cdot \left( x \cdot y \right) \right) + \left( \left( x \cdot y \right) \cdot \overline{y} \right) \quad&=\\ 
\text{Doble negación es identidad }\quad&=\; \left( \overline{x} \cdot \left( \overline{\overline{x}} \cdot y \right) \right) + \left( \left( x \cdot y \right) \cdot \overline{y} \right) \quad&=\\ 
\text{Lemas }\quad&=\; 0 + \left( \left( x \cdot y \right) \cdot \overline{y} \right) \quad&=\\
\text{Conmutativa del producto }\quad&=\; 0 + \left( \overline{y} \cdot \left( x \cdot y \right) \right) \quad&=\\
\text{Conmutativa del producto }\quad&=\; 0 + \left( \overline{y} \cdot \left( y \cdot x \right) \right) \quad&=\\
\text{Doble negación es identidad }\quad&=\; 0 + \left( \overline{y} \cdot \left( \overline{\overline{y}} \cdot x \right) \right) \quad&=\\
\text{Lemas }\quad&=\; 0 + 0 \quad&=\\
\text{Elemento Absorbente de la suma }\quad&=\; 0 &\\
\end{align}



\begin{align}
\forall x \qquad& \forall y \qquad \qquad x \in \setB \qquad& y \in \setB\\
{} & \left( x \cdot y \right) + \left( \overline{x} + \overline{y} \right) \;&=\; 1\\
{} & \left( x \cdot y \right) + \left( \overline{x} + \overline{y} \right) \;&=\; 0\\
\\ \nonumber
\text{Por unicidad del complementario }\quad&\overline{ {x \cdot y} } &\;=\; \overline{x} + \overline{y}
\end{align}
\end{proof}



\begin{thm}[Suma como complementario del producto de complementarios]
Dado un álgebra de Boole \(\mathscr{B}\) con conjunto subyacente \(\setB\), cualquier pareja de elementos de \(\setB\), \(x\) e \(y\), podemos poner \(x+y=\overline{\overline{x}\cdot\overline{y}}\).
\end{thm}



\begin{proof}
Por simple negación de los términos de la igualdad del teorema de Morgan correspondiente, \(\overline{x+y} = \overline{x}\cdot\overline{y}\).
\end{proof}



\begin{thm}[Producto como complementario de la suma de complementarios]
Dado un álgebra de Boole \(\mathscr{B}\) con conjunto subyacente \(\setB\), cualquier pareja de elementos de \(\setB\), \(x\) e \(y\), podemos poner \(x \cdot y=\overline{\overline{x} + \overline{y}}\).
\end{thm}


\begin{proof}
Por simple negación de los términos de la igualdad del teorema de Morgan correspondiente, \(\overline{x \cdot y} = \overline{x} + \overline{y}\).
\end{proof}



\section{Propiedad asociativa.}



\begin{lem}[Pseudo-asociativa suma 1]
Dado un álgebra de Boole, \(\mathscr{B}\), con conjunto subyacente \(\setB\). Dados tres elementos cualesquiera \(a\), \(b\) y \(c\) de \(\setB\), \( a \cdot ( ( a + b ) + c ) = a \cdot ( a + ( b + c ) )\).
\end{lem}



\begin{proof}
\begin{align}
\text{Definimos Término 1}\qquad& T_1 \;&\mathdef\; \left( \left( a + b \right) + c \right) \\
\text{Definimos Término 2}\qquad& T_2 \;&\mathdef\; \left( a + \left( b + c \right) \right)
\end{align}
\begin{align}
\text{Multiplicamos \(a \cdot T_1\)}\qquad& a \cdot T_1 \;&=\; a \cdot \left( \left( a + b \right) + c \right) \;=\\
\text{Distribuimos el producto}\qquad&{}&=\; \left( a \cdot \left( a + b \right) \right) + \left( a \cdot c \right) \;=\\
\text{Simplificación en retículos}\qquad&{}&=\; a + \left( a \cdot c \right) \;=\\
\text{Simplificación en retículos}\qquad&{}&=\; a \;
\end{align}
\begin{align}
a \cdot T_1 \;=\; a \cdot \left( \left( a + b \right) + c \right) \;=\; a
\end{align}
\begin{align}
\text{Multiplicamos \(a \cdot T_2\)}\qquad& a \cdot T_2 \;&=\; a \cdot \left( a + \left( b  + c \right) \right) \;=\\
\text{Distribuimos el producto}\qquad&{}&=\; \left( a \cdot a \right) + \left( a \cdot \left( b + c \right) \right) \;=\\
\text{Idempotencia de la suma}\qquad&{}&=\; a + \left( a \cdot \left( b + c \right) \right) \;=\\
\text{Simplificación en retículos}\qquad&{}&=\; a \;
\end{align}
\begin{align}
a \cdot T_1 \;=\; a \cdot \left( \left( a + b \right) + c \right) \;=\; a\\
a \cdot T_2 \;=\; a \cdot \left( a + \left( b + c \right) \right) \;=\; a\\
a \cdot \left( \left( a + b \right) + c \right) \;=\; a \cdot \left( a + \left( b + c \right) \right) 
\end{align}
\end{proof}



\begin{lem}[Pseudo-asociativa suma 2]
Dado un álgebra de Boole, \(\mathscr{B}\), con conjunto subyacente \(\setB\). Dados tres elementos cualesquiera \(a\), \(b\) y \(c\) de \(\setB\), \( \overline{a} \cdot ( ( a + b ) + c ) = \overline{a} \cdot ( a + ( b + c ) )\).
\end{lem}



\begin{proof}
\begin{align}
\text{Definimos Término 1}\qquad& T_1 \;&\mathdef\; \left( \left( a + b \right) + c \right) \\
\text{Definimos Término 2}\qquad& T_2 \;&\mathdef\; \left( a + \left( b + c \right) \right)
\end{align}
\begin{align}
\text{Multiplicamos \(\overline{a} \cdot T_1\)}\qquad& \overline{a} \cdot T_1 \;&=\; \overline{a} \cdot \left( \left( a + b \right) + c \right) \;=\\
\text{Distribuimos el producto}\qquad&{}&=\; \left( \overline{a} \cdot \left( a + b \right) \right) + \left( \overline{a} \cdot c \right) \;=\\
\text{Doble negación}\qquad&{}&=\; \left( \overline{a} \cdot \left( \overline{\overline{a}} + b \right) \right) + \left( \overline{a} \cdot c \right) \;=\\
\text{Simplificación en retículos}\qquad&{}&=\; \left( \overline{a} \cdot b \right) + \left( \overline{a} \cdot c \right) \;=\\
\text{Factor común}\qquad&{}&=\; \overline{a} \cdot \left( b +  c \right) \;\;\\
\end{align}
\begin{align}
\overline{a} \cdot T_1 \;=\; \overline{a} \cdot \left( \left( a + b \right) + c \right) \;=\; \overline{a} \cdot \left( b + c \right)
\end{align}
\begin{align}
\text{Multiplicamos \(\overline{a} \cdot T_2\)}\qquad& \overline{a} \cdot T_2 \;&=\; \overline{a} \cdot \left( a + \left( b  + c \right) \right) \;=\\
\text{Distribuimos el producto}\qquad&{}&=\; \left( \overline{a} \cdot a \right) + \left( \overline{a} \cdot \left( b + c \right) \right) \;=\\
\text{Elemento complementario}\qquad&{}&=\; 0 + \left( \overline{a} \cdot \left( b + c \right) \right) \;=\\
\text{Conmutativa de la suma}\qquad&{}&=\; \left( \overline{a} \cdot \left( b + c \right) \right) + 0 \;=\\
\text{Neutro de la suma}\qquad&{}&=\; \overline{a} \cdot \left( b + c \right) \;\;\\
\end{align}
\begin{align}
\overline{a} \cdot T_1 \;=\; \overline{a} \cdot \left( \left( a + b \right) + c \right) \;=\; \overline{a} \cdot \left( b + c \right)\\
\overline{a} \cdot T_2 \;=\; \overline{a} \cdot \left( a + \left( b + c \right) \right) \;=\; \overline{a} \cdot \left( b + c \right)\\
\overline{a} \cdot \left( \left( a + b \right) + c \right) \;=\; \overline{a} \cdot \left( a + \left( b + c \right) \right) 
\end{align}
\end{proof}



\begin{thm}[Asociativa de la suma]
Dado un álgebra de Boole, \(\mathscr{B}\), con conjunto subyacente \(\setB\). Dados tres elementos cualesquiera \(a\), \(b\) y \(c\) de \(\setB\), \( ( a + b ) + c = a + ( b + c ) \).
\end{thm}



\begin{proof}
\begin{align}
\text{Del lema pseudo-asociativa de la suma 1}\\
\qquad\quad a \cdot \left( \left( a + b \right) + c \right) \;=\; a \cdot \left( a + \left( b + c \right) \right) \\ 
\text{Del lema pseudo-asociativa de la suma 2}\\
\qquad\quad \overline{a} \cdot \left( \left( a + b \right) + c \right) \;=\; \overline{a} \cdot \left( a + \left( b + c \right) \right) \\
\text{Sumando los términos derechos de la igualdad}\\
\qquad\quad a \cdot \left( \left( a + b \right) + c \right) + \overline{a} \cdot \left( \left( a + b \right) + c \right) \;=\; \left( a + b \right) + c\\
\text{Sumando los términos izquierdos de las igualdades}\\ 
\qquad\quad a \cdot \left( a + \left( b + c \right) \right) + \overline{a} \cdot \left( a + \left( b + c \right) \right) \;=\; a + \left( b + c \right)\\
\text{Rehaciendo los términos derecho e izquierdo de las igualdades}\\ 
\left( a + b \right) + c \quad = \quad  a + \left( b + c \right)\\
\end{align}
\end{proof}



\begin{lem}[Pseudo-asociativa producto 1]
Dado un álgebra de Boole, \(\mathscr{B}\), con conjunto subyacente \(\setB\). Dados tres elementos cualesquiera \(a\), \(b\) y \(c\) de \(\setB\), \( a + ( ( a \cdot b ) \cdot c ) = a + ( a \cdot ( b \cdot c ) )\).
\end{lem}



\begin{proof}

\begin{align}
\text{Definimos Término 1}\qquad& T_1 \;&\mathdef\; \left( \left( a \cdot b \right) \cdot c \right) \\
\text{Definimos Término 2}\qquad& T_2 \;&\mathdef\; \left( a \cdot \left( b \cdot c \right) \right)
\end{align}

\begin{align}
\text{Sumamos \(a + T_1\)}\qquad& a + T_1 \;&=\; a + \left( \left( a \cdot b \right) \cdot c \right) \;=\\
\text{Distribuimos la suma}\qquad&{}&=\; \left( a + \left( a \cdot b \right) \right) \cdot \left( a + c \right) \;=\\
\text{Simplificación en retículos}\qquad&{}&=\; a \cdot \left( a + c \right) \;=\\
\text{Simplificación en retículos}\qquad&{}&=\; a \;
\end{align}

\begin{align}
a + T_1 \;=\; a + \left( \left( a \cdot b \right) \cdot c \right) \;=\; a
\end{align}

\begin{align}
\text{Sumamos \(a + T_2\)}\qquad& a + T_2 \;&=\; a + \left( a \cdot \left( b  \cdot c \right) \right) \;=\\
\text{Distribuimos la suma}\qquad&{}&=\; \left( a + a \right) \cdot \left( a + \left( b \cdot c \right) \right) \;=\\
\text{Idempotencia del producto}\qquad&{}&=\; a \cdot \left( a + \left( b \cdot c \right) \right) \;=\\
\text{Simplificación en retículos}\qquad&{}&=\; a \;
\end{align}

\begin{align}
a + T_1 \;=\; a + \left( \left( a \cdot b \right) \cdot c \right) \;=\; a\\
a + T_2 \;=\; a + \left( a \cdot \left( b \cdot c \right) \right) \;=\; a\\
a + \left( \left( a \cdot b \right) \cdot c \right) \;=\; a + \left( a \cdot \left( b \cdot c \right) \right) 
\end{align}

\end{proof}



\begin{lem}[Pseudo-asociativa producto 2]
Dado un álgebra de Boole, \(\mathscr{B}\), con conjunto subyacente \(\setB\). Dados tres elementos cualesquiera \(a\), \(b\) y \(c\) de \(\setB\), \( \overline{a} + ( ( a \cdot b ) \cdot c ) = \overline{a} + ( a \cdot ( b \cdot c ) )\).
\end{lem}



\begin{proof}
\begin{align}
\text{Definimos Término 1}\qquad& T_1 \;&\mathdef\; \left( \left( a \cdot  b \right) \cdot  c \right) \\
\text{Definimos Término 2}\qquad& T_2 \;&\mathdef\; \left( a \cdot \left( b \cdot c \right) \right)
\end{align}
\begin{align}
\text{Sumamos \(\overline{a} + T_1\)}\qquad& \overline{a} + T_1 \;&=\; \overline{a} + \left( \left( a \cdot b \right) \cdot c \right) \;=\\
\text{Distribuimos la suma}\qquad&{}&=\; \left( \overline{a} + \left( a \cdot b \right) \right) \cdot \left( \overline{a} + c \right) \;=\\
\text{Doble negación}\qquad&{}&=\; \left( \overline{a} + \left( \overline{\overline{a}} \cdot b \right) \right) \cdot \left( \overline{a} + c \right) \;=\\
\text{Simplificación en retículos}\qquad&{}&=\; \left( \overline{a} + b \right) \cdot \left( \overline{a} + c \right) \;=\\
\text{Sumando común}\qquad&{}&=\; \overline{a} + \left( b \cdot  c \right) \;\;\\
\end{align}
\begin{align}
\overline{a} + T_1 \;=\; \overline{a} + \left( \left( a \cdot b \right) \cdot c \right) \;=\; \overline{a} + \left( b \cdot c \right)
\end{align}
\begin{align}
\text{Sumamos \(\overline{a} + T_2\)}\qquad& \overline{a} + T_2 \;&=\; \overline{a} + \left( a \cdot \left( b \cdot c \right) \right) \;=\\
\text{Distribuimos la suma}\qquad&{}&=\; \left( \overline{a} + a \right) + \left( \overline{a} + \left( b \cdot c \right) \right) \;=\\
\text{Elemento complementario}\qquad&{}&=\; 1 \cdot \left( \overline{a} + \left( b \cdot c \right) \right) \;=\\
\text{Conmutativa del producto}\qquad&{}&=\; \left( \overline{a} + \left( b \cdot c \right) \right) \cdot 1 \;=\\
\text{Neutro del producto}\qquad&{}&=\; \overline{a} + \left( b \cdot c \right) \;\;\\
\end{align}
\begin{align}
\overline{a} + T_1 \;=\; \overline{a} + \left( \left( a \cdot b \right) \cdot c \right) \;=\; \overline{a} + \left( b \cdot c \right)\\
\overline{a} + T_2 \;=\; \overline{a} + \left( a \cdot \left( b \cdot c \right) \right) \;=\; \overline{a} + \left( b \cdot c \right)\\
\overline{a} + \left( \left( a \cdot b \right) \cdot c \right) \;=\; \overline{a} + \left( a \cdot \left( b \cdot c \right) \right) 
\end{align}
\end{proof}



\begin{thm}[Asociativa del producto]
Dado un álgebra de Boole, \(\mathscr{B}\), con conjunto subyacente \(\setB\). Dados tres elementos cualesquiera \(a\), \(b\) y \(c\) de \(\setB\), \( ( a \cdot b ) \cdot c = a \cdot ( b \cdot c ) \).
\end{thm}



\begin{proof}
\begin{align}
\text{Del lema pseudo-asociativa del producto 1}\\
\qquad\quad a + \left( \left( a \cdot b \right) \cdot c \right) \;=\; a + \left( a \cdot \left( b \cdot c \right) \right) \\ 
\text{Del lema pseudo-asociativa del producto 2}\\
\qquad\quad \overline{a} + \left( \left( a \cdot b \right) \cdot c \right) \;=\; \overline{a} + \left( a \cdot \left( b \cdot c \right) \right) \\
\text{Multiplicando los términos derechos de la igualdad}\\
\qquad\quad \left(a + \left( \left( a \cdot b \right) \cdot c \right) \right) \cdot \left( \overline{a} + \left( \left( a \cdot b \right) \cdot c \right) \right) \;=\; \left( a \cdot b \right) \cdot c\\
\text{Multiplicando los términos izquierdos de las igualdades}\\ 
\qquad\quad \left( a + \left( a \cdot \left( b \cdot c \right) \right) \right) \cdot \left( \overline{a} + \left( a \cdot \left( b \cdot c \right) \right) \right) \;=\; a \cdot \left( b \cdot c \right)\\
\text{Rehaciendo los términos derecho e izquierdo de las igualdades}\\ 
\left( a \cdot b \right) \cdot c \quad = \quad  a \cdot \left( b \cdot c \right)
\end{align}
\end{proof}




\chapter{Retículos de Boole.}



	La estructura de retículo \(\retL\) es una estructura dónde,
primeramente, se da una relación de orden sobre un conjunto base \(\setL\).
Llamaremos a esta relación de orden \(\cdot \leqq \cdot\). Como cualquier
relación de orden ha de cumplir la propiedades reflexiva, antisimétrica y
transitiva. Esto es, si el orden se realiza sobre un conjunto \(\setL\),
\(a,b,c\;\in\setL\), se establece que \(a\leqq a\), \(b\leqq b\) y \(c \leqq
c\) (propiedad \textsf{reflexiva}). Además, si \(a \leqq b\) y \(b \leqq a\)
entonces tenemos que \(a = b\), e igualmente con \(a,c\) y \(b,c\) (propiedad
\textsf{antisimétrica}). Por último, si \(a\leqq b\) y \(b \leqq c\) entonces
\(a \leqq c\) (propiedad \textsf{transitiva}).



	De esta forma, lo primero que pedimos a un conjunto es una relación de
orden. Este orden no tiene porqué ser total (o, lo que es lo mismo, lineal),
sino que en la mayor parte de los casos será un orden parcial, esto es, dos
elementos del conjunto no tienen porqué estar relacionados mediante este
orden.



	Sin embargo pediremos algunas propiedades de conexión para el orden. La
fundamental será que existan un elemento mínimo y un elemento máximo en el
propio conjunto, esto es, un elemento \(0 \in \setL\) tal que cualquier
elemento del retículo ha de estar relacionado con él, \(\forall \;x\; \in
\setL\) se tiene que cumplir que \(0\leqq x\). Para el máximo \(1\in\setL\),
pediremos que \(\forall \; x \; \in \setL\) se cumpla \( x \leqq 1 \). Con
esto podemos decir nuestro sistema con orden tiene máximo y mínimo
\(\mathscr{L} =\left\langle \setL , \leqq , \left\lbrace 0,1 \right\rbrace
\right\rangle\).



	Pero con esto aún no tenemos un retículo. Para que sea un retículo,
necesitamos que \(\forall \;x,y\; \in \setL\) exista un elemento mínimo tal
que \(\alpha  = \min\setCompr{z\in\setL}{x \leqq z \textsf{ y } y \leqq z}\),
y que denotaremos por \(\alpha = x \vee y\). Ahora sí que a este conjunto
ordenado le podemos llamar semi-retículo con \(\sup\). Para que sea retículo
ha de existir también \(\omega = \max\setCompr{z\in\setL}{z \leqq x \textsf{
y } z \leqq y}\) y que representamos por \(\omega = x \wedge y\). Un retículo
tiene \(\sup\) e \(\inf\).



	Además necesitamos elementos complementarios y leyes distributivas.



	El elemento complementario consiste en que a cada elemento del retículo
\(a\in\setL\) corresponde otro elemento del retículo
\(b\in\setL\) que \(a\wedge b = 0\) y \(a\vee b = 1\).



	Existe una generalización muy importante de los retículos de Boole, los
retículos de Heyting, en que solo se pide un complementario parcial. Esta
generalización es importante para formalizar lógicas no clásicas como las
lógicas intuicionistas. Así, la lógica clásica o buleana sería un caso
extremo de las lógicas intuicionistas.



	Por último nos queda la distribución de la operación \(\wedge\) sobre la
\(\vee\). Ésta trae consigo la distribución de la operación \(\vee\) sobre la
\(\wedge\). Los retículos en general no cumplen estas distributividades.



\section{Relaciones de orden a partir de un álgebra de Boole.}



	Viendo esto, podemos reinterpretar todo lo visto en las álgebras de
Boole, como un retículo \(\algB\), dónde se ha establecido un orden en su
conjunto subyacente \(\setB\), \(\leq\) tal que \(x+y = y \;\iff \; x \leq
y\) y \(x \cdot y = x \;\iff \; x \leq y\). Estas condiciones son
equivalentes y está probado como teorema que \(x+y = y \;\iff \; x\cdot y =
x\). De esta forma podemos trocar la suma buleana en el ínfimo de los
elementos mayores que dos elementos \(x\;y\;\;\in\setB\), \(x\vee y \mathdef
x+y\), y también \( x \wedge y \mathdef x \cdot y \). Los elementos neutros
pasan a ser el mínimo \(0 = \min\setB\) y \(1 = \max\setB\). Además, para un
elemento \(x\in\setB\setminus\set{0,1}\), \(\neg x\) es un elemento no
relacionado con \(x\) por la relación de orden. Esto es, \(x \nleq \neg x\) y
\(\neg x \nleq x\). El caso de \( \set{0,1} = \setB_2 \subseteq \setB \) es
la única exclusión por representar un retículo lineal (el único).



%	De alguna manera podemos considerar que entre un elemento y su
%complementario existe una distancia máxima. Esta distancia sería mínima si
%los elementos \(x\) e \(y\) fuesen consecutivos: \(x \precsim y\) si \( y
%\in \inf \setCompr{z \in \setB}{x \leq z} \) y \( x \in \sup \setCompr{z \in
%\setB}{z \leq x}\) y mayor cuanto más sucesivos intermedios haya. Cuando no
%haya sucesivos intermedios crecería de nuevo, siendo la distancia mayor la
%que conecta con los elementos que sumados dan más cerca del máximo y
%multiplicados más cerca del mínimo. La idea es informal.



\begin{defn}
Dado un álgebra de Boole \(\algB = \left\langle \setB, +, \cdot
\right\rangle\), que cumple los postulados de Huntington, obtenemos una
relación entre sus elementos que definimos por \(\leqB\), definida
como sigue.
\begin{equation}
\forall x \in \setB \; \forall y \in \setB \quad \left( x \leqB y \logicdef x 
\cdot y = x \right)
\end{equation}
\end{defn}



\begin{thm}
Dado un álgebra de Boole \(\algB = \left\langle \setB, +, \cdot
\right\rangle\), que cumple los postulados de Huntington, obtenemos que el
conjunto subyacente junto la relación interna anteriormente definida
\(\langle\setB,\leqB\rangle\) constituye un orden.
\end{thm}



\begin{proof}
\begin{eqnarray}
\nonumber
&&\textsf{Reflexiva en \(\langle\setB,\leqB\rangle\)  }\\
&&\text{Por idempotencia }\qquad\quad x = x \cdot x \then x \leqB x\\
&&\textsf{Antisimétrica en \(\langle \setB , \leqB \rangle\)  }\\
&&\text{Por transitividad  de \(=\) } \quad x \cdot y = x \textsf{ y } x
\cdot y = y \then x = y\\
&&\text{Por definición de \( \leqB \) } \;\qquad x \leqB y
\textsf{ y } y \leqB x \then x = y\\
&&\textsf{Transitiva en \( \langle \setB, \leqB \rangle \)  }\\
&&\text{Hipótesis de transitividad } \quad x \leqB y \textsf{ y } y
\leqB z\\
&&\text{Por definición de \( \leqB \) } \;\qquad x \cdot y = x \textsf{
y } y \cdot z = y\\
&&\qquad\;\qquad x \cdot y = x \then (x \cdot y) \cdot z = x \cdot z\\
&&\qquad\;\qquad x \cdot y = x \then x \cdot (y \cdot z) = x \cdot z\\
&&\qquad\;\qquad x \cdot y = x \textsf{ y } y \cdot z = y \then x \cdot y = x
\cdot z\\
&&\qquad\;\qquad x \cdot y = x \textsf{ y } y \cdot z = y \then x = x \cdot y
= x \cdot z\\
&&\qquad\;\qquad x \cdot y = x \textsf{ y } y \cdot z = y \then x \cdot z =
x\\
&&\qquad\;\qquad x \leqB y \textsf{ y } y \leqB z \then x
\leqB z
\end{eqnarray}
\end{proof}



\begin{defn}
Dado un álgebra de Boole \(\algB = \left\langle \setB, +, \cdot
\right\rangle\), que cumple los postulados de Huntington, obtenemos una
relación entre sus elementos que definimos por \( \geqB \), definida
como sigue.
\begin{equation}
\forall x \in \setB \; \forall y \in \setB \quad \left( x \geqB y \logicdef x
+ y = x \right)
\end{equation}
\end{defn}



\begin{thm}
Dado un álgebra de Boole \(\algB = \left\langle \setB, +, \cdot
\right\rangle\), que cumple los postulados de Huntington, obtenemos que el
conjunto subyacente junto la relación interna anteriormente definida
\(\langle \setB , \geqB \rangle\) constituye un orden.
\end{thm}



\begin{proof}
\begin{eqnarray}
\nonumber
&&\textsf{Reflexiva en \(\langle\setB,\geqB\rangle\)  }\\
&&\text{Por idempotencia }\qquad\quad x = x + x \then x \geqB x\\
&&\textsf{Antisimétrica en \(\langle\setB,\geqB\rangle\)  }\\
&&\text{Por transitividad  de \(=\) } \quad x + y = x \textsf{ y } x + y = y
\then x = y\\
&&\text{Por definición de \(\geqB\) } \;\qquad x \geqB y
\textsf{ y } y \geqB x \then x = y\\
&&\textsf{Transitiva en \(\langle\setB,\geqB\rangle\)  }\\
&&\text{Hipótesis de transitividad } \quad x \geqB y \textsf{ y } y
\geqB z\\
&&\text{Por definición de \(\geqB\) } \;\qquad x + y = x \textsf{ y }
y + z = y\\
&&\qquad\;\qquad x + y = x \then (x + y) + z = x + z\\
&&\qquad\;\qquad x + y = x \then x + (y + z) = x + z\\
&&\qquad\;\qquad x + y = x \textsf{ y } y + z = y \then x + y = x + z\\
&&\qquad\;\qquad x + y = x \textsf{ y } y + z = y \then x = x + y = x + z\\
&&\qquad\;\qquad x + y = x \textsf{ y } y + z = y \then x + z = x\\
&&\qquad\;\qquad x \geqB y \textsf{ y } y \geqB z \then x
\geqB z
\end{eqnarray}
\end{proof}



\begin{thm}
Los órdenes anteriormente definidos son compatibles entre sí, de forma que \(
\forall x,y \in \setB \) se cumple que \( x \leqB y \iff y
\geqB x \).
\end{thm}



\begin{proof}
\begin{eqnarray}
\nonumber
&&\text{Por \ref{thm:TEquivRet} de equivalencia reticular }\\
&&\forall x,y \in \setB \quad x \cdot y = x \iff x + y = y\\ \nonumber
&&\text{Por las definiciones de orden anteriores }\\
&&\forall x,y \in \setB \quad x \leqB y \iff y \geqB x
\end{eqnarray}
\end{proof}



\begin{thm}
Para todo par de elementos en un álgebra de Boole \(\algB\), \(x,y \in
\setB\), se cumple que, \(x\cdot y \leqB x\) y \(x\cdot y \leqB
y\).
\end{thm}



\begin{proof}
\begin{eqnarray}
&& (x \cdot y) \cdot y = x \cdot (y \cdot y) = x \cdot y\\
&& x \cdot y \leqB y\\
&& (x \cdot y) \cdot x = x \cdot (x \cdot y) = (x \cdot x) \cdot y = x \cdot
y\\
&& x \cdot y \leqB x\\
&& x \cdot y \leqB y \textsf{  y  } x \cdot y \leqB x
\end{eqnarray}
\end{proof}



\begin{thm}
Para todo par de elementos en un álgebra de Boole \( \algB \), \( x , y \in
\setB \), se cumple que, \( x + y \geqB x \) y \( x + y \geqB
y \).
\end{thm}



\begin{proof}
\begin{eqnarray}
&& (x + y) + y = x + (y + y) = x + y\\
&& x + y \geqB y\\
&& (x + y) + x = x + (x + y) = (x + x) + y = x + y\\
&& x + y \geqB x\\
&& x + y \geqB y \textsf{  y  } x + y \geqB x
\end{eqnarray}
\end{proof}



\begin{thm}
Para todo par de elementos en un álgebra de Boole \( \algB \), \( x , y \in
\setB \), se cumple que, \( x \cdot y \leqB x \) y \( x \cdot y \leqB y \).
\end{thm}



\begin{proof}
\begin{eqnarray}
&& (x \cdot y) \cdot y = x \cdot (y \cdot y) = x \cdot y\\
&& x \cdot y \leqB y\\
&& (x \cdot y) \cdot x = x \cdot (x \cdot y) = (x \cdot x) \cdot y = x \cdot
y\\
&& x \cdot y \leqB x\\
&& x \cdot y \leqB y \textsf{  y  } x \cdot y \leqB x
\end{eqnarray}
\end{proof}



\section{Cadenas de elementos con orden total.}



	Sea ahora una cadena de elementos relacionados bien por \(\leqB\), bien
por \(\geqB\), de manera excluyente, de forma que vaya desde el mínimo
\(0_{\setB}\) hasta el máximo \(1_{\setB}\). Esta cadena de elementos
conforma un orden total que es un suborden del general del retículo de Boole,
bien \(\left\langle \setB , \leqB \right\rangle\), bien  \(\left\langle \setB
, \geqB \right\rangle\) (procedentes del álgebra de Boole \(\algB\)).
Establezcamos las definiciones.



\newcommand{\PartesDeSinVacio}[1]{\ensuremath{\mathscr{P}\left(#1\right)\setminus{\set{\emptyset}}}}
\newcommand{\PartesDe}[1]{\ensuremath{\mathscr{P}\left(#1\right)}}
\newcommand{\EsCadena}[2]{\ensuremath{\mathsf{Chain}_{#1}\left(#2\right)}}
\newcommand{\EsCadenaSat}[2]{\ensuremath{\mathsf{SChain}_{#1}\left(#2\right)}}
\newcommand{\EsCadenaDisc}[2]{\ensuremath{\mathsf{DChain}_{#1}\left(#2\right)}}
\newcommand{\EsCadenaAsc}[2]{\ensuremath{\mathscr{AC}_{#1}\left(#2\right)}}
\newcommand{\EsCadenaDesc}[2]{\ensuremath{\mathscr{DC}_{#1}\left(#2\right)}}
\newcommand{\EsCadenaAscSat}[2]{\ensuremath{\mathscr{SAC}_{#1}\left(#2\right)}}
\newcommand{\EsCadenaDescSat}[2]{\ensuremath{\mathscr{SDC}_{#1}\left(#2\right)}}



\begin{defn}
Sea el subconjunto no vacío \(\emptyset \neq A \subseteq
\setB\). Diremos que constituye una \emph{ cadena ascendente },
\(\EsCadenaAsc{\algB}{A}\) , si y solo si,
para cualquier pareja de elementos de \(x y \in A\), bien
\(x \leqB y\) ó \(y \leqB x\). Formalmente:
\begin{eqnarray}
\nonumber
&&A \in \;\PartesDeSinVacio{\setB}\\
&&\EsCadenaAsc{\algB}{A} \logicdef \forall x \in A \forall y \in A \left( x
\leqB y \textsf{ ó } y \leqB x \right)
\end{eqnarray}
\end{defn}



\begin{defn}
Sea el subconjunto no vacío \(\emptyset \neq A \subseteq
\setB\). Diremos que constituye una \emph{ cadena descendente },
\(\EsCadenaDesc{\algB}{A}\) , si y solo si,
para cualquier pareja de elementos de \(x y \in A\), bien
\(x \geqB y\) ó \(y \geqB x\). Formalmente:
\begin{eqnarray}
\nonumber
&&A \in \;\PartesDeSinVacio{\setB}\\
&&\EsCadenaDesc{\algB}{A} \logicdef \forall x \in A \forall y \in A \left( x
\geqB y \textsf{ ó } y \geqB x \right)
\end{eqnarray}
\end{defn}



Estas cadenas, conjuntos totalmente ordenados dentro del retículo, a veces
pueden ser saturadas. Esto quiere decir que dada la cadena \(A \in
\PartesDeSinVacio{\setB}\), \(\EsCadenaAsc{\algB}{A}\), para todo otro
elemento de \(x \in \setB\), \(\neg\EsCadenaAsc{\algB}{A\cup\set{x}}\).



\begin{defn}
Sea el subconjunto no vacío \(\emptyset \neq A \subseteq
\setB\). Diremos que constituye una \emph{ cadena ascendente saturada },
\(\EsCadenaAscSat{\setB}{A}\) , si y solo si, \(\EsCadenaAsc{\setB}{A}\) y
además, \(\forall X \subset \setB\setminus A \textsf{ y } X\neq\emptyset\)
tenemos que \(A\cup{X}\) no puede ser cadena.
\begin{eqnarray}
\nonumber
&&A \in\PartesDeSinVacio{\setB}\\
&&\EsCadenaAscSat{\algB}{A} \logicdef \forall X \in
\PartesDeSinVacio{\setB\setminus{A}}
\quad\neg\left(\EsCadenaAsc{\algB}{A\cup{X}}\right)
\end{eqnarray}
\end{defn}



\begin{defn}
Sea el subconjunto no vacío \(\emptyset \neq A \subseteq
\setB\). Diremos que constituye una \emph{ cadena descendente saturada },
\(\EsCadenaDescSat{\algB}{A}\) , si y solo si, \(\EsCadenaDesc{\algB}{A}\) y
además, \(\forall X \subset \setB\setminus A \textsf{ y } X\neq\emptyset\)
tenemos que \(A\cup{X}\) no puede ser cadena.
\begin{eqnarray}
\nonumber
&&A \in\PartesDeSinVacio{\setB}\\
&&\EsCadenaDescSat{\algB}{A} \logicdef \forall X \in
\PartesDeSinVacio{\setB\setminus{A}}
\quad\neg\left(\EsCadenaDesc{\algB}{A\cup{X}}\right)
\end{eqnarray}
\end{defn}



	Otra propiedad que puede tener una cadena es la de ser \emph{cadena
discreta}. Esto ocurre si, siendo un elemento de una cadena saturada, existe
solo un elemento estrictamente mayor y uno estrictamente menor y ningún
elemento entre ellos. De hecho podría ser discreta por abajo o por arriba.



\begin{defn}
	Decimos que un elemento \(x \in \mathscr{C}\) tal que \( \EsCadena
{\algB} {\mathscr{C}} \), \emph{es discreto por arriba  en \(\mathscr{C}\)}
si \( \exists y \in \mathscr{C} \left( x < y \textsf{ y } x \neq y  \textsf{
y } \forall z \in \mathscr{C} \left( z \leq x \textsf{ ó } y \leq z \right)
\right) \).
\end{defn}



\begin{defn}
	Decimos que un elemento \(x \in \mathscr{C}\) tal que \( \EsCadena
{\algB} {\mathscr{C}} \), \emph{es discreto por abajo en \(\mathscr{C}\)} si
\( \exists y \in \mathscr{C} \left( y \leq x \textsf{ y } y \neq x  \textsf{
y } \forall z \in \mathscr{C} \left( z \leq y \textsf{ ó } x \leq z \right)
\right) \).
\end{defn}



\begin{defn}
	Decimos que un elemento \(x \in \mathscr{C}\) tal que \( \EsCadena
{\algB} {\mathscr{C}} \), \emph{es discreto en \(\mathscr{C}\)} si es \emph{
discreto por arriba en \(\mathscr{C}\)} y es \emph{ discreto por abajo en
\(\mathscr{C}\)}.
\end{defn}



\begin{defn}
	Decimos de una cadena \(\mathscr{C}\) en \(\algB\), que es cadena
discreta, \(\EsCadenaDisc{\algB}{\mathscr{C}}\), si \(\forall x \in
\mathscr{C} \left(x \textsf{ es discreto en } \mathscr{C}\right)\).
\end{defn}



\begin{defn}
	Dado un álgebra de Boole \(\algB\), decimos que un elemento \(x \in
\setB\) \emph{ es discreto } si \(\forall C \in \PartesDe{\setB}
\left(\EsCadena{\algB}{C} \textsf{ y } x \in C \then x \textsf{ es discreto
en } C\right)\).
\end{defn}



\begin{defn}
	Dado un álgebra de Boole \(\algB\), decimos que \emph{ \(\algB\) es
discreta } si \(\forall x \in \setB \left(x \textsf{ es discreto en } \algB
\right)\).
\end{defn}




	Ahora podemos dar el siguiente enunciado: si un álgebra de Boole es
finita, entonces es discreta. Este resultado es fácil de demostrar e
intuitivo por sí mismo, pero nos vendrá bien para reparar en algunas
propiedades de las álgebras de Boole finitas.



\begin{thm}
	Sea \(\mathscr{C} \subseteq \algB\) una cadena. Esta cadena siempre puede
embeberse en otra cadena de \(\algB\), digamos \(\mathscr{D}\), que será una
cadena saturada en \(\mathscr{C}\).
\end{thm}



\section{Conjuntos de átomos y de hiperátomos.}

	Si un elemento del álgebra de Boole \(\algB\), \(x \in \setB\), es tal
que \[ Z \mathdef \set{\inf\setCompr{x+y}{ x+y \neq x , x+y \neq y , y \in
\setB\setminus\set{1,0,x}}}\], el conjunto de elementos minimales
estrictamente mayores que él mismo, si \(Z \neq \emptyset\), diremos que
tiene la relación de orden \(\preccurlyeq\) discretizada en \(x\).
Igualmente, si es tal que  \[ W \mathdef \set{\sup\setCompr{x \cdot y}{ x
\cdot y \neq x , x \cdot y \neq y , y \in \setB\setminus\set{1,0,x}}} \], el
conjunto de elementos maximales estrictamente menores que él mismo, si \( W
\neq \emptyset \), decimos que la relación de orden \(\succcurlyeq\) está
discretizada en \(x\).




\begin{thm}
Sea \(\algB\) un álgebra de Boole no trivial de conjunto subyacente \(\setB\)
de cardinal finito. Sea \(x \in \setB\setminus\set{0} \) un elemento del
álgebra. Entonces existirá una serie de elementos \(y_1, y_2, \ldots ,
y_k\;\in\setB\setminus\set{0}\) con \(k \in \setN\), tal que \( x \gneq x
\cdot y_1 \gneq x \cdot y_1 \cdot y_2 \gneq \ldots \gneq x \cdot
{\prod}_{\imath=1}^{k} y_{\imath} = 0_{\setB}\).
\end{thm}



\begin{proof}
Por contradicción.
\begin{eqnarray}
\textsf{Hipótesis  } \exists x \in \setB\setminus\set{0} \nexists y \in \setB\setminus\set{0,\neg x}\quad x \cdot y \lneqq x
\end{eqnarray}
\end{proof}


	El concepto de sucesor o predecesor puede establecerse dentro del álgebra
de Boole, con las operaciones de suma y producto. Primero estableceremos los
elementos minimales que no son \(0_\setB\). Los llamaremos átomos del
álgebra:



\begin{equation}
\mathsf{Atom} \left( \algB \right) \mathdef \setCompr{x \in \setB \setminus
\set{0_\setB}}{ \forall y \in \setB \; \left( x \cdot y \in \set{x,0_\setB}
\right) }
\end{equation}



	Lo primero es tener en cuenta que no siempre este conjunto será no vacío.
Por ejemplo, en el ejemplo de las uniones finitas de intervalos genéricos
disjuntos en el intervalo unidad racional. En ese caso podéis ver que no
existen esos elemento atómicos a los que nos estamos refiriendo. Por otra
parte es claro que en las álgebras sobre las partes de un conjunto, los
átomos son los conjuntos con un solo elemento del conjunto de partida.



	Como estas últimas son las álgebras que más nos interesan, ya veremos
porqué, nos centraremos en las álgebras que tienen este conjunto
\(\mathsf{Atom}\left(\algB\right) \;\neq\; \emptyset\). Por ejemplo en el
álgebra de conmutación de Shannon, \(\algB_2\), podemos considerar que \(1\)
es el único átomo (es isomorfa al álgebra de las partes de un conjunto con un
solo elemento). Para álgebras de más elementos ya no será \(1\) un átomo.
Para el álgebra de las partes de un conjunto de dos elementos, \(\mathscr{P}
\set{a,\neg a} = \set{\emptyset, \set{a}, \set{\neg a}, \set{a,\neg a}}\), el
conjunto \(\mathsf{Atom}\left(\mathscr{P} \set{a,\neg
a}\right)=\set{\set{a},\set{\neg a}}\), será el conjunto de los átomos.



\begin{defn}
Sea \(\algB\) un álgebra de Boole, tal que tenga más de dos elementos.
\begin{equation}
  \mathsf{Atom}\left(\algB\right) \mathdef \setCompr{x \in \setB \setminus
  \set{0_\setB}} {\forall y \in \setB \; \left( x \cdot y \in \set{ x ,
  0_\setB } \right)}
\end{equation}
\end{defn}




\begin{thm}
Sea \(\algB\) un álgebra de Boole que admita conjunto de átomos no vacío.
Sean \(x \, y \; \in \mathsf{Atom}\left({\algB}\right)\quad x\neq y\),
entonces \(x \cdot y = 0_\setB\).
\end{thm}



\begin{proof}
\begin{eqnarray}
&&x \in \mathsf{Atom}\left({\algB}\right) \iff \forall z \in \setB \left(x
\cdot z = x \textsf{ ó } x \cdot z = 0\right) \\
&&y \in \mathsf{Atom}\left({\algB}\right) \iff \forall w \in \setB \left(y
\cdot w = y \textsf{ ó } y \cdot w = 0\right) \\
&&\textsf{Hipótesis   1   } \qquad\qquad x \neq y \\
&&\textsf{Hipótesis   2   } \qquad\qquad x = x \cdot y = \\
&&\text{Por conmutatividad   } \quad = y \cdot x = \\
&&\text{Por ser \(y\) átomo  } \;\;\qquad = y \\
&& x = y \textsf{ y } x \neq y \\
&& x \neq y \then x \cdot y \neq x \\
&& x \neq y \then x \cdot y = 0
\end{eqnarray}
\end{proof}





	Bien dado el conjunto de átomos de una algebra de Boole (no vacío),
tendremos que ver si mediante sumas podemos generar un álgebra de Boole.



	Sea \( \algB \) que admite un conjunto de átomos no vacío \( \mathbb{A}
\mathdef \mathsf{Atom} \left( \algB \right) \). Sea \( x \in \mathbb{A} \).






	También podemos desarrollar el concepto dual, el de hiperátomo y conjunto de hiperátomos. 



\begin{defn}
Sea \(\algB\) un álgebra de Boole, tal que tenga más de dos elementos.
\begin{equation}
  \mathsf{Htom}\left(\algB\right) \mathdef \setCompr{x \in \setB \setminus
  \set{1_\setB}} {\forall y \in \setB \; \left( x + y \in \set{ x ,
  1_\setB } \right)}
\end{equation}
\end{defn}




\begin{thm}
Sea \(\algB\) un álgebra de Boole que admita conjunto de h-tomos no vacío.
Sean \( x \, y \; \in \mathsf{Htom}\left({\algB}\right) \quad x \neq y \),
entonces \(x + y = 1_\setB \).
\end{thm}



\begin{proof}
\begin{eqnarray}
&&x \in \mathsf{Htom}\left({\algB}\right) \iff \forall z \in \setB \left(x
+ z = x \textsf{ ó } x + z = 1 \right) \\
&&y \in \mathsf{Htom}\left({\algB}\right) \iff \forall w \in \setB \left(y
+ w = y \textsf{ ó } y + w = 1 \right) \\
&&\textsf{Hipótesis   1   } \qquad\qquad x \neq y \\
&&\textsf{Hipótesis   2   } \qquad\qquad x = x + y = \\
&&\text{Por conmutatividad   } \quad = y + x = \\
&&\text{Por ser \(y\) h-tomo} \;\,\qquad = y \\
&& x = y \textsf{ y } x \neq y \\
&& x \neq y \then x + y \neq x \\
&& x \neq y \then x + y = 1
\end{eqnarray}
\end{proof}



	Supongamos ahora que \(\algB\) tiene un conjunto subyacente finito (y no
trivial),  \( \natural \setB \in \setN \setminus \set{0,1} \). Supongamos que
\(\natural \setB = n\). Pero ya hemos demostrado que \( \forall x \in \setB
\; \exists ! \neg x \in \setB \) y que se cumple, por ser \(\natural \setB >
1\), que \(x \neq \neg x\).



	Conociendo este resultado podemos definir una partición junto con la
relación de equivalencia que es el cociente de la partición.



\begin{defn}
Sean \(\algB\) un álgebra de Boole no trivial y \(\setB\) su conjunto
subyacente. Entonces \(\forall x,y \in \setB\), decimos que \(x
R_{\complement} y\) si y solo si \(y = \neg x\) ó \(y = x\).
\end{defn}



\begin{thm}
Sean \(\algB\) un álgebra de Boole no trivial y \(\setB\) su conjunto
subyacente. La relación interna antes definida \(R_{\complement}\) es una
relación de equivalencia.
\end{thm}


\begin{proof}
\begin{eqnarray}
&&\textsf{\(R_{\complement}\) es Reflexiva}\\
&&\text{Semántica de igualdad       }\quad\forall x \in \setB \quad x = x\\
&&\text{Definición de \(R_{\complement}\)}\quad\forall x \in \setB \quad x = x \then x R_{\complement} x\\ \nonumber
&&\text{por \verb"modus ponens"     }\quad\forall x \in \setB \quad x R_{\complement} x\\
&&\textsf{\(R_{\complement}\) es Simétrica}\\
&&\text{Hipótesis de la condición de simetría       } \quad x,y \in \setB \quad x R_{\complement} y\\
&&\text{Definición de \(R_{\complement}\)}\quad\forall x,y \in \setB \quad x R_{\complement} y \then y = \neg x \textsf{ ó } y = x\\
&&\text{por teorema de la doble negación     }\quad x R_{\complement} y \then \neg y = \neg \neg x = x \textsf{ ó } y = x\\
&&\text{por simetría de la igualdad     }\quad\forall x,y \in \setB \quad x R_{\complement} y \then x = \neg y \textsf{ ó } x = y\\
&&\text{Definición de \(R_{\complement}\)}\quad\forall x,y \in \setB \quad x R_{\complement} y \then y R_{\complement} x\\ \nonumber
&&\textsf{\(R_{\complement}\) es Transitiva}\\
&&\text{Hipótesis de la condición de transitividad       } \quad x,y,z \in \setB \quad x R_{\complement} y \textsf{ y } y R_{\complement} z\\
&&\text{Definición de \(R_{\complement}\)}\quad x R_{\complement} y \then y = \neg x \textsf{ ó } y = x\\
&&\text{Definición de \(R_{\complement}\)}\quad y R_{\complement} z \then z = \neg y \textsf{ ó } z = y\\
&&\text{\(1^{er}\) Caso de 4   }\quad \neg y = x \textsf{ y } z = \neg y\\
&&\text{\(1^{er}\) Caso de 4   }\quad \neg y = x \textsf{ y } z = \neg y \then z = x\\
&&\text{\(1^{er}\) Caso de 4   }\quad z = x \then x R_{\complement} z\\
&&\text{\(2^{do}\) Caso de 4   }\quad \neg y = x \textsf{ y } z = y\\
&&\text{\(2^{do}\) Caso de 4   }\quad \neg z = x \then x R_{\complement} z\\
&&\text{\(3^{er}\) Caso de 4   }\quad y = x \textsf{ y } \neg z = y\\
&&\text{\(3^{er}\) Caso de 4   }\quad \neg z = x \then x R_{\complement} z\\
&&\text{\(4^{to}\) Caso de 4   }\quad y = x \textsf{ y } z = y\\
&&\text{\(4^{to}\) Caso de 4   }\quad z = x \then x R_{\complement} z
\end{eqnarray}
\end{proof}



\begin{thm}
Sean \(\algB\) un álgebra de Boole no trivial y \(\setB\) su conjunto
subyacente. \({[x]}_{R_{\complement}}\) tiene cardinalidad
\(\natural{[x]}_{R_{\complement}}=2\), siendo la clase de equivalencia,
exactamente, \({[x]}_{R_{\complement}} = \set{x,\neg x}\).
\end{thm}



\begin{thm}
Sean \(\algB\) un álgebra de Boole no trivial y \(\setB\) su conjunto
subyacente. Si \(\natural\setB = n \in \setN\setminus\set{0,1}\), entonces
\(\natural\left(\setB / R_{\complement}\right) = n / 2\). Luego \(2 \mid n\).
\end{thm}



	Tenemos ahora herramientas par investigar las álgebras de Boole finitas.
Una primera pregunta es ¿son las álgebras de Boole finitas álgebras que
admiten conjunto de átomos no vacío? Veremos que sí.



	Para comenzar, el álgebra de conmutación, \(\algB_{2}\) es la más pequeña
no trivial y siempre hay un homomorfismo inyectivo que la incrusta en
cualquier álgebra de Boole. \(\algB_{2}\) decimos que tiene un solo átomo, el
que coincide con \(1_{{\algB}_{2}}\). Utilizamos el término átomo para
álgebras de más de \(2\) elementos. Por ejemplo, el álgebra de las partes de
un conjunto de cardinal \(2\) tiene cardinal \(4\), como vimos en un ejemplo
detallado al principio de este escrito. Este álgebra se construye sobre
\(\setB \mathdef \mathscr{P}\set{a, b}\). Este álgebra es atómica, y sus
átomos son \(\mathsf{Atom}\left(\mathscr{P}\set{a, b}\right) =
\set{\set{a},\set{b}}\). Este álgebra tiene un átomo más que el anterior
\({\algB}_{2}\). Así hemos detallado las álgebras con
\(\mathsf{Atom}\left(\mathscr{P}\set{a, b, c}\right) = \set{\set{a},\set{b},
\set{c}}\) y \(\mathsf{Atom}\left(\mathscr{P}\set{a, b, c, d}\right) =
\set{\set{a},\set{b}, \set{c}, \set{d}}\). Aún más, al describir las álgebras
sobre el conjunto potencia de un conjunto \(U\), \(\left\langle \mathscr{P}
U, \cup, \cap \right\rangle\) tiene como conjunto de átomos
\(\mathsf{Atom}\left(\mathscr{P} U\right) = \setCompr{\set{x} \in \mathscr{P}
U}{x \in U}\), independientemente del cardinal del conjunto \(U\).



	Ahora nos hace falta saber algo más sobre las álgebra de Boole atómicas.
Para el caso \({\algB}_{2}\) está claro que todo el álgebra es su elemento
atómico y su complementario. Como hay un solo átomo parece que todo se reduce
a ese elemento y su complementario. Cuando el álgebra crece algo más, ponemos
un átomo más, la cosa se puede complicar. Por ejemplo, en
\( \mathscr{P} \left[ \set{a,b} \right] \), \( 1_{\mathscr{P} \left[
\set{a,b} \right]} \) deja de ser un átomo, para ser la unión de los átomos,
y cada átomo es el complementario del otro átomo. Esta situación también será
pasajera. Volviendo a añadir un átomo obtenemos \( \mathscr{P} \left[
\set{a,b,c} \right] \), dónde ningún átomo es complementario de otro. El \(
1_\setB \) vuelve a ser la unión de los átomos. Ahora el complementario de
cada átomo es un h-tomo. Esta situación se mantendrá ya en todas las álgebras
atómicas con más de \(2\) átomos, que es cuando aparecen los conjuntos \(
\mathsf{Atom} \left( \algB \right) \cap \mathsf{Htom} \left( \algB \right) =
\emptyset \). 





\printbibliography



\end{document}



